\documentclass[11pt,a4paper]{article}
\usepackage[T1]{fontenc}
\usepackage{lmodern,amsmath,amssymb,amsthm,mathtools}
\usepackage[margin=26mm]{geometry}
\usepackage{microtype,booktabs,longtable,array,enumitem}
\usepackage[authoryear,round]{natbib}
\usepackage[hidelinks]{hyperref}
\usepackage{url}
\hypersetup{pdftitle={Counting satisfying assignments on smooth projective varieties, after Katz},pdfauthor={}}
\newtheorem{theorem}{Theorem}[section]
\newtheorem{lemma}[theorem]{Lemma}
\newtheorem{proposition}[theorem]{Proposition}
\newtheorem{corollary}[theorem]{Corollary}
\theoremstyle{definition}
\newtheorem{definition}[theorem]{Definition}
\newtheorem{example}[theorem]{Example}
\newtheorem{question}[theorem]{Question}
\theoremstyle{remark}\newtheorem{remark}[theorem]{Remark}
\newcommand{\F}{\mathbb F}
\newcommand{\Fp}{\mathbb F_p}
\newcommand{\A}{\mathbb A}
\newcommand{\PP}{\mathbb P}
\newcommand{\OO}{\mathcal O}
\newcommand{\Z}{\mathbb Z}
\newcommand{\Q}{\mathbb Q}
\newcommand{\SAT}{\#\mathrm{SAT}}
\newcommand{\ModP}{\mathrm{Mod}_p\mathrm P}
\newcommand{\workbenchurl}{https://github.com/KokunoYumeto/arithmetic-geometry-complexity-workbench}
\newcommand{\recordurl}{https://github.com/KokunoYumeto/mathematics-workbenches}
\DeclareMathOperator{\Spec}{Spec}
\DeclareMathOperator{\Proj}{Proj}
\DeclareMathOperator{\Pic}{Pic}
\DeclareMathOperator{\Tr}{Tr}
\DeclareMathOperator{\rank}{rank}
\DeclareMathOperator{\ord}{ord}
\DeclareMathOperator{\CH}{CH}
\DeclareMathOperator{\Span}{span}
\setlength{\emergencystretch}{2em}
\setlength{\parskip}{2pt}
\setlist{nosep}
\raggedbottom
\title{Counting satisfying assignments on smooth projective varieties, after Katz}
\author{}
\date{}
\begin{document}
\maketitle
\begin{abstract}
Katz expressed the Frobenius action on the top coherent cohomology of an anticanonical hypersurface as a coefficient of its equation. We combine this description with a Boolean encoding and an explicit smoothing construction. At every fixed prime, modular counting is complete for smooth geometrically integral circuit hypersurfaces whose top coherent cohomology is a line. For their Artin--Mazur formal groups, every fixed exact finite height and every fixed lower threshold at least two is modular-counting complete; the next height counts cyclic orbits of balanced assignment configurations. The full height, including infinity, is computable in polynomial space. Integral invariant derivatives satisfy sharp height-dependent divisibility laws. Explicit sparse equations admit further algorithms: a small residual cofactor gives deterministic scalar recovery after saturation, and bounded complete-intersection excess after certified coordinate, digit and prime-field relation reductions gives an always-correct randomized algorithm of expected polynomial time. Other constructions record exact counts in Tate motives, Chow groups, algebraic K-theory and fixed-field Frobenius multiplicities. Unrestricted saturation, effective projective compression, infinite-height hardness and small-Betti modular encoding remain open. The mathematics and text were produced by GPT-6 Astra (ChatGPT Pro mode), with separate same-model proof checks; no human has checked this paper.
\end{abstract}
\section{The main statements}
The starting point is a coefficient formula of Nicholas Katz. Katz's Algorithm (2.3.7.14) and Special Case (2.3.7.17), printed pages 35--36 of \citet{core-katz}, calculate the Frobenius action on a smooth hypersurface from its defining polynomial. The work here is an encoding of Boolean assignments into that coefficient, together with equations whose derivatives prove smoothness over the ground field for every Boolean input.

\begin{theorem}[Counting on a coherent line]\label{thm:main}
Fix a prime $p$. Given a Boolean formula $\Phi$ with an explicitly declared set of $n$ variables and $S$ satisfying assignments, one can construct in deterministic polynomial bit time a homogeneous arithmetic circuit over $\Fp$ for a degree-$d$ polynomial in $d$ variables, defining
\[
X\subset\PP^{d-1}_{\Fp},\qquad d\ge3,
\]
such that $X$ is smooth and geometrically integral, $\omega_X\simeq\OO_X$, and
\[
H^i(X,\OO_X)=
\begin{cases}\Fp,&i=0,d-2,\\0,&0<i<d-2.\end{cases}
\]
On the top line the $p$-power Frobenius is multiplication by $S\bmod p$, and
\[
\#X(\Fp)\equiv1+S\pmod p.
\]
The construction is uniform in a varying prime with running time polynomial in $|\Phi|+p$. Evaluation of the scalar is complete for the function class consisting of accepting-path counts modulo $p$; deciding its nonvanishing is $\ModP$-complete under deterministic many-one reductions on the stated smoothness promise. Smoothness recognition is not required of either oracle.
\end{theorem}

The input representation is part of this statement. A circuit can multiply many short local polynomials without listing the resulting monomials. When the same equation is given by an explicit list of $t$ monomials, its coefficient can instead be computed in $t^{p-1}\operatorname{poly}(L)$ bit operations at a fixed prime, where $L$ is the sparse input length. The theorem therefore distinguishes presentations of a smooth variety, even when the requested coherent group has dimension one.

\begin{theorem}[Stabilization over the coefficient space]\label{thm:stabilization-intro}
Let $V_u$ be the affine space of homogeneous degree-$u$ polynomials in $u\ge1$ variables over $\Fp$. If $Q$ is any positive multiple of $p$ with $Q\ge u+2$, there is an explicit affine-linear embedding of $V_u$ into the coefficient chart of degree-$d$ forms in $d$ variables with one prescribed coefficient equal to one, where $d=Q+1$. Every geometric fibre is a smooth anticanonical hypersurface. The top coherent direct image is a canonically trivialized line, and its Frobenius coefficient pulls back to
\[
h_p(G)=[(x_1\cdots x_u)^{p-1}]G^{p-1}.
\]
The least admissible $Q$ gives $d\le u+p+2$. A sparse input gains only $O(d)$ terms; a circuit gains $O(d\log d)$ gates.
\end{theorem}

This theorem applies to singular and zero input forms. The input form itself is not being asserted to define a smooth variety. The universal family is smooth after adding coordinates. In this family the rational-point congruence is $1+(-1)^d h_p(G)$; the prime-power choice in Theorem~\ref{thm:main} makes this $1+S$.

\paragraph{Formal-group heights.}
Artin and Mazur attach a formal group to the top coherent line;
Stienstra expresses its logarithm by integral coefficient numerators.
For a smooth degree-$d$ hypersurface $X\subset\PP^{d-1}_{\Fp}$,
$d\ge3$, write $\Phi_X$ for its degree-$(d-2)$ Artin--Mazur formal
group and $h(X)$ for its height, allowing $h(X)=\infty$.
It is a smooth one-dimensional formal group, including the plane cubic
case through its formal Picard group. If $\widetilde F$ is any
homogeneous integral lift, its normalized logarithm is
\[
 L(T)=\sum_{m\ge1}b_m\frac{T^m}{m},\qquad
 b_m=[(x_1\cdots x_d)^{m-1}]\widetilde F^{m-1}.
\]
These are \citet[II, Corollaries~2.11--2.12 and~4.2,
pp.~106--110]{fg-artin-mazur} and
\citet[Theorem~1, p.~907]{fg-stienstra}. In particular the first height
is detected by $b_p\bmod p$, not by reduction of the possibly
nonintegral coefficient $b_p/p$.

\begin{theorem}[Fixed height predicates]\label{fg:thm:complexity-intro}
Fix a prime $p$. Inputs are division-free arithmetic circuits over $\Fp$
promised to compute homogeneous degree-$d$ polynomials in $d\ge3$
variables defining smooth geometrically integral projective
hypersurfaces, with dimension allowed to grow. The following statements
hold under deterministic polynomial-time many-one reductions.
\begin{enumerate}
\item Deciding $h(X)=1$ is $\ModP$-complete, and
      $h(X_\Phi)=1$ in Theorem~\ref{thm:main} precisely when
      $S\not\equiv0\pmod p$.
\item For every fixed integer $r\ge2$, deciding $h(X)\ge r$ is
      $\ModP$-complete; the predicate includes infinite height.
\item For every fixed finite $h\ge1$, deciding $h(X)=h$ is
      $\ModP$-complete. Hardness in these statements persists after
      imposing any fixed lower bound on the dimension.
\end{enumerate}
At fixed $p$, every fixed finite height predicate is polynomial time for
an explicit sparse equation. In fixed dimension the full height of a
circuit equation is polynomial-time computable. No polynomial bound in
a varying binary-encoded prime or threshold is asserted by the fixed
parameter statements.
\end{theorem}

The coefficient and reduction proofs are in
Sections~\ref{fg:sec:height}--\ref{fg:sec:space}.
Exact-height hardness uses an explicit prime-field existence theorem.

\begin{theorem}[Explicit realization of every finite height]
\label{fg:thm:seeds-intro}
For every prime $p$ and every finite integer $h\ge1$, there is an
explicit smooth geometrically integral hypersurface over $\Fp$ of
degree $p^h+p+1$ in $p^h+p+1$ homogeneous variables whose
one-dimensional top Artin--Mazur formal group has exact height $h$.
Its dimension is $p^h+p-1$, and its canonical bundle is trivial.
Its equation can be constructed in time polynomial in its displayed
degree, with $p$ and $h$ as parameters; no polynomial bound in the
binary length of a varying $h$ is asserted.
\end{theorem}

The construction in Section~\ref{fg:sec:seeds} prescribes a unique
multinomial contribution to each relevant logarithm numerator and then
preserves those coefficients under smoothing. It gives equations over
the prime field. The all-height construction of
\citet[Theorem~4.18]{fg-kty-two} works over an algebraically closed
field and uses a smaller degree; its nonempty parameter loci alone
do not specify a prime-field point.

The next height has an explicit counting interpretation. For a
width-three CNF $\phi$, use two variables $w_{f,i,0},w_{f,i,1}$ for
each clause-variable incidence. If $I_f$ is the set of variables in
clause $f$ and $R_f\subseteq\{0,1\}^{I_f}$ is its set of satisfying
local assignments, define
\[
 P_f=\sum_{a\in R_f}\prod_{i\in I_f}w_{f,i,a_i}.
\]
Put
\[
 P=\prod_fP_f\prod_i
       \left(\prod_{f\ni i}w_{f,i,0}+\prod_{f\ni i}w_{f,i,1}\right),
 \qquad Z=\prod_{(f,i)}w_{f,i,0}w_{f,i,1}.
\]
An unused declared variable receives a tautological unary relation.
If there is no incidence, first add a forced variable. A local
configuration chooses one term from each displayed factor, so
coefficient multiplicities distinguish these choices.

\begin{theorem}[The second height counts cyclic orbits]\label{fg:thm:orbit-intro}
Fix $p$, and let $S$ be the number of satisfying assignments of $\phi$.
Set $A_2(\phi)=[Z^{p+1}]P^{p+1}$. There is a function
$C(\phi)\in\#\mathrm P$ satisfying the exact integer identity
\[
 A_2(\phi)=S^2+pC(\phi).
\]
It counts the nonfixed cyclic orbits of balanced ordered $(p+1)$-tuples
of local configurations: their product monomial is $Z^{p+1}$, and the
cyclic group rotates the first $p$ configurations while fixing the last.
For the hypersurface of Theorem~\ref{thm:main},
\[
 h(X_\phi)\ge3
 \iff S=0\pmod p\ \text{and }C(\phi)=0\pmod p
 \iff A_2(\phi)=0\pmod{p^2}.
\]
On $S=0\bmod p$, its second normalized height invariant is
$b_{p^2}/p=C(\phi)\bmod p$. This invariant can differ for two CNF
presentations of the same Boolean function, even with the same output
degree and ambient dimension.
\end{theorem}

\begin{theorem}[The full height in polynomial space]\label{fg:thm:space-intro}
Fix $p$ and the circuit promise of Theorem~\ref{fg:thm:complexity-intro}.
If $L$ is the complete binary input length, the finite height can be
output in binary, or infinity identified, with
$O_p(L^2\log(L+2))$ working bits. Every finite height is at most
$\binom{pd-1}{d-1}$. The algorithm halts on all syntactically valid
inputs, with correctness required on the promise. In particular,
infinite-height decision belongs to polynomial space.
\end{theorem}

The geometric finite operator is due to
\citet[Theorem~C]{fg-kty-one}, and its matrix formulation to
\citet[\S3.4]{fg-bgp}. The space bound follows by accessing its entries
from the original circuit and recursively computing one matrix-power
entry at a time. A matching hardness classification for infinite height
is not established here.

\begin{theorem}[Integral derivatives of a height-$h$ lift]\label{fg:thm:derivative-intro}
Let $R=W(k)$ for a perfect characteristic-$p$ field, let $\mathcal F$ be
a one-dimensional commutative formal-group law over $R$, and let
$D=(L'(T))^{-1}d/dT$ be its normalized integral invariant derivation.
If the reduction has height at least $h\ge1$, put
\[
 c_h=\frac1{p-1}-\frac1{p^h-1},\qquad
 E_h(a)=\sum_{\substack{j\ge1\\h\nmid j}}\lfloor a/p^j\rfloor.
\]
For every $f\in R[[T]]$ and $a\ge0$,
\[
 D^af\in p^{\max\{0,E_h(a),\lceil c_ha-(h-1)\rceil\}}R[[T]].
\]
At exact height $h$, for $j=hr+s$, $0\le s<h$,
$v_p(D^{p^j}(T^{p^s})(0))=E_h(p^j)$. The slope $c_h$ and uniform
affine error $h-1$ are optimal. At infinite height the exact uniform
exponent is $v_p(a!)$.
For a Stienstra lift of height at least $r$, with $q=p^r$,
\[
 h(X)\ge r+1\iff b_q\in p^rR
      \iff D^qR[[T]]\subseteq q!R[[T]].
\]
\end{theorem}

Thus $S=0\bmod p$ implies the height-two derivative law, with slope
$p/(p^2-1)$ and affine error one, for every integral Stienstra lift over
$\Z_p$. In particular $\overline D^{p}=0$. The next height is detected
by the additional divisibility of $D^{p^2}T(0)$: its valuation is $p$
at height two and at least $p+1$ at height at least three. The generic
proofs are given in Section~\ref{fg:sec:derivatives}; directional
statements for a specified Honda Cartier dual require its additional
integral model data.



\paragraph{Presentation boundaries.}
For a fixed prime and a fixed number of homogeneous sparse equations, counting projective points modulo that prime is polynomial time, with no smoothness or reducedness assumption. A logarithmic number of trinomials also suffices for this algorithm. An arbitrary number of sparse equations can nevertheless present the smooth circuit hypersurface isomorphically. A second distinction concerns linear equations: over $\F_2$, any number of linear equations and a fixed number of sparse nonlinear equations admit polynomial-time counting modulo every fixed power of two. At odd primes the smooth anticanonical complete-intersection problem is already modular-counting complete with many trinomials that are linear and one sparse nonlinear equation. Its nonlinear monomial support matters: for fixed $p,c$, support at most $\sigma$ gives a counting bound $p^{c(p-1)\sigma}\operatorname{poly}(L)$ in the presence of arbitrary linear constraints.

For a smooth anticanonical complete intersection on an $m$-dimensional linear solution space, a sharper bound uses the degree constraint. Let $c$ be the fixed number of nonlinear equations, and let $\sigma_i$ be the largest support of a monomial in the $i$th equation after discarding monomials with an individual coordinate exponent at least $p$. Put $K=(p-1)\sum_i\sigma_i-m$. The scalar is zero if $K<0$, and otherwise takes $T^{c(p-1)}p^K\operatorname{poly}(L)$ operations, where $T$ is the total number of displayed terms. Thus an excess of $O(\log L)$ over the minimum possible support permits a polynomial-time algorithm. The proof reduces each selected product to a weighted distribution on a vector space of dimension at most $K$.

The support algorithm has a complementary hardness result. At each fixed odd prime and for every fixed $\eta>0$, the one-nonlinear-equation smooth family can be constructed with $K\le d^\eta$. In particular its largest monomial support can be at most $(1/(p-1)+\varepsilon)d$ for any fixed $\varepsilon>0$. Distributing the padding among the tagged permanent terms preserves both the scalar and the derivative proof of smoothness. The interval between logarithmic excess and arbitrarily small fixed powers of $d$ is not classified.

\begin{theorem}[A small cofactor in a saturated presentation]\label{thm:saturation-intro}
Fix a prime $p$ and a nonnegative integer $B$. Let $f_1,\ldots,f_s\in\Fp[x_0,\ldots,x_N]$ be explicit homogeneous sparse polynomials, with binary exponents, such that
\[
(f_1,\ldots,f_s):(x_0,\ldots,x_N)^\infty=(g),
\]
where $g$ defines a smooth geometrically integral hypersurface of degree $N+1$ and positive dimension. If at least one nonzero $f_j$ has degree at most $N+1+B$, its cofactor $f_j/g$ can be recovered up to a ground-field unit, and the scalar of Frobenius on $H^{N-1}(X,\OO_X)$ can be computed, in deterministic polynomial bit time. The degrees of the other inputs are unrestricted. Consequently the scalar is polynomial-time computable if every cofactor $f_i/g$ has individual degree at most $B$.
\end{theorem}

The saturation promise means that the cofactors have no common projective zero. On every affine line they therefore generate the unit ideal. Univariate gcd computations recover the values of the one small cofactor, while reducing all other equations modulo the low-degree selected input keeps their binary degrees affordable. A truncated quotient then gives the scalar without expanding $g$. A designated cofactor of polynomially bounded numerical degree and sparsity also suffices, by ordinary sparse interpolation. The remaining principal-ideal question allows every residual cofactor to lie outside these bounds.

Before stating the higher-codimension result, fix the following sparse
preprocessing rule on the displayed generating list of $K$.
For a set $H$ of coordinate indices, repeatedly
remove each row's full monomial content outside $H$, then apply the
largest rank-one base-$p^r$ digit reduction $AB^{p^r}\mapsto AB$
with $r\ge1$ and $B$ nonconstant, and then a minimum-degree
replacement from Theorems~\ref{sp:thm:linearized-preprocessing},
\ref{sp:thm:affine-preprocessing}, and~\ref{sp:thm:zero-sum-preprocessing}.
The first recognizes
$F=U^{p^e}+\sum_{i=0}^{s}a_i V^{p^e-p^i}U^{p^i}$, where
$V$ is a monomial, $\deg U=\deg V=d\ge1$, $s<e$,
$(d+1)p^s\le p^e$, and $1+\sum_i a_i\ne0$, and replaces $F$ by $U$.
With the same degree gap, the second recognizes the additional term
$\beta V^{p^e}$ when $1+\sum_i a_i=0$ and $\beta\ne0$, and replaces
$F$ by $U,V$, after removing the $V$-term of $U$.
The third treats the homogeneous zero-sum case without the additional
term, requires $p^e>p$, and replaces $F$ by $U^p-V^{p-1}U$ after the
same normalization. This retains every possible prime-field component.
All coefficients lie in $\Fp$, and all these data are recovered from
the row. Normalize each output by its first nonzero coefficient,
skip zero rows, and repeat on the list until no step applies.
Write $\mathcal P_H(K)$ for the resulting ideal. The recognition
tests and their iteration have polynomial sparse-bit cost.

\begin{theorem}[Bounded excess after certified sparse reduction]\label{thm:saturation-random-intro}
Fix a prime $p$, a codimension $c\ge2$, and an integer $E\ge0$.
Let $K\subset\Fp[x_0,\ldots,x_N]$ be given by an explicit
homogeneous sparse list with binary exponents. Suppose that its
saturation is the ideal $I_X$ of a smooth geometrically integral
anticanonical global complete intersection
$X\subset\PP^N_{\Fp}$ of dimension $N-c\ge1$.
Set $H_X=\{j:x_j\in I_X\}$ and $K^\circ=\mathcal P_{H_X}(K)$.
Suppose that, after a finite field extension, $K^\circ$ contains a
homogeneous regular sequence of length $c$ whose degree sum is at
most $N+1+E$. Neither $H_X$, the sequence, its profile, nor the
extension need be given.
For every positive integer $a$, a randomized algorithm of polynomial
bit cost in the sparse input length and $a$ returns the Frobenius
scalar on $H^{N-c}(X,\OO_X)$ or reports failure, with failure
probability at most $2^{-a}$. Every returned scalar is correct on
the promise. Repetition has expected polynomial bit cost, and the
other input degrees are unrestricted. The only random choices are
projection centres; all coordinate divisions are certified from the
original equations, using a deterministic grid of at most $N+1$
fibres for each tested tuple and centre.
If $\dim(K^\circ)_1=c-1$, the scalar is computable deterministically
in polynomial bit time.
\end{theorem}

The bound concerns a sequence in a computably processed input ideal,
not merely the small generators of its saturation. Since
$K\subseteq K^\circ$, bounded excess in the original ideal is a
special case. The algorithm examines the polynomially many possible
sets of at most $c-1$ vanishing coordinates. A nonzero coordinate
class in a finite algebra obtained by imposing the original rows on a
tested fibre certifies that the
coordinate does not vanish identically on $X$, and hence that division
by its powers is valid. For the correct branch a degree-one grid finds
all these certificates. Theorem~\ref{sci:thm:certified-preprocessing}
and Corollary~\ref{sci:cor:certified-linear} give the full proof.
Example~\ref{sci:ex:certified-fermat} has unbounded least excess after
digit reduction alone and excess zero after the certified divisions.
Example~\ref{sp:ex:constant-field-family} survives both those
operations and reaches excess zero through the prime-field relation.
Example~\ref{sp:ex:zero-sum-family} survives all the earlier relation
rules as well; retaining its prime-field components gives a cofactor
of polynomial degree and support, so the designated sparse-cofactor
algorithm applies.

For a safely processed list, deterministic degree-profile candidates
and a fixed finite-fibre grid recover the required annihilator once a
good centre is known. The fibre algebras retain their nilpotents.
A centre over an extension of polynomial degree coprime to $p-1$
gives deterministic scalar recovery by
Proposition~\ref{sci:prop:fibre-grid}. When the tuple itself is defined
over $\Fp$, normalization of its colon line removes the coprimality
requirement on the centre field.

There is a further deterministic case under the same geometric
hypotheses: if the input designates an original regular $c$-tuple
of degree sum at most $N+1+E$, and gives factorizations of its first
$c-1$ forms into nonzero $\Fp$-linear forms with binary multiplicities,
Theorem~\ref{sci:thm:factored-prefix} constructs a centre and recovers
the scalar in polynomial bit time. The factors need not belong to
$K$, and $K_1$ may be zero. In both linear-prefix results $c$ is fixed;
the odd-prime hard construction above has an unbounded number of
independent linear equations.

\begin{theorem}[A maximum-degree bound for saturation]\label{thm:saturation-ci-intro}
Fix a prime $p$, a codimension $c\ge2$, and an integer $E\ge0$. Let an explicit homogeneous sparse list over $\Fp$, with binary exponents, have saturation equal to the ideal of a smooth geometrically integral anticanonical global complete intersection $X\subset\PP^N$, of dimension $N-c\ge1$. Let $\delta$ be the largest degree of a nonzero input. If
\[
0\le c\delta-(N+1)\le E,
\]
the Frobenius scalar on $H^{N-c}(X,\OO_X)$ is computable in deterministic polynomial bit time. Neither minimal generators nor a regular sublist need be given.
\end{theorem}

This deterministic proof constructs regular combinations over a finite
extension of polynomial degree. A graded annihilator of bounded degree
identifies the desired line inside the dual cohomology of that complete
intersection, and its semilinear eigenvalue descends by a field norm.
Wiebe's determinant identities and the Koszul coefficient formula carry
the algebraic step. A deterministic algorithm also applies when a regular
sublist of bounded total excess is given and all input degrees lie within
a fixed gap above its smallest degree. The preceding randomized theorem
removes that gap and discovers the regular sequence under its existential
bounded-excess promise.

\paragraph{Other counting presentations.}
A smooth sparse cubic over $\mathbb Z$ carries the exact witness count as the multiplicity of a Tate motive; its Chow groups and algebraic K-theory give corresponding readouts. This affine construction complements the modular scalar on a projective coherent line. Explicit Boolean wiring can expose disjoint affine partitions and permit exact counting. Sparse Walsh presentations can permit witness recovery even when every coherent trace vanishes. Fixed-field sparse univariate constructions encode a count in the number of equal-length Frobenius orbits, and smooth Artin--Schreier curves turn the same count into genus and geometric $p$-rank. Smooth affine hyperplane complements admit parsimonious reductions over every fixed field with more than two elements. Each assertion requires its own input model; cohomological rank, number of equations, support size, and extension-field degree cannot be interchanged.

A hidden presentation can sometimes be recovered. Over a fixed odd prime, suppose the base coordinates are marked and a designated quadratic subspace is exactly the span of $\ell_i^2-1$, where the private linear parts of the affine forms $\ell_i$ are independent. Two deterministically chosen quadratic combinations recover those forms, up to signs and permutation, in polynomial time. With a fixed number of further equations, fixed monomial support, and bounded rank of their coefficient-row powers, this gives a polynomial-time empty-fibre test. An arbitrary equation list need not exhibit the designated subspace, and general private permutations need not preserve the rank bound.

\paragraph{Why the distinction matters.}
The Weil conjectures control the absolute values of Frobenius eigenvalues. Recovering their multiplicities from a compact system of equations is a separate computational question. Our finite split examples isolate this distinction exactly, while the projective construction shows that even a single coherent Frobenius entry can carry modular counting. The reductions are unconditional. Their obstruction to a polynomial-time compression into a bounded number of sparse equations is conditional on $\ModP\not\subseteq\mathrm P$ at the fixed prime.

\paragraph{Credit and scope.}
Katz supplies the coherent coefficient formula. The binary incidence encoding predates the smoothing used here; the earlier smooth parity constructions over binary extensions, over $\F_4$, and by restriction of scalars to $\F_2$ are compared below with their precise coherent dimensions. Wan's modular sparse-counting theorem, Valiant's permanent reduction, the complete-intersection Picard theorem of SGA~2, and the coefficient calculations for regular sequences supply distinct ingredients of the presentation results. We give the proofs needed for the exact forms asserted here and state the unresolved saturation and compression problems separately.

Recent work also distinguishes several meanings of circuit complexity. The OpenAI preprint \citep[Chapter~5, Theorems~1.1--1.3, pp.~115--116]{root-openai-permanent} reports lower bounds of $\Omega(n^2\log\log n)$ gates for division-free circuits computing the formal permanent over $\mathbb C$, and $\Omega(n^4/\log n)$ variable leaves for formulas, including valid divisions. These statements concern symbolic computation of a polynomial. Valiant's reduction used here concerns evaluation on finite-field matrices and transfers a Boolean count to geometry; neither assertion supplies the other. The geometric reductions below do not depend on those preprint lower bounds.

\section{Proof overview}
The Boolean part uses two variables for each clause-variable incidence. A local satisfying choice and a consistent variable choice together contribute a squarefree monomial precisely when they describe a satisfying assignment. Distinct auxiliary tags then force each designated summand to occur exactly once in a $(p-1)$st power. Wilson's theorem fixes the sign at odd primes.

The geometric part pads the original form and adds pure powers with derivatives that are especially simple in characteristic $p$. Two derivatives force one coordinate to vanish. The remaining derivatives leave only a small possible gradient locus, on which the derivative in the first added coordinate is nonzero. This proves smoothness uniformly in every input coefficient. The hypersurface sequence identifies its top coherent group with a one-dimensional Laurent-monomial quotient, so the same coefficient becomes the Frobenius scalar. Summing a polynomial indicator over the affine cone independently checks the sign in the rational-point congruence.

The formal-group argument first identifies multiplication by the prime on principal-unit cocycles with coherent Frobenius. Stienstra's integral logarithm then identifies every later height test with a successively more precise coefficient. For the second test, cyclic rotation of balanced collections has exactly the squared assignment count as its fixed locus, and the remaining orbits give an explicit counting function. Products in disjoint variable blocks amplify divisibility while the smoothing construction preserves every coefficient below its chosen scale. To realize every exact finite height over the prime field, we prescribe term multiplicities and use binary prefix vectors to force a unique multinomial contribution. Its factorial valuation gives the desired first nonzero height invariant, and smoothing preserves all coefficients needed for that test. A finite monomial operator from quasi-F-splitting supplies a stopping rule for infinite height; recursive access to its matrix powers keeps the working space polynomial. The derivative bounds follow from integral translation operators and a rescaled exponential, with separate witnesses for the slope and its affine error.

The presentation results compare two different operations. Expanding a bounded number of sparse equations is cheap at a fixed prime. Expanding a circuit product can be exponential. To list the latter projectively we add coordinates for gate values and use several affine charts to prevent extra scheme structure at infinity. Many linear equations can also hide products of linear forms inside a single sparse monomial. In characteristic two each such monomial simply fixes some coordinates to one; at odd primes its values carry additional information. Restricting its coordinate support restores an explicit algorithm by projecting the linear solution space to that support.

For saturation, a low-degree listed multiple becomes a modulus for every other equation. The absence of a common zero among the cofactors prevents an affine-line restriction from introducing an extra common factor, which is the usual difficulty in using projected gcds. The resulting quotient evaluates the small cofactor at arbitrary points. Total-degree interpolation and a determinant-one coordinate translation then reduce the desired Frobenius coefficient to a polynomial number of sparse term choices. This argument needs only one bounded-degree cofactor, because all other degrees enter through modular binary exponentiation.

In higher codimension, the scalar is recovered from an annihilator line.
A complete intersection whose degree sum exceeds the anticanonical sum by
a fixed amount has a canonical section space of polynomial dimension.
Multiplication relations identify in that space the line belonging to $X$,
and the Frobenius adjoint is again given by sparse coefficient extraction.
To handle other equations of arbitrarily large degree, a finite linear
projection gives a common monomial basis for every fibre algebra, including
its nilpotents. The residual quotient is free over the same projection
base. Consequently the obstruction to annihilation has degree bounded by
the excess, regardless of the degrees of the input equations, and a fixed
Newton interpolation grid detects it exactly. Enumerating degree
profiles and using prime-degree field extensions constructs a short
candidate list deterministically: the chosen coefficients remain linearly
independent over the function field of every unwanted component. The
random step chooses a good projection centre for a regular candidate.
A field norm and an inverse exponent remove the choice of basis and the
coprime extension degree.

Before this recovery, finite quotients of the original equations certify which coordinate factors can be removed. On the correct branch, the degree-one coordinates have affine-linear coefficients in a graded free basis, so a fixed degree-one grid detects every nonzero class. Wrong branches may fail to provide certificates but cannot certify a forbidden division. Geometric integrality also forces algebraic constants to lie in the prime field; a recognized linearized relation then permits further degree reduction. Their hidden factors are proposed by exponent sections and checked by complete identities. This reduces the remaining excess condition beyond coordinate and digit reduction. For a factored linear prefix, one finds a suitable point on each linear component by the fixed-parameter semidiagonal identity test, then merges those points by one-variable searches to construct a common centre.

Finally, the structured examples make information loss visible. Walsh coefficients describe a function on a split Boolean fibre, while a trace sees only one aggregate. Private multipliers can erase several zeta coefficients without changing which fibres are empty. In the fixed-field univariate examples, purity leaves a whole orbit multiplicity undetermined. These examples motivate precise recovery questions instead of attributing computational power to smoothness alone.

\section{Boolean assignments as a homogeneous coefficient}
This section gives the entire counting encoding, including its sign, before any geometry is imposed. A formula can first be converted to a Boolean circuit with fan-in at most two. Introducing a variable for the output of each gate and enforcing its truth table gives a width-three conjunctive normal form with a unique extension of each primary input. We give the exact gate clauses in Section~\ref{sec:uniform}; thus this conversion preserves the number of satisfying assignments, not merely satisfiability.

Let $\Phi$ now be a width-three CNF. Treat a clause $f$ as a relation on its distinct variable scope $I_f$; repeated literals and opposite signs on the same variable are evaluated in that relation. For every declared variable that appears nowhere, add a tautological unary relation. If necessary add a fresh variable with a positive unit clause so that at least one incidence is present. These operations preserve $S$.

For each incidence $e=(f,i)$, $i\in I_f$, introduce two variables $w_{e,0},w_{e,1}$. Write $\mathcal I$ for the incidence set and $b=2|\mathcal I|$. Define polynomials over $\Z$ by
\begin{align}
P_f&=\sum_{\substack{a\in\{0,1\}^{I_f}\\a\text{ satisfies }f}}
                 \prod_{i\in I_f}w_{(f,i),a_i},\label{eq:local-clause}\\
Q_i&=\prod_{f:i\in I_f}w_{(f,i),0}
       +\prod_{f:i\in I_f}w_{(f,i),1},\nonumber\\
P&=\prod_f P_f\prod_iQ_i,\qquad
Z=\prod_{e\in\mathcal I}w_{e,0}w_{e,1}.\label{eq:incidence}
\end{align}
An empty clause contributes the zero polynomial, with its scope degree recorded formally. Both $P$ and $Z$ are homogeneous of degree $b$.

\begin{lemma}[Incidence coefficient]\label{lem:incidence}
The integer coefficient $[Z]P$ equals $S$. The factored expression for $P$ is a homogeneous arithmetic circuit of size polynomial in $|\Phi|$. Every clause factor has at most eight monomials and every variable factor has two.
\end{lemma}
\begin{proof}
A term selected from $Q_i$ uses one common label at all incidences of $i$. To obtain $Z$, every adjacent clause factor must use the complementary label. These complementary labels give one consistent Boolean assignment. The available monomials in each $P_f$ say exactly that its local restriction satisfies $f$. Conversely a satisfying assignment selects its local clause monomials and the complementary variable monomials uniquely. All their coefficients are one, so this is a bijection with the contributions to the coefficient, including multiplicities. Each local scope has size at most three. Products for the variable factors, followed by the product of the local factors, require a number of multiplication gates linear in the incidence count. Balanced products give logarithmic depth.
\end{proof}

Now reduce coefficients modulo $p$, introduce $p-1$ further variables $t_1,\ldots,t_{p-1}$, and put $u=b+p-1$. The polynomial
\begin{equation}\label{eq:normalization}
G=-t_1^{p-1}P+\sum_{j=2}^{p-1}t_j^{p-1}Z
\end{equation}
has degree $u$ in exactly $u$ variables. It remains a compact homogeneous circuit even if the product $P$ has many monomials.

\begin{lemma}[Odd-prime normalization]\label{lem:normalization}
For $G$ in \eqref{eq:normalization},
\[
h_p(G)=[(Zt_1\cdots t_{p-1})^{p-1}]G^{p-1}=S\quad\text{in }\Fp.
\]
\end{lemma}
\begin{proof}
Each tag occurs in just one summand, with exponent $p-1$. Its target exponent forces that summand to be selected exactly once among the $p-1$ factors. There are $(p-1)!$ orderings. The remaining product is $-PZ^{p-2}$, and its target coefficient is $-[Z]P$. Wilson's theorem gives $-(p-1)!=1$ in $\Fp$. At $p=2$ there is one tag, the displayed sum is empty, and the same equality holds since minus equals plus. At an odd prime, replacing the first minus sign by a plus would give $-S$.
\end{proof}

\section{A smooth stabilization in families}
This section proves smoothness by the derivative equations and shows that the coefficient identity survives arbitrary base change.

\begin{theorem}[Explicit stabilization]\label{thm:stabilizer}
Let $R$ be an $\Fp$-algebra, let $G\in R[x_1,\ldots,x_u]$ be homogeneous of degree $u\ge1$, and let $Q$ be a positive multiple of $p$ such that $Q\ge u+2$. Set $m=Q-1$ and $d=Q+1=m+2$. In $d$ variables define
\begin{align}
A_G&=G(x_1,\ldots,x_u)\prod_{i=u+1}^{m}x_i+x_1^m,\label{eq:padding}\\
F_G&=yzA_G+y^{Q+1}+\sum_{i=1}^{m}x_i^{Q+1}-x_1z^Q.\label{eq:stabilizer}
\end{align}
Then $F_G$ is homogeneous of degree $d$, and
\begin{equation}\label{eq:hasse-stabilizer}
h_p(F_G)=h_p(G)
\end{equation}
as an identity in the coefficients of $G$. For the universal coefficient ring, \eqref{eq:stabilizer} defines a smooth projective family with geometrically integral fibres and a trivialized top coherent direct image. The map of coefficient spaces is affine-linear and injective into the chart in which the coefficient of $y^d$ is one.
\end{theorem}
\begin{proof}
Both summands of $A_G$ have degree $m$, so every summand of $F_G$ has degree $m+2$. To contribute to $(x_1\cdots x_myz)^{p-1}$ in $F_G^{p-1}$, no factor can use $-x_1z^Q$, since $Q\ge p$. The only remaining source of $z$ is $yzA_G$. Therefore all $p-1$ factors must select it. The padding coordinate $x_{u+1}$ occurs with exponent one in the first summand of $A_G$ and not at all in $x_1^m$. Its target exponent forces the first summand in every factor. Removing the common padding gives exactly \eqref{eq:hasse-stabilizer}. This argument also works when $Q=p$ and $m=p-1$; one must not discard $x_1^m$ merely by claiming that $m\ge p$ in that case.

Over an algebraic closure of any residue field, differentiation gives
\[
(F_G)_y=zA_G+y^Q,\qquad (F_G)_z=yA_G,
\qquad y(F_G)_y-z(F_G)_z=y^{Q+1}.
\]
At a common zero of the partial derivatives, $y=0$. The other derivatives then become
\[
(F_G)_{x_i}=x_i^Q-\delta_{i1}z^Q.
\]
Consequently $x_2=\cdots=x_m=0$ and $x_1^Q=z^Q$. A projective point satisfying these equalities has both $x_1$ and $z$ nonzero. Because $m>u$, the padding product vanishes on the first coordinate axis and $A_G(x_1,0,\ldots,0)=x_1^m$. Hence
\[
(F_G)_y=z x_1^m\ne0,
\]
a contradiction. This proof is valid in characteristic two. For a general multiple $Q$ of $p$, it does not assert that $x_1^Q=z^Q$ implies $x_1=z$; that extra assertion is available only for a pure $p$-power and is unnecessary. Euler's identity has coefficient $d=1$ in $R$, so the entire projective gradient locus is empty.

Let $R_0=\Fp[V_u]$ be the universal coefficient ring. The polynomial $F_G$ is monic in $y$, so $R_0[x,z,y]/(F_G)$ is free as an $R_0[x,z]$-module. It is therefore flat over $R_0$. Homogeneous localization is flat, and taking its degree-zero summand preserves flatness. The projective family is flat and finitely presented, and the calculation above proves that all geometric fibres are smooth. The fibrewise criterion for smoothness now proves relative smoothness.

Distinct monomials of $G$ go to distinct monomials after multiplication by $yz\prod_{i>u}x_i$. Their exponent positions do not coincide with those of the fixed summands in \eqref{eq:stabilizer}. The input coefficients are therefore recovered by projecting onto these coefficient positions. This proves the affine-linear embedding. The coherent line and geometric integrality follow from the calculation in Lemma~\ref{lem:coherent} below, which is valid after every base change.
\end{proof}

The least multiple $Q\ge u+2$ satisfies $Q<u+2+p$, and therefore $d\le u+p+2$. A list of $T$ input terms gains at most $m+3$ terms. A circuit gains $O(d\log d)$ gates by balanced products and repeated squaring. Each gate index and exponent has polynomial binary length. The number of degree-$u$ monomials in $u$ variables can be exponential in $u$; the circuit statement does not materialize that coefficient space in the input.

For Theorem~\ref{thm:main}, choose instead the least power $Q$ of $p$ with $Q-1>u$ and $Q-1\ge p$. It satisfies
\[
Q\le p\max\{u+2,p+1\}.
\]
At an odd prime this $Q$ is odd and $d=Q+1$ is even. This choice preserves the required sign in the rational-point formula without changing polynomial bit size. The sharper stabilization statement concerns the Hasse coefficient itself, so it does not require this parity restriction.

\subsection{The coherent generator and the point-count check}
\begin{lemma}[Katz's line in the hypersurface sequence]\label{lem:coherent}
Let $F$ be a degree-$d$ homogeneous equation in $d\ge3$ variables $v_1,\ldots,v_d$, defining a smooth hypersurface $X$ over an $\Fp$-algebra as in Theorem~\ref{thm:stabilizer}. The top coherent line is trivialized, compatibly with base change, by the class corresponding to
\[
e=(v_1\cdots v_d)^{-1}\in H^{d-1}(\PP^{d-1},\OO(-d)).
\]
Absolute Frobenius sends $e$ to $h_p(F)e$; on a general base this means $ce\mapsto c^p h_p(F)e$.
\end{lemma}
\begin{proof}
Write $R$ for the base algebra, and let
$\pi:X\to\operatorname{Spec}R$ and
$\operatorname{pr}:\PP^{d-1}_R\to\operatorname{Spec}R$ be the
structure morphisms. Use the exact sequence
\[
0\longrightarrow\OO_{\PP}(-d)\xrightarrow{\ F\ }\OO_{\PP}
\longrightarrow\OO_X\longrightarrow0.
\]
The cohomology of projective space gives $\pi_*\OO_X=\OO_{\mathrm{base}}$, vanishing of the intermediate direct images, and a connecting isomorphism
\[
R^{d-2}\pi_*\OO_X\simeq R^{d-1}\operatorname{pr}_*\OO_{\PP}(-d).
\]
These statements follow directly from the standard affine-cover complex and hold under arbitrary base change. Its top quotient has degree-$(-d)$ Laurent monomials modulo those with some nonnegative exponent. The sole surviving monomial is $e$, so it is a free rank-one module.

The Frobenius map on the ideal lifts as $s\mapsto F^{p-1}s^p$, because $(Fs)^p=F(F^{p-1}s^p)$. Naturality of the connecting map therefore gives the class
\[
\frac{F^{p-1}}{(v_1\cdots v_d)^p}.
\]
Only a numerator monomial with every exponent $p-1$ survives the Laurent quotient. Its coefficient is $h_p(F)$. This is the rank-one specialization of Katz's Algorithm (2.3.7.14) and Special Case (2.3.7.17), pp. 35--36 \citep{core-katz}.

On each geometric fibre, $H^0(\OO_X)$ is the ground field, so the fibre is connected. A smooth scheme is regular, and intersecting irreducible components would make a local ring have more than one minimal prime. Its irreducible components are thus disjoint; connectedness makes the fibre integral. Adjunction gives $\omega_X\simeq\OO_X$.
\end{proof}

For an independent check of the sign, sum the indicator $1-F(v)^{p-1}$ over $\Fp^d$. For $j\ge0$,
\[
\sum_{a\in\Fp}a^j=
\begin{cases}-1,&j>0\text{ and }p-1\mid j,\\0,&\text{otherwise}\end{cases}
\quad\text{in }\Fp,
\]
where $a^0=1$ also at $a=0$. The total degree of $F^{p-1}$ is $d(p-1)$, so the only contributing exponent vector is $(p-1,\ldots,p-1)$. If $C$ is the affine cone count, then
\[
C\equiv(-1)^{d+1}h_p(F),\qquad
C=1+(p-1)\#X(\Fp).
\]
It follows that
\begin{equation}\label{eq:point-sign}
\#X(\Fp)\equiv1+(-1)^d h_p(F)\pmod p.
\end{equation}
The prime-power choice above yields $1+S$, and at $p=2$ the sign has no effect. Lemmas~\ref{lem:incidence}--\ref{lem:coherent} establish the geometric and arithmetic assertions of Theorem~\ref{thm:main}.

\subsection{Entire parameter lines and quartic K3 surfaces}
\label{sec:entire-lines}
The gradient argument also works for any homogeneous degree-$m$ form $A$ with $A(1,0,\ldots,0)\ne0$. In that case the $z$-exponent argument gives $h_p(F)=h_p(A)$ without padding.

For odd $p$, take $Q=p$, $m=p-1$, and
\[
A_a=x_1^{p-1}+\cdots+x_{p-2}^{p-1}-a x_{p-1}^{p-1}
\quad(a\in\A^1_{\Fp}).
\]
Its first-axis value is one. Each of its $p-1$ terms must occur once in the target, giving $h_p(A_a)=-(p-1)!a=a$. Thus \eqref{eq:stabilizer} gives a smooth family over the entire affine line, of degree $p+1$ in $\PP^p$, with Hasse section $a$. In characteristic two take $G=ax_1$, $Q=4$, and $A_a=ax_1x_2x_3+x_1^3$ instead; the resulting quintics in $\PP^4$ have the same property.

For $p=3$ write the variables as $r,s,y,z$. The family is
\begin{equation}\label{eq:k3}
yz(r^2-a s^2)+y^4+r^4+s^4-rz^3=0.
\end{equation}
It consists of smooth quartic surfaces in $\PP^3$, with trivial canonical sheaf and $H^1(\OO)=0$, hence K3 surfaces. Its top coherent generator corresponds to $(rsyz)^{-1}$, and Frobenius is multiplication by $a$.

Here is an explicit count over $\F_3$. On the chart $r=1$, the equation becomes
\[
1+s^2+y^2+z\{y(1-a s^2)-1\}=0,
\]
because $t^4=t^2$ and $t^3=t$ for $t\in\F_3$. A nonzero coefficient of $z$ gives one solution; a zero coefficient gives three or zero according as the constant is zero or nonzero. For $s=0$ the contribution is $2$. For $s=1$ or $2$ the contribution per value of $s$ is $5,3,5$ for $a=0,1,2$, respectively. Thus this chart contributes $12,8,12$.

On $r=0$, normalize the first nonzero coordinate among $s,y,z$ to one. If $s=1$, the equation is $1+y^2-a y z=0$: it contributes $0,2,2$ in the three parameter values. If $s=0$ and $y=1$, the equation is $1=0$ and contributes nothing. The final point $[0:0:0:1]$ always lies on the surface. The total counts are consequently
\[
\begin{array}{c|ccc}
a&0&1&2\\\hline
\#X_a(\F_3)&13&11&15\\
1+a\pmod3&1&2&0
\end{array}
\]
as required. The fibre at $a=0$ is smooth with zero coherent Frobenius. This computation alone does not determine its full crystalline Newton polygon or assert that it is supersingular.

The construction gives entire lines in degree $p+1$ for odd $p$, in degree five for $p=2$, and in all larger degrees congruent to one modulo $p$ obtainable by stabilization with a padding variable. More explicitly, odd $p$ admits every $d\equiv1\pmod p$ with $d\ge p+1$: keep the degree-$p-1$ form $A_a$ above as input and use padding for $d\ge2p+1$. At $p=2$ every odd $d\ge5$ is obtained from $G=ax_1$. These are constructive sufficient degree classes, not a classification of all possible smooth pencils in other degrees.

Katz's Dwork-pencil calculation, (2.3.7.18)--(2.3.7.23), p. 36, evaluates the same coefficient for $\sum X_i^d-d\lambda\prod X_i$ on its smooth parameter open set. In \eqref{eq:k3} and its higher-dimensional versions the gradient calculation proves smoothness at every finite parameter, including the zero of the Hasse section. Knutson's point-count and degeneration identities and Poonen's use of characteristic-$p$ powers for derivative control are related ingredients, with different hypotheses and conclusions \citep[Theorems 1--2]{core-knutson}\citep[\S2.3]{core-poonen}.

\subsection{A unit clause through the whole construction}
\label{sec:unit-clause}
Take $p=3$ and the positive unit clause $(s)$, for which $S=1$. Let $a,b$ denote the incidence variables for zero and one, and let $c,e$ be the tags. Then
\[
P=b(a+b),\qquad Z=ab,\qquad
G=-c^2b(a+b)+e^2ab.
\]
The sole contribution to $(abce)^2$ in $G^2$ is the cross product of $-abc^2$ and $abe^2$, with coefficient $-2=1$ in $\F_3$. Here $u=4$. Choose the original prime-power $Q=9$, so $m=8$ and $d=10$, and identify $(x_1,x_2,x_3,x_4)=(a,b,c,e)$. The equation is
\[
F=yz\{Gx_5x_6x_7x_8+a^8\}
  +y^{10}+\sum_{i=1}^{8}x_i^{10}-az^9.
\]
It is smooth by the derivative argument. The connecting map carries the generator of $H^8(X,\OO_X)$ to
\[
\frac1{abcex_5x_6x_7x_8yz},
\]
and Frobenius multiplies this line by one. Over $\F_3$, put $A=Gx_5x_6x_7x_8+a^8$. The equation is linear in $z$:
\[
(yA-a)z+y^2+\sum_{i=1}^{8}x_i^2=0.
\]
Summing the contributions $1$, $3$, or $0$ over the $3^9$ values of $(x,y)$ gives $19\,907$ affine points, hence $(19\,907-1)/2=9\,953$ projective points. The reproducible finite count and the later two-variable cubic example use the same convention that a projective point has its first nonzero coordinate normalized to one.

\section{Residue evaluation, nonvanishing, and sparse input}
This section states the complexity claims as promise problems and proves their membership as well as their hardness.

Write $\#_p\mathrm P$ for functions whose value is the least nonnegative residue modulo $p$ of the number of accepting paths of a polynomial-time nondeterministic machine. Write $\mathrm{Mod}_p\mathrm P$ for the corresponding decision class, with acceptance when that residue is nonzero. A promise algorithm need only agree with the specified problem on inputs satisfying its promise; membership below is witnessed by a total counting function that agrees there.

The parsimonious gate clauses in Section~\ref{sec:uniform} reduce accepting-path counts to width-three CNF counts. Lemmas~\ref{lem:incidence} and \ref{lem:normalization}, followed by the proven smooth stabilization, give $\#_p\mathrm P$-hardness of scalar evaluation and $\mathrm{Mod}_p\mathrm P$-hardness of nonvanishing. The construction produces promised instances without a smoothness oracle.

For membership, a nondeterministic machine guesses a point of $\PP^{d-1}(\Fp)$ in its unique representative with first nonzero coordinate one. It can guess the coordinate index and the later coordinates in fixed-width binary words, reject out-of-range encodings, and evaluate the circuit. It therefore has exactly $C_X=\#X(\Fp)$ accepting paths. A disjoint fixed-size nondeterministic gadget with $p-1$ accepting paths realizes $C_X-1$ modulo $p$. If needed, making $p-1$ copies realizes its negative. By \eqref{eq:point-sign}, this gives the scalar $(-1)^d(C_X-1)$ modulo $p$ on every promised input. The machine uses polynomial time for a fixed $p$. Nothing in this membership proof tests smoothness, checks a circuit identity, or expands the polynomial. This completes Theorem~\ref{thm:main}.

\begin{proposition}[Frobenius splitting]\label{prop:splitting}
For a smooth geometrically connected projective variety $X$ over $\Fp$ with $\omega_X\simeq\OO_X$ and $H^{\dim X}(X,\OO_X)$ a line, nonvanishing of its coherent Frobenius is equivalent to Frobenius splitting. Thus the nonvanishing hardness applies to Frobenius splitting on the same smooth circuit hypersurfaces.
\end{proposition}
\begin{proof}
Here $F_X:X\to X$ denotes the absolute Frobenius morphism. A splitting of $\OO_X\to (F_X)_*\OO_X$ makes its map on top cohomology injective, hence nonzero. Conversely, Serre duality identifies the dual of this map with the map
\[
\operatorname{Hom}((F_X)_*\OO_X,\OO_X)\longrightarrow
\operatorname{Hom}(\OO_X,\OO_X)
\]
given by composing with $\OO_X\to (F_X)_*\OO_X$, after choosing a trivialization of $\omega_X$. The target is $\Fp$, since $X$ is proper and geometrically connected. If coherent Frobenius is nonzero, this map is onto; a homomorphism mapping to $1$ is a splitting. This is the rank-one form of the criterion in \citet[Corollary 2.5 and Remark 2.6]{core-hlyy}.
\end{proof}

For comparison, let $F$ be listed as $t$ monomials with binary exponent vectors and binary coefficients reduced modulo $p$. Enumerate the $t^{p-1}$ ordered choices in $F^{p-1}$. For each, sum the exponent vectors, test equality with $(p-1,\ldots,p-1)$, and multiply the coefficients modulo $p$. If $E$ is the largest input exponent, each sum uses at most $O(\log(2+(p-1)E))$ bits per coordinate. With an explicit coordinate list of length $d$, the cost is
\[
O\bigl(t^{p-1}d(p-1)(\log(2+E)+\log p)^2\bigr)
\]
with elementary bit arithmetic, plus input parsing and coefficient reduction; a coarser $t^{p-1}\operatorname{poly}(L)$ bound accommodates compressed sparse exponent lists. At fixed $p$ this is polynomial. There is no inconsistency with circuit hardness: the incidence product $\prod_fP_f\prod_iQ_i$ is represented compactly by multiplication gates, and its expanded support need not be polynomial.

\citet[Theorem 1.1, p. 3 of the author manuscript]{core-wan} gives, for one sparse affine polynomial in $n$ variables with $t$ monomials over $\F_{p^h}$, a deterministic computation of its zero count modulo $p^b$ in
\[
O\!\left(ntpb\left(\frac{3th}{(h+b)(p-1)}+3\right)^{(h+b)(p-1)}\right)
\subseteq O\bigl(n(8t)^{(h+b)p}\bigr)
\]
bit operations after reducing each positive exponent to its equivalent value in $\{1,\ldots,p^h-1\}$. Binary exponent reduction adds polynomial input-reading cost. Thus fixed characteristic, fixed extension degree, and fixed precision give polynomial time in sparse input length. The theorem concerns one explicitly sparse equation; a circuit product is a different representation.

\section{Formal-group height and the next counting invariant}\label{fg:sec:height}
The coherent line also detects the first height of an Artin--Mazur formal
group. Higher heights require integral coefficient data of greater
$p$-adic precision. We first establish the geometric interpretation and
the exact precision needed, before analyzing the counting construction.

\subsection{The top formal group and its first coefficient}
Let $k=\Fp$, let $d\ge3$, and let $X=V(F)\subset\PP^{d-1}_k$ be a smooth
geometrically integral hypersurface defined by a homogeneous polynomial
of degree $d$. Write $n=d-2$. For a local Artinian $k$-algebra $A$ with
residue field $k$, define
\[
 \Phi_X(A)=\ker\bigl(H^n_{\mathrm{et}}(X_A,\mathbb G_m)
                    \longrightarrow H^n_{\mathrm{et}}(X,\mathbb G_m)\bigr).
\]
Here $X_A=X\times_k\Spec A$. The height of a one-dimensional formal group
in characteristic $p$ is $h$ if the first nonzero term of its
multiplication-by-$p$ series has degree $p^h$; it is infinite if that series
is zero. Thus $h\ge r$ always includes infinite height.

\begin{proposition}[The coherent line and height one]\label{fg:prop:height-one}
The functor $\Phi_X$ is a smooth one-dimensional commutative formal group.
In a coordinate with tangent generator as in Lemma~\ref{lem:coherent},
\[
 [p]_{\Phi_X}(T)=H_p(F)T^p\pmod{T^{p+1}},\qquad
 H_p(F)=[(x_1\cdots x_d)^{p-1}]F^{p-1}.
\]
Consequently
\[
 h(X)=1\iff H_p(F)\ne0\iff\#X(\Fp)\not\equiv1\pmod p.
\]
On the family of Theorem~\ref{thm:main}, this scalar is $S\bmod p$.
For fixed $p$, both $h(X)=1$ and $h(X)\ge2$ are $\ModP$-complete
promise problems for circuit input in growing dimension. Each is in
deterministic polynomial time for explicit sparse input, or for circuit
input in fixed dimension. In dimension one the only heights are one and
two.
\end{proposition}

\begin{proof}
The hypersurface sequence gives
\[
 H^0(X,\OO_X)=k,\quad H^i(X,\OO_X)=0\ (0<i<n),\quad
 H^n(X,\OO_X)=k,\quad H^{n+1}(X,\OO_X)=0.
\]
For $n\ge2$, pro-representability follows from
\citet[II, Corollary~4.2, p.~109]{fg-artin-mazur}; the next-degree
vanishing gives smoothness, and the tangent space is $H^n(X,\OO_X)$
\citep[II, Corollaries~2.4--2.5 and~2.11--2.12,
pp.~106--107]{fg-artin-mazur}. Their complete-intersection special case
on p.~110 gives the same conclusion.

For $n=1$, the previous-degree vanishing would be false. Instead the
degree-zero deformation functor is $A\mapsto1+\mathfrak m_A$, since
$H^0(X_A,\OO_{X_A})=A$. It is formally smooth. Corollary~2.12 of
Artin--Mazur therefore applies again, since $H^2(X,\OO_X)=0$.
Their identification in II, \S4, p.~109, makes $\Phi_X$ the formal
completion of the Picard functor at the identity. After extension to an
algebraic closure, choosing a point identifies this formal group with
that of the resulting elliptic curve. Its height is one or two
\citep[IV, Theorem~7.4 and Corollary~7.5, p.~134]{fg-silverman}.
Height is unchanged by a field extension, so this holds over $k$ as well.

We verify the first coefficient directly, including its normalization.
For $I=\mathfrak m_A$, the constant family splits the reduction map on
unit sheaves. Thus
\[
 \Phi_X(A)=H^n(X,1+I\OO_X).
\]
The successive quotients of the $I$-adic filtration are coherent additive
sheaves. They are acyclic on affine opens, so a finite affine cover
computes this cohomology by multiplicative \v Cech cocycles; it also
identifies its Zariski and \'etale descriptions. This is the
principal-unit description of \citet[II.(1.3)--(1.4), pp.~103--104,
and Corollary~2.6, p.~106]{fg-artin-mazur}.

Choose a nonzero tangent vector $\eta\in H^n(X,\OO_X)$ and the matching
formal coordinate. Over $k[t]/(t^{p+1})$, the point with coordinate $t$
has a cocycle representative with components $1+t a_I+O(t^2)$, where
$(a_I)$ represents $\eta$. Multiplication by $p$ raises each component
to its $p$-th power, giving
\[
 (1+t a_I+O(t^2))^p=1+t^p a_I^p\pmod{t^{p+1}}.
\]
If coherent Frobenius carries $\eta$ to $c\eta$, naturality of the tangent
identification under $k[\epsilon]/(\epsilon^2)\to k[t]/(t^{p+1})$,
$\epsilon\mapsto t^p$, identifies this point with coordinate $ct^p$.
Hence $[p](T)=cT^p\bmod T^{p+1}$. Lemma~\ref{lem:coherent} gives
$c=H_p(F)$, and equation~\eqref{eq:point-sign} gives the point-count
equivalence. Over $\Fp$ this scalar is unchanged by a nonzero change of
tangent basis: the semilinear rescaling factor is $\lambda^{p-1}=1$.

For completeness, the first nonzero degree of a homomorphism of formal
groups in characteristic $p$ is a power of $p$. Its leading term
$aT^e$ satisfies $(X+Y)^e=X^e+Y^e$ by comparing total degree $e$ in the
homomorphism identity. The base-$p$ expansion of $(1+X)^e$ forces $e$ to
be a $p$-power. The linear term of $[p]$ is zero. This proves the claimed
height criterion, including the zero-series case.

A nondeterministic machine counts the projective rational points by
guessing the unique representative whose first nonzero coordinate is
one. Adding $p-1$ accepting paths gives a count congruent to
$\#X(\Fp)-1$. Smoothness forces every ambient variable to occur in the
circuit, since an omitted variable produces a singular coordinate
vertex. The number of guessed coordinates is therefore polynomial in
the input length, even when their labels are binary encoded.
This proves membership for height one. If $a$ is any accepting-path
count, then
\[
 1+(p-1)a^{p-1}\equiv
 \begin{cases}1,&a\equiv0\pmod p,\\0,&a\not\equiv0\pmod p.
 \end{cases}
\]
Sums and fixed powers are again accepting-path counts. Thus $\ModP$ is
closed under complement, with this explicit fixed-prime construction.
Theorem~\ref{thm:main} and parsimonious width-three encodings give
hardness for both predicates. The same reduction permits any fixed
lower bound on the output dimension by choosing a larger $p$-power $Q$.
For sparse input, enumerate the $t^{p-1}$ ordered term selections in
$F^{p-1}$. In fixed dimension, enumerate the fixed finite set of
projective $\Fp$-points. These give the asserted polynomial algorithms.
\end{proof}

\subsection{Integral logarithms and every height test}
Let now $k$ be any perfect field of characteristic $p$, put $R=W(k)$,
and normalize $v_p(p)=1$. Choose a homogeneous lift
$\widetilde F\in R[x_1,\ldots,x_d]$ of the equation, and define
\[
 b_m=[(x_1\cdots x_d)^{m-1}]\widetilde F^{m-1},\qquad b_1=1.
\]
\citet[Theorem~1, p.~907; specialization on p.~908 and in
\S3.9, p.~914]{fg-stienstra} constructs an integral formal-group law
$\mathcal F$ lifting $\Phi_X$, with normalized logarithm
\begin{equation}\label{fg:eq:stienstra}
 L(T)=\sum_{m\ge1}\frac{b_m}{m}T^m.
\end{equation}
Indeed, the positive integer vectors of length $d$ and total $d$ consist
only of $(1,\ldots,1)$, so his index set has one element. The lift is
flat: if $pA=\widetilde F B$, reduction modulo $p$ in the polynomial
domain gives $B=pB_1$, and cancellation gives $A=\widetilde F B_1$.
Thus the quotient is torsion-free over the discrete valuation ring.
The defining sequence has one nonzero element and is regular. These
verify the base and flatness hypotheses in Stienstra's theorem.
The coefficient of $T^p$ in \eqref{fg:eq:stienstra} is $b_p/p$;
the coherent scalar is the reduction of its integral numerator $b_p$.

\begin{lemma}[Logarithmic test for the next height]\label{fg:lem:height-log}
Let $\mathcal F$ be any one-dimensional commutative formal-group law
over $R$, with normalized logarithm
$L(T)=T+\sum_{m\ge2}\ell_mT^m$. If its reduction has height at least
$r\ge1$, and $q=p^r$, then
\[
 \ell_m\in R\ (m<q),\qquad p\ell_q\in R,
 \qquad p\ell_q\equiv[T^q][p]_{\mathcal F}(T)\pmod p.
\]
Its height is at least $r+1$ precisely when $p\ell_q\in pR$.
For the logarithm \eqref{fg:eq:stienstra}, this says
\begin{equation}\label{fg:eq:next-height}
 b_{p^r}\in p^{r-1}R,\qquad
 h(X)\ge r+1\iff b_{p^r}\in p^rR
 \quad\text{if }h(X)\ge r.
\end{equation}
In particular,
\begin{align}
 h(X)\ge r&\iff p^j\mid b_{p^j}\quad(1\le j<r),\label{fg:eq:all-height}\\
 h(X)=r&\iff p^j\mid b_{p^j}\quad(1\le j<r),
                  \quad p^r\nmid b_{p^r},\label{fg:eq:exact-height}\\
 h(X)=\infty&\iff p^j\mid b_{p^j}\quad\text{for every }j\ge1.
 \label{fg:eq:infinite-height}
\end{align}
\end{lemma}

\begin{proof}
Write $g(T)=[p]_{\mathcal F}(T)=pT+\sum_{m\ge2}\gamma_mT^m$.
Comparing coefficients in $L(g(T))=pL(T)$ gives
\begin{equation}\label{fg:eq:log-recursion}
 (p-p^m)\ell_m=\gamma_m+
       \sum_{i=2}^{m-1}\ell_i[T^m]g(T)^i.
\end{equation}
The height hypothesis means $\gamma_m\in pR$ for $m<q$.
Inductively, every summand on the right for $m<q$ is in $pR$,
and $(p-p^m)/p$ is a unit. Thus $\ell_m$ is integral there.
At $m=q$, each coefficient in $[T^q]g^i$ for $i\ge2$ uses only
degrees smaller than $q$ and is divisible by $p^i$. Hence
$p\ell_q$ is integral and reduces to $\gamma_q$. The first possible
nonzero term of the reduced $p$-series has $p$-power degree, as proved
above. Its coefficient at $q$ vanishes exactly when the height is at
least $r+1$. Substituting $\ell_q=b_q/q$ and induction in $r$ proves
all the assertions, including the infinite case.
\end{proof}

This proof uses the original coordinate and does not discard the
non-$p$-power terms of its logarithm. For comparison, an integral
$p$-typicalization keeps the coefficients
$\lambda_i=b_{p^i}/p^i$. By
\citet[Theorem~2.8 and (4.3.1), pp.~134,138]{fg-hazewinkel},
its Hazewinkel parameters satisfy
\[
 v_n=p\lambda_n-\sum_{i=1}^{n-1}\lambda_i v_{n-i}^{p^i}\in R.
\]
In particular
\[
 v_1=b_p,\qquad v_2=\frac{b_{p^2}-b_p^{p+1}}p.
\]
They are parameters of the integral $p$-typical law, rather than all the
ordinary power-series coefficients of $[p]$ in the original coordinate.
The first nonzero one modulo $p$ detects the same height. The recursion
above is not needed to infer \eqref{fg:eq:next-height}.

\begin{corollary}[The second height invariant]\label{fg:cor:second}
On the locus $H_p(F)=0$, the element $b_{p^2}$ is divisible by $p$, and
\[
 h(X)\ge3\iff\frac{b_{p^2}}p=0\pmod p
            \iff b_{p^2}=0\pmod{p^2}.
\]
If $a_{p^2}$ denotes the literal coefficient of $T^{p^2}$ in $L$, this
invariant is $p a_{p^2}\bmod p$. The coefficient $a_{p^2}$ itself need
not be integral. For a fixed defining polynomial over $k$, the residue
$b_{p^2}/p\bmod p$ on this locus is independent of its homogeneous lift.
\end{corollary}

\begin{proof}
Only lift independence remains to be checked. Replace
$\widetilde F$ by $\widetilde F+pH$, with $H$ homogeneous of degree $d$.
Modulo $p^2$, the difference of the numerators is
\[
 p(p^2-1)[(x_1\cdots x_d)^{p^2-1}]H\widetilde F^{p^2-2}.
\]
Modulo $p$, the product after the coefficient sign is
$H F^{p-2}(F^{p-1})^p$. A contributing exponent of $HF^{p-2}$ is
congruent to $(p-1,\ldots,p-1)$ coordinatewise modulo $p$.
Its total degree is $d(p-1)$, so nonnegativity forces equality with
that vector. The coefficient is therefore
\[
 [(x_1\cdots x_d)^{p-1}]HF^{p-2}\, H_p(F)^p=0.
\]
This proves equality of the numerators modulo $p^2$, and hence of the
quotient residues. General height tests are intrinsic by their
equivalence to the reduced $p$-series. Under a nonstrict rescaling
$U=uT$, its first nonzero coefficient has weight $u^{1-p^h}$, so it is
its vanishing that is independent of every tangent normalization.
\end{proof}

\begin{example}[Why the second numerator needs two digits]\label{fg:ex:fermat}
For the Fermat cubic over $\F_2$, the integer lift gives $b_2=0$ and
$b_4=3!=6$. Thus $b_4/2=1\bmod2$ and its height is exactly two.
Both numerators vanish modulo two, so their bare reductions cannot
distinguish the two stages. Replacing the lift by $F+2xyz$ gives
$b'_2=2$ and $b'_4=14$, with the same second residue; the unrestricted
parameter is $(14-2^3)/2=3$.

For the Fermat quartic over $\F_3$, $b_3=0$ and
$b_9=8!/(2!)^4=2520$, divisible by nine. Its height is infinite.
Indeed, $b_{3^r}=0$ for odd $r$. For $r=2s$ put
$a=(3^{2s}-1)/4=2(1+9+\cdots+9^{s-1})$. Its base-three digit sum is
$2s=r$, and that of $3^r-1$ is $2r$. Legendre's factorial formula gives
\[
 v_3\left(\frac{(3^r-1)!}{(a!)^4}\right)=r.
\]
Equation~\eqref{fg:eq:infinite-height} applies at every stage.
These are smooth equations, since their coordinate derivatives have
no common projective zero. The multinomial formula also appears in
\citet[Example~4.13, p.~919]{fg-stienstra}.
\end{example}

\section{Higher heights as modular counting problems}\label{fg:sec:counting}
The next height test is a residue of a count associated with the incidence
presentation. It depends on more than the number of satisfying assignments.
Nevertheless, every fixed threshold at least two is modular-counting complete.
We prove the coefficient identities first, then establish membership for
arbitrary arithmetic circuits and give the reductions.

Throughout this section $p$ is fixed. For a homogeneous integer form $H$
of degree $u$ in $u$ variables, write
\[
 B_t(H)=[(x_1\cdots x_u)^{t-1}]H^{t-1}\qquad(t\ge1).
\]
When $H$ is a chosen lift of a defining equation, this is the logarithm
numerator $b_t$ of \eqref{fg:eq:stienstra}. The notation also applies to
forms whose zero schemes are singular or reducible; no formal group is
attached to such an intermediate form.

\subsection{Replica coefficients and the second invariant}
Let $\varphi$ be a width-three CNF with $S$ satisfying assignments, and use
the integer incidence polynomials $P,Z$ of \eqref{eq:incidence}. Every
declared variable occurs in an incidence, after the prescribed
tautological relations and, if necessary, a forced variable are added.
Both $P$ and $Z$ have degree $b$ in $b$ variables, and $[Z]P=S$.
A \emph{local configuration} chooses one allowed term in every factor
$P_f$ and one of the two terms in every factor $Q_i$. Distinct choices
remain distinct even when they produce the same monomial.

Regard \eqref{eq:normalization} as an integer polynomial $G_\varphi$ in
$u=b+p-1$ variables. We likewise use the displayed integer lift of
\eqref{eq:stabilizer} when its input is an integer form.

\begin{lemma}[Replica numerators]\label{fg:lem:replica}
Let $G\in\Z[x_1,\ldots,x_u]$ be homogeneous of degree $u\ge1$, and let
$Q$ be a positive multiple of $p$ with $Q\ge u+2$.
For every $1\le t\le Q$,
\begin{equation}\label{fg:eq:stabilized-numerator}
 b_t(F_G)=B_t(G).
\end{equation}
For $i\ge1$, set
\[
 k_i=\frac{p^i-1}{p-1},\qquad
 A_i(\varphi)=[Z^{k_i}]P^{k_i},\qquad
 C_i=(-1)^{k_i}\frac{(p^i-1)!}{(k_i!)^{p-1}}.
\]
If $Q\ge p^i$, then
\begin{equation}\label{fg:eq:replica-numerator}
 b_{p^i}(F_{G_\varphi})=C_iA_i(\varphi)
\end{equation}
over the integers. Moreover
\begin{equation}\label{fg:eq:replica-residue}
 v_p(C_i)=0,\qquad C_i\equiv1\pmod p,\qquad
 A_i(\varphi)\equiv S^i\pmod p.
\end{equation}
For each fixed $p,i$, the function $A_i$ belongs to $\#\mathrm P$.
\end{lemma}

\begin{proof}
Put $K=t-1$. In a term contributing to
$(x_1\cdots x_myz)^K$ in $F_G^K$, the summand $-x_1z^Q$ cannot occur,
since $Q>K$. The only remaining source of $z$ is $yzA_G$, whose
$z$-exponent is one. Thus all $K$ factors select $yzA_G$. The padding
coordinate $x_{u+1}$ exists because $Q\ge u+2$. Its exponent is one in
$G\prod_{j>u}x_j$ and zero in $x_1^{Q-1}$, so its target exponent forces
the first summand in every factor. Removing the padding gives
\eqref{fg:eq:stabilized-numerator}. This also treats the endpoint $Q=t$.

Now take $t=p^i$. Each tag variable occurs in one summand of
$G_\varphi$, with exponent $p-1$. Its target exponent $p^i-1$ forces
exactly $k_i$ copies of that summand. The multinomial coefficient is
$(p^i-1)!/(k_i!)^{p-1}$. The first summand contributes
$(-1)^{k_i}P^{k_i}$ and the remaining summands contribute
$Z^{(p-2)k_i}$. Removing them leaves the coefficient
$[Z^{k_i}]P^{k_i}$, proving \eqref{fg:eq:replica-numerator}. At $p=2$
there is only the first tag summand, with the same conclusion.

For an integer $a\ge0$, let $s_p(a)$ be its base-$p$ digit sum. Counting
the multiples of each power of $p$ in $a!$ gives
$v_p(a!)=(a-s_p(a))/(p-1)$. The digits of $p^i-1$ are all $p-1$,
whereas those of $k_i$ are all one. Hence
\[
 v_p((p^i-1)!)=k_i-i,\qquad
 (p-1)v_p(k_i!)=k_i-i.
\]
This proves the unit assertion. The digitwise multinomial expansion
gives the factorial quotient modulo $p$ as
$((p-1)!)^i=(-1)^i$. For odd $p$, $k_i\equiv i\pmod2$; at $p=2$
the sign is immaterial. Therefore $C_i\equiv1$.

In characteristic $p$,
\[
 P^{k_i}=\prod_{j=0}^{i-1}P(w^{p^j}).
\]
For a contribution to $Z^{k_i}$, every exponent in the lowest block is
congruent to one modulo $p$, and is therefore at least one. The block
has total degree $b$ in $b$ variables, so all its exponents are exactly
one. Subtract this exponent vector and divide the remaining target by
$p$. Induction forces the selected monomial in every block to be $Z$.
Its coefficient is $S$ modulo $p$, proving $A_i\equiv S^i$.

Finally, guess $k_i$ ordered local configurations and accept exactly
when their monomial product is $Z^{k_i}$. Each clause factor has at most
eight choices and each variable factor two. The exponent test and the
witness length are polynomial because $p,i$, and hence $k_i$, are
fixed. This counts every coefficient contribution with its local-choice
multiplicity, proving the $\#\mathrm P$ assertion.
\end{proof}

The same degree argument shows that for every homogeneous degree-$u$
form $H$ in $u$ variables,
\[
 B_{p^i}(H)\equiv B_p(H)^i\pmod p.
\]
Indeed, split $H^{p^i-1}$ into Frobenius twists of $H^{p-1}$. In the
lowest block every exponent is congruent to $p-1$ and hence at least
$p-1$; its total degree $u(p-1)$ forces equality in every coordinate.
No individual-exponent bound is needed for this congruence.

\begin{theorem}[The second invariant as an orbit count]\label{fg:thm:orbit}
There is an explicit function $\mathcal C\in\#\mathrm P$ such that
\begin{equation}\label{fg:eq:orbit-count}
 A_2(\varphi)=S^2+p\mathcal C(\varphi).
\end{equation}
It counts the nontrivial orbits of the cyclic group of order $p$ on
balanced ordered collections of $p+1$ local configurations, rotating
the first $p$ configurations and fixing the last. Balance means that
their monomial product is $Z^{p+1}$.

For the original formula construction, or any admissible choice
$Q\ge p^2$, the second invariant on $S\equiv0\pmod p$ is
\[
 \frac{b_{p^2}}p\equiv\mathcal C(\varphi)\pmod p.
\]
Consequently
\begin{equation}\label{fg:eq:height-three-count}
 h(X_\varphi)\ge3
 \iff S\equiv\mathcal C(\varphi)\equiv0\pmod p
 \iff A_2(\varphi)\equiv0\pmod{p^2}.
\end{equation}
\end{theorem}

\begin{proof}
The set of balanced ordered collections has cardinality $A_2$, since
$k_2=p+1$. Its cyclic action preserves balance. Since $p$ is prime,
each orbit has size one or $p$, and a fixed record has its first $p$
configurations identical.

Every incidence coordinate appears in exactly one clause factor and
one variable factor. Thus its exponent in a local configuration is
zero, one, or two. If the common first monomial has exponent vector
$a$ and the last has exponent vector $c$, balance gives
\[
 pa_j+c_j=p+1,\qquad 0\le a_j,c_j\le2.
\]
The only possibility is $a_j=c_j=1$: $a_j=0$ would require
$c_j=p+1>2$, and $a_j=2$ would require $c_j=1-p<0$.
Both selected monomials are therefore $Z$. Lemma~\ref{lem:incidence}
gives $S$ local configurations for each, so there are $S^2$ fixed
records. Counting the remaining size-$p$ orbits proves
\eqref{fg:eq:orbit-count}.

To realize $\mathcal C$ itself in $\#\mathrm P$, encode each local
configuration injectively in a fixed-length string by ordering its
factor choices. Guess a balanced collection, reject fixed records,
and accept exactly when its encoding is strictly lexicographically
smaller than each of its nonidentity rotations. Every nontrivial orbit
has exactly one such record. The test is polynomial time for fixed $p$.
This is a canonical-representative count, not an assumed closure under
division by $p$.

The original choice of $Q$ is a power of $p$ with $Q-1\ge p$, hence
$Q\ge p^2$. Lemma~\ref{fg:lem:replica} gives $b_p=C_1S$ and
$b_{p^2}=C_2A_2$. If $p\mid S$, then
\[
 \frac{b_{p^2}}p
 =C_2\left(\frac{S^2}p+\mathcal C(\varphi)\right)
 \equiv\mathcal C(\varphi)\pmod p,
\]
using $p^2\mid S^2$ and $C_2\equiv1$. Apply
Corollary~\ref{fg:cor:second}. Finally, $p^2\mid A_2$ already implies
$S^2\equiv0\pmod p$ by \eqref{fg:eq:replica-residue}, so it forces
$p\mid S$. Equation~\eqref{fg:eq:orbit-count} then gives the remaining
equivalence in \eqref{fg:eq:height-three-count}.
\end{proof}

\begin{example}[Equivalent formulas with different second heights]
\label{fg:ex:same-function}
Work over $\F_2$ with the three declared Boolean variables
$x_0,x_1,x_2$. Consider
\begin{align*}
 \varphi_0={}&x_1\wedge x_2,\\
 \varphi_1={}&(\neg x_0\vee\neg x_1\vee x_2)
 \wedge(x_0\vee\neg x_1\vee x_2)\\
 &{}\wedge(\neg x_0\vee x_2)\wedge x_1.
\end{align*}
The unit clause forces $x_1=1$. The first two clauses of $\varphi_1$
then force $x_2=1$, and the third is redundant. Thus both formulas
compute the same Boolean function and have $S=2$.

Here $k_2=3$. We compute the coefficients by local marginal counts.
For a clause $f$, let $N_f((a_i)_{i\in I_f})$ count its ordered triples
of satisfying assignments having exactly $a_i$ ones at variable $i$.
In $Q_i^3$, choosing the zero-labelled monomial $a_i$ times has
coefficient $\binom3{a_i}$. Balance at each incidence forces the same
ones count $a_i$ in its clause triples. Therefore
\[
 A_2(\varphi)=
 \sum_{0\le a_i\le3}
   \prod_i\binom3{a_i}
   \prod_fN_f((a_i)_{i\in I_f}).
\]

For $\varphi_0$, the unit clauses force $a_1=a_2=3$. The required
tautological unary relation for the unused variable $x_0$ contributes
$\binom3{a_0}$, in addition to the same factor from $Q_0^3$. Hence
\[
 A_2(\varphi_0)=\sum_{a=0}^3\binom3a^2
 =1+9+9+1=20.
\]

For $\varphi_1$, put $a=a_0$, $b=a_2$, and again $a_1=3$.
Two clause relations are $\neg x_0\vee x_2$ and one is $x_0\vee x_2$.
Their ordered-triple counts are
\[
 N_-(a,b)=\binom3b\binom ba,\qquad
 N_+(a,b)=\binom3b\binom b{a+b-3},
\]
with out-of-range binomial coefficients zero. The first formula chooses
the $b$ positions of $x_2=1$ and the $a$ positions of $x_0=1$ within
them. For the second, $x_0$ must equal one wherever $x_2$ is zero.
Thus
\[
 A_2(\varphi_1)=
 \sum_{a,b=0}^3\binom3a\binom3b\,N_-(a,b)^2N_+(a,b).
\]
Exactly six summands are nonzero:
\[
 \begin{array}{c|r}
 (a,b)&\binom3a\binom3b\,N_-(a,b)^2N_+(a,b)\\ \hline
 (1,2)&972\\
 (2,2)&486\\
 (0,3)&1\\
 (1,3)&81\\
 (2,3)&81\\
 (3,3)&1
 \end{array}
\]
Their sum is $1622$. In particular,
\[
 \mathcal C(\varphi_0)=8,\qquad
 \mathcal C(\varphi_1)=809.
\]
Theorem~\ref{fg:thm:orbit} and \eqref{fg:eq:exact-height} give height
at least three for the first constructed group and height exactly two
for the second. The second invariant is therefore not determined by
$S$, or even by the Boolean function.

The example can use a common ambient space. The incidence counts are
three and nine, so the tagged input dimensions are seven and nineteen.
The common choice $Q=32$ is admissible for both. Their stabilized
equations then both have degree $33$ in $\PP^{32}$.
\end{example}

\subsection{Counting coefficient digits for arbitrary circuits}
\begin{lemma}[Fixed prime-power tests]\label{fg:lem:digits}
Let $f$ be a $\#\mathrm P$ function, and fix an integer $a\ge1$.
The function
\[
 \Delta_a(f)=
 \prod_{j=0}^{a-1}
 \left(1+(p-1)\binom{f}{p^j}^{p-1}\right)
\]
belongs to $\#\mathrm P$ and satisfies
\[
 \Delta_a(f)\equiv
 \begin{cases}
 1,&p^a\mid f,\\
 0,&p^a\nmid f
 \end{cases}\pmod p.
\]
Fixed finite conjunctions and complements of these tests also have
such zero-or-one counting representatives modulo $p$.
\end{lemma}

\begin{proof}
Encode accepting paths by witness strings of one polynomial length,
fixing unused positions to zero. For a fixed integer $q\ge0$,
guessing $q$ accepting strings in strictly increasing lexicographic
order gives exactly $\binom{f}{q}$ accepting tuples. Thus this
binomial is a $\#\mathrm P$ function.

Write the nonnegative integer $f$ in base $p$ as
$f=\sum_j f_jp^j$. In characteristic $p$,
\[
 (1+T)^f=\prod_j(1+T^{p^j})^{f_j}.
\]
The lower-index factors have combined degree less than $p^j$, and
higher-index factors have no positive term of degree at most $p^j$.
Comparison at degree $p^j$ therefore gives
$\binom{f}{p^j}\equiv f_j\pmod p$.

For any accepting-path count $c$, the nonnegative count
$1+(p-1)c^{p-1}$ is one modulo $p$ when $c$ is zero, and zero
otherwise. Products and sums are realized by tuples of accepting
computations and disjoint labelled branches. Their number is fixed
here. Applying this zero test to all the first $a$ digits proves the
formula. Products of zero-or-one residues give conjunctions, and
$t\mapsto1+(p-1)t$ complements such a residue. No count division is
used.
\end{proof}

\begin{lemma}[Coefficient witnesses for circuit input]
\label{fg:lem:circuit-coeff}
Fix $J\ge1$. Suppose an arithmetic circuit over $\Fp$ is promised to
define a homogeneous degree-$d$ polynomial in $d\ge3$ variables whose
projective hypersurface is smooth and geometrically integral.
There is a deterministic polynomial-time construction of a circuit for a
homogeneous $p$-primitive integer lift $\widetilde F$ with nonnegative coefficients
such that each function
\[
 B_{p^j}(\widetilde F)
 =[(x_1\cdots x_d)^{p^j-1}]
       \widetilde F^{p^j-1},\qquad 1\le j\le J,
\]
is represented by a polynomial-time counting verifier. In particular
these are $\#\mathrm P$ functions of the circuit input.
\end{lemma}

\begin{proof}
First, the numerical dimension is polynomially bounded on the promise.
If some variable is absent, its coordinate point lies on the
hypersurface and every first derivative vanishes there, since $d\ge3$.
This would be a singular point. Thus every ambient variable occurs
syntactically. Even with binary coordinate labels, $d$ is at most the
number of distinct variable inputs, hence at most the input length.
A total verifier may reject larger invalid declarations before any
dimension-dependent construction.

Replace subtraction by addition and multiplication by $p-1$. If
needed, homogenize through degree $d$: retain degrees zero through $d$
at every gate, and use convolution of homogeneous components at a
multiplication gate. This uses polynomially many gates, without
expanding higher-degree components. Evaluate each degree-zero
subcircuit in $\Fp$ and replace it by its representative in
$\{0,\ldots,p-1\}$. Interpret the remaining component circuit over the
nonnegative integers and take its degree-$d$ output. This defines
$\widetilde F$ and reduces to $F$ modulo $p$. It is $p$-primitive, meaning that some coefficient is not divisible by $p$,
because $F$ is nonzero.

Folding the degree-zero subcircuits is necessary for the witness
bound. Repeated squaring of an integer constant could otherwise give
a double-exponentially large coefficient even in a circuit of bounded
positive degree. We are choosing a lift of the given polynomial over
$\Fp$, so such integer evaluations need not be preserved. Every
nonzero scalar leaf of the chosen lift is bounded by $p-1$.

Consider now a homogeneous nonnegative circuit with $M$ gates and
output degree $D$, with degree-zero subcircuits already replaced by
scalar leaves. A distributive parse selects one input at an addition
gate and both inputs at a multiplication gate. A contributing parse
has exactly $D$ variable leaves. Every multiplication with two
positive-degree inputs splits positive degree, so there are at most
$D-1$ such branchings. Between them, a path has length at most $M$.
Scalar children are leaves. Thus the entire unfolded parse has
$O(MD)$ vertices, despite sharing in the original circuit.

A nondeterministic verifier follows the parse in a fixed depth-first
order. At a scalar leaf $c$, it makes exactly $c$ accepting labelled
choices, rejecting unused binary labels; zero rejects. At variable
leaves it accumulates the exponent vector, and at the end it checks
the target monomial. This counts every distributive contribution with
its coefficient multiplicity, in polynomial time.

Apply the construction to $\widetilde F^{p^j-1}$. For fixed $p,J$,
its circuit size and degree $D=d(p^j-1)$ are polynomial in the input
length. The same chosen lift works for all $j\le J$.
Lemma~\ref{fg:lem:height-log} applies to that lift; no specified
integer evaluation of the original circuit is part of the input.
\end{proof}

\subsection{Fixed thresholds and exact-height seeds}
\begin{theorem}[Complexity of every fixed height threshold]
\label{fg:thm:threshold}
Fix a prime $p$ and an integer $r\ge2$. The following promise problem
is $\ModP$-complete under deterministic polynomial-time many-one
reductions: given a homogeneous arithmetic circuit over $\Fp$ for a
degree-$d$ polynomial in $d\ge3$ variables, promised to define a smooth
geometrically integral hypersurface $X\subset\PP^{d-1}$, decide whether
the top Artin--Mazur formal group has $h(X)\ge r$.
Infinite height is included in this predicate. The dimension is
allowed to grow, and hardness persists after imposing any fixed lower
bound on it.

For every fixed finite $h\ge1$, the exact-height predicate $h(X)=h$
on the same circuit promise belongs to $\ModP$.
\end{theorem}

\begin{proof}
Use the lift of Lemma~\ref{fg:lem:circuit-coeff}, taking $J=r-1$ for the
threshold predicate and $J=h$ for the exact-height predicate, and put
$f_j=B_{p^j}(\widetilde F)$. By \eqref{fg:eq:all-height},
height at least $r$ is equivalent to $p^j\mid f_j$ for $1\le j<r$.
Lemma~\ref{fg:lem:digits} supplies the accepting-path count
\[
 \prod_{j=1}^{r-1}\Delta_j(f_j),
\]
whose residue is one exactly on those inputs. Likewise
\[
 \left(\prod_{j=1}^{h-1}\Delta_j(f_j)\right)
       \left(1+(p-1)\Delta_h(f_h)\right)
\]
represents exact height $h$ by \eqref{fg:eq:exact-height}, with an
empty product equal to one. These total counting machines need only
agree on the promise; they do not recognize smoothness.

For hardness, start with a CNF $\varphi$ and take $a=r-1$ copies of
$G_\varphi$ in disjoint variable blocks. Their product
\[
 G^{(a)}=\prod_{\nu=1}^aG_\varphi(x^{(\nu)})
\]
has degree $au$ in $au$ variables. Apply Theorem~\ref{thm:stabilizer}
with a $p$-power $Q\ge\max\{au+2,p^{r-1}\}$.
It gives a smooth geometrically integral anticanonical hypersurface.
Coefficients separate across the blocks, and
Lemma~\ref{fg:lem:replica} yields
\[
 b_{p^j}(F_{G^{(a)}})=
       \bigl(C_jA_j(\varphi)\bigr)^a\qquad(1\le j<r).
\]
If $S\equiv0\pmod p$, every $A_j$ is divisible by $p$, so these
numerators have valuation at least $a=r-1\ge j$. The output height
is at least $r$. If $S\not\equiv0\pmod p$, its first numerator
$(C_1S)^a$ is a unit, so the height is one. Thus
\[
 h(F_{G^{(a)}})\ge r\iff S\equiv0\pmod p.
\]
The product input itself need not be smooth; its stabilization is the
promised output.

To give the reduction in the nonzero-residue convention for $\ModP$,
let $q$ be any source accepting-path count and form
$q'=1+(p-1)q^{p-1}$. This is a $\#\mathrm P$ count and
$q'\equiv0\pmod p$ exactly when $q\not\equiv0\pmod p$.
The parsimonious Boolean encoding used for Theorem~\ref{thm:main}
gives a width-three CNF with exactly $q'$ satisfying assignments.
Applying the preceding construction gives a many-one reduction in
the required direction.

At fixed $p,r$, the number of blocks and the lower bound $p^{r-1}$
are constants. The least admissible $p$-power $Q$ is less than $p$
times its numerical threshold, so the circuit size is polynomial.
Choosing $Q$ also to exceed any fixed dimension bound preserves
this cost.
\end{proof}

\begin{theorem}[Exact height from one fixed seed]\label{fg:thm:exact-seed}
Fix a prime $p$ and a finite integer $h\ge1$. If there exists a smooth
geometrically integral anticanonical hypersurface over $\Fp$ of
positive dimension whose top formal group has height $h$, then
deciding exact height $h$ in the growing-dimensional circuit class
of Theorem~\ref{fg:thm:threshold} is $\ModP$-complete. If no such
hypersurface exists, its promise yes-locus is empty.
\end{theorem}

\begin{proof}
Fix one seed equation $H$ of degree $e$ in $e\ge3$ variables and a
homogeneous $p$-primitive integer lift $\widetilde H$. Its equation is
fixed data, because $p,h$, and the seed are fixed. Write
$\beta_j=B_{p^j}(\widetilde H)$. Equations
\eqref{fg:eq:all-height}--\eqref{fg:eq:exact-height} give
\[
 p^j\mid\beta_j\quad(1\le j<h),\qquad
 v_p(\beta_h)=h-1.
\]

Given $\varphi$, form $G=\widetilde H(v)G_\varphi(w)$ in disjoint
variable blocks. It has degree $e+u$ in $e+u$ variables. Stabilize
its reduction modulo $p$ with a $p$-power
$Q\ge\max\{e+u+2,p^h\}$, using the displayed integer form for its
lift. Coefficient separation and
\eqref{fg:eq:stabilized-numerator} give
\[
 b_{p^j}(F_G)=\beta_jC_jA_j(\varphi)\qquad(1\le j\le h).
\]
All prior divisibilities hold for every input. At the $h$th stage,
$\beta_h$ has valuation $h-1$, $C_h$ is a unit, and
$A_h(\varphi)\equiv S^h\pmod p$. Consequently
\[
 h(F_G)=h\iff S\not\equiv0\pmod p.
\]
When $S\equiv0$, the output instead has height at least $h+1$,
possibly infinite. The map is polynomial time with the fixed seed
hard-coded, and every output satisfies the geometric promise.
This proves hardness; membership is
Theorem~\ref{fg:thm:threshold}.
\end{proof}

Theorem~\ref{fg:thm:all-heights} below supplies a seed over $\Fp$
for every finite height. Corollary~\ref{fg:cor:exact-height-complete}
therefore makes exact-height completeness unconditional.
The counting algorithms above treat $p$ and the finite threshold as
fixed parameters. They assert neither polynomial dependence on a
binary-encoded variable threshold nor a finite stopping rule for the
infinite list in \eqref{fg:eq:infinite-height}.

\section{Explicit prime-field seeds of every finite height}\label{fg:sec:seeds}
The seed required by Theorem~\ref{fg:thm:exact-seed} exists for every
prime and every finite height. We first construct an integer form whose
relevant diagonal coefficient is a single multinomial. The coefficient
preservation of Lemma~\ref{fg:lem:replica} then gives a smooth
anticanonical hypersurface over the prime field. We use the notation
$B_t$ of Section~\ref{fg:sec:counting}; the intermediate form need not
define a smooth variety, and we attach no formal group to it.

\subsection{An integer coefficient seed}
Fix a prime $p$ and an integer $h\ge2$, and put
\[
 N=p^h-1,\qquad s=(2h-1)(p-1).
\]
Let $\nu=(\nu_1,\ldots,\nu_s)$ be the nondecreasing list containing
\[
 \begin{array}{c|c}
 \text{entry}&\text{multiplicity}\\ \hline
 1&2p-1\\
 p^k,\quad 1\le k\le h-2&2p-2\\
 p^{h-1}&p-2 .
 \end{array}
\]
The middle range is empty when $h=2$, and the final multiplicity is
zero when $p=2$. Distinguish $\nu_1=1$, and write
\[
 w_j=\nu_{j+1}\quad(1\le j\le s-1),\qquad
 S_0=0,\quad S_j=w_1+\cdots+w_j.
\]
We have
\begin{equation}\label{fg:eq:seed-weight}
 \sum_{j=1}^s\nu_j
 =(2p-1)+2(p-1)\sum_{k=1}^{h-2}p^k+(p-2)p^{h-1}
 =N,\qquad S_{s-1}=N-1.
\end{equation}

For every integer $0\le t\le N-1$, define a binary vector
$u_t\in\{0,1\}^{s-1}$ as follows. If $t=S_j$, let $u_t$ have its first
$j$ entries equal to one and its remaining entries zero. At all other
indices, apply descending greedy selection to the list $w$: start with
residual $t$, visit the indices $s-1,s-2,\ldots,1$, and select $w_j$
exactly when it does not exceed the current residual, subtracting it
when selected. The selected indices are the one entries of $u_t$.
The proof below verifies that the residual ends at zero.

In $N$ variables define
\begin{equation}\label{fg:eq:seed-monomials}
 M_1=x_1^N,\qquad
 M_{j+1}=\prod_{t=1}^{N-1}
       x_{t+1}^{\,1+u_{t,j}-u_{t-1,j}}\quad(1\le j\le s-1),
 \qquad G_{p,h}=\sum_{j=1}^sM_j.
\end{equation}

\begin{proposition}[Exact diagonal coefficients of the seed]
\label{fg:prop:coefficient-seed}
The construction satisfies $w\cdot u_t=t$ for $0\le t\le N-1$.
The monomials in \eqref{fg:eq:seed-monomials} are pairwise distinct
of degree $N$, and every exponent in $M_2,\ldots,M_s$ is at most two.
Thus $G_{p,h}$ is a degree-$N$ integer form in $N$ variables, with
exactly $s$ terms, all with coefficient one.

For every integer $r\ge0$,
\begin{equation}\label{fg:eq:seed-period}
 B_{r+1}(G_{p,h})=
 \begin{cases}
 0,&N\nmid r,\\[1mm]
 \displaystyle\frac{(aN)!}{\prod_{j=1}^s(a\nu_j)!},
       &r=aN,\quad a\ge0.
 \end{cases}
\end{equation}
In particular,
\begin{equation}\label{fg:eq:seed-height-data}
 B_{p^i}(G_{p,h})=0\quad(1\le i<h),\qquad
 v_p\bigl(B_{p^h}(G_{p,h})\bigr)=h-1.
\end{equation}
The form has a division-free circuit with
$O(Ns+\log N)$ addition and multiplication gates, and its full exponent
table and such a circuit can be generated in time polynomial in
$N,s,\log N$.
\end{proposition}

\begin{proof}
The length of the prescribed list is
\[
 (2p-1)+(h-2)(2p-2)+(p-2)=(2h-1)(p-1)=s,
\]
and summing its entries gives \eqref{fg:eq:seed-weight}.
The remaining list $w$ has $2p-2$ unit entries. Before the first entry
of weight $p^k$, its lower-weight entries have total
\[
 2(p-1)(1+p+\cdots+p^{k-1})=2(p^k-1).
\]
Consequently, each entry of this sorted list is at most one plus the
sum of all preceding entries.

For any sorted positive integer list satisfying this condition and
starting with one, descending greedy selection represents every
integer between zero and the total weight. Indeed, if a target is at
least the final entry, selecting that entry leaves a nonnegative
target at most the preceding total. If it is smaller than the final
entry, it is already at most the preceding total, by the displayed
condition. Induction proves the assertion. The prescribed prefix
vectors also have their required weighted sums, and the positive
weights make all prefix indices distinct. Thus
\[
 w\cdot u_t=t,\qquad u_0=0,\qquad
 u_{N-1}=(1,\ldots,1).
\]

Each exponent in $M_{j+1}$ lies in $\{0,1,2\}$, and telescoping gives
\[
 \deg M_{j+1}
 =\sum_{t=1}^{N-1}(1+u_{t,j}-u_{t-1,j})
 =N-1+u_{N-1,j}-u_{0,j}=N.
\]
Only $M_1$ involves $x_1$. If two other exponent columns, indexed by
$j<k$, were equal, their successive differences would give
$u_{t,j}=u_{t,k}$ for every $t$, starting from $u_0=0$. At $t=S_j$,
however, the prescribed prefix vector has $u_{t,j}=1$ and $u_{t,k}=0$.
This proves pairwise distinctness and the coefficient-one assertion.

Consider a term contributing to
$(x_1\cdots x_N)^r$ in $G_{p,h}^r$, and let $n_j$ be the number of
selected copies of $M_j$. Its $x_1$ exponent is $Nn_1$. Therefore
there is no contribution unless $r=aN$ for an integer $a\ge0$.
In that case
\[
 n_1=a,\qquad \sum_{j=1}^{s-1}n_{j+1}=a(N-1).
\]
Write $n'=(n_2,\ldots,n_s)$. The required exponent $aN$ at
$x_{t+1}$ gives
\[
 a(N-1)+(u_t-u_{t-1})\cdot n'=aN,
 \qquad (u_t-u_{t-1})\cdot n'=a.
\]
Summing from $u_0=0$ yields $u_t\cdot n'=at$. At the prefix indices
$t=S_j$, this is
\[
 n_2+\cdots+n_{j+1}=aS_j.
\]
Successive differences force $n_{j+1}=aw_j$, so the only possible
count vector is $n=a\nu$. Conversely, this vector has the required
total and all the required exponents because $w\cdot u_t=t$.
The multinomial theorem therefore proves the exact integer identity
\eqref{fg:eq:seed-period}, including its $a=0$ case.

For $1\le i<h$, we have $0<p^i-1<N$, so
$B_{p^i}(G_{p,h})=0$. For $i=h$, the coefficient is
$N!/\prod_j\nu_j!$. Legendre's formula, already used in
Lemma~\ref{fg:lem:replica}, gives
\[
 v_p\left(\frac{N!}{\prod_j\nu_j!}\right)
 =\frac{\sum_j s_p(\nu_j)-s_p(N)}{p-1}
 =\frac{(2h-1)(p-1)-h(p-1)}{p-1}=h-1.
\]
Here each $\nu_j$ is a power of $p$, and the $h$ base-$p$ digits of
$N=p^h-1$ are all $p-1$. This proves
\eqref{fg:eq:seed-height-data}.

Finally, there are $N$ vectors $u_t$, each obtained by at most
$s-1$ comparisons and subtractions on $O(\log N)$-bit integers.
The exponent table follows by adjacent differences. Precompute the
squares $x_i^2$; each of the last $s-1$ monomials then uses $O(N)$
multiplications, while $x_1^N$ uses $O(\log N)$ by repeated squaring.
Adding the terms proves the gate and construction bounds.
\end{proof}

For instance, $p=2,h=2$ gives
$G_{2,2}=x_1^3+x_2^2x_3+x_2x_3^2$, with $B_2=0$ and $B_4=6$.
At $p=2,h=3$, the multiplicity list is $(1,1,1,2,2)$, and
$B_8=7!/(2!2!)=1260$ has valuation two. These values illustrate the
integer identity; no smoothness claim about an intermediate form is
needed.

\subsection{Smooth seeds and exact-height completeness}
\begin{theorem}[Every finite height over the prime field]
\label{fg:thm:all-heights}
For every prime $p$ and every finite integer $h\ge1$, there is an
explicit smooth geometrically integral anticanonical projective
hypersurface over $\Fp$ whose top Artin--Mazur formal group has exact
height $h$. One may take its degree and number of homogeneous
variables to be
\begin{equation}\label{fg:eq:smooth-seed-size}
 d_{p,h}=p^h+p+1.
\end{equation}
Its dimension is then $p^h+p-1$.
\end{theorem}

\begin{proof}
First assume $h\ge2$ and take the form $G_{p,h}$ above, in
$N=p^h-1$ variables. Choose
\[
 Q=p^h+p.
\]
This is a positive multiple of $p$ with
$Q\ge N+2$ and $Q\ge p^h$. Theorem~\ref{thm:stabilizer} applies to
the reduction of $G_{p,h}$ over $\Fp$ and gives a smooth
geometrically integral hypersurface $X=V(F_{G_{p,h}})$ of degree
$Q+1$ in $Q+1$ variables. Thus it is anticanonical and has the
degree and dimension in \eqref{fg:eq:smooth-seed-size}.
The theorem allows every positive multiple $Q$ of $p$ satisfying
the size condition; $Q$ need not be a power of $p$.

Use the displayed integer stabilizer equation as the homogeneous
lift. Lemma~\ref{fg:lem:replica} gives exact equalities over the
integers
\[
 b_{p^i}(F_{G_{p,h}})=B_{p^i}(G_{p,h})
 \qquad(1\le i\le h).
\]
The earlier numerators vanish and the last has valuation $h-1$ by
\eqref{fg:eq:seed-height-data}. The integral Stienstra logarithm
and \eqref{fg:eq:exact-height} therefore give $h(X)=h$.
For $h=1$, use instead the degree-one form $G=x_1$ in one
variable and choose $Q=2p$. Then $Q\ge3$ and $Q\ge p$, so the
same stabilization and coefficient identity apply. Since
$B_p(x_1)=1$, the resulting smooth hypersurface has $b_p=1$ and
height one by Proposition~\ref{fg:prop:height-one}. Its degree and
variable count are $Q+1=2p+1=p^1+p+1$ as well. The formal-group
interpretation is used only after stabilization to positive dimension.

The construction is explicit by Proposition~\ref{fg:prop:coefficient-seed}
and the circuit construction in Theorem~\ref{thm:stabilizer}.
For each fixed $p,h$, its equation is fixed finite data.
\end{proof}

\begin{corollary}[Completeness at every fixed exact height]
\label{fg:cor:exact-height-complete}
Fix a prime $p$ and a finite integer $h\ge1$.
For division-free arithmetic circuits over $\Fp$ promised to define
smooth geometrically integral anticanonical hypersurfaces of positive
dimension, with dimension allowed to grow, deciding $h(X)=h$ is
$\ModP$-complete under deterministic polynomial-time many-one
reductions. Hardness persists after imposing any fixed lower bound
on the dimension.
\end{corollary}

\begin{proof}
Theorem~\ref{fg:thm:all-heights} supplies the fixed seed in
Theorem~\ref{fg:thm:exact-seed}, which proves completeness.
In its reduction, the smoothing parameter can also exceed any
prescribed fixed dimension bound. This preserves all coefficient
identities through $p^h$ and the polynomial circuit-size bound.
\end{proof}

The degree in \eqref{fg:eq:smooth-seed-size} is exponential in $h$.
The completeness statement fixes $p$ and $h$; it does not assert a
polynomial-size seed as a function of a binary-encoded varying height.

For the literature comparison,
\citet[Convention~4.15 and Theorem~4.18, pp.~25--28]{fg-kty-two}
construct anticanonical hypersurfaces of every finite height over an
algebraically closed field, using degree $p^h-1$ for $h>1$.
That proof intersects nonempty smoothness and height loci in a
parameter space. Its stated field convention does not supply a
prime-field point in that intersection. The coefficient identity
\eqref{fg:eq:seed-period} and the prescribed stabilization above
give equations over $\Fp$, with the explicit larger degree
\eqref{fg:eq:smooth-seed-size}.

\section{Computing the full height in polynomial space}\label{fg:sec:space}
Fixed thresholds use a constant number of Stienstra coefficients. Infinite
height requires a finite stopping argument. A finite operator from the
quasi-$F$-splitting criterion supplies one, and its entries can be computed
from a circuit without storing its exponentially large matrix.

We use the following theorem in its established geometric form.
\citet[Theorem~C, p.~4, and Theorem~5.8]{fg-kty-one} give a Fedder type
criterion for quasi-$F$-split height. For smooth Calabi--Yau varieties this
height agrees with Artin--Mazur height by
\citet[Theorem~4.5, pp.~9--10 of the preprint]{fg-yobuko}.
That source assumes an algebraically closed ground field. Apply its
comparison after extension to $\overline{\Fp}$: the first nonzero
coefficient of a formal group's $p$-series is unchanged by field
extension, and the coefficient recurrence below for an equation over
$\Fp$ stays in $\Fp$. Its first nonzero index is therefore unchanged
as well; see also \citet[Theorem~5.13]{fg-kty-one}.
The concrete monomial formula and finite operator
are given by \citet[Proposition~3.2, Remark~3.4 and Algorithms~2--4,
pp.~5--7]{fg-bgp}. We recall the exact criterion and prove its complexity
consequences for circuit input.

Let $F\in\Fp[x_1,\ldots,x_d]$ be homogeneous of degree $d\ge3$, and
assume its projective hypersurface is smooth and geometrically connected.
Put
\[
 D=d(p-1),\qquad g=F^{p-1},\qquad
 V=\Fp[x_1,\ldots,x_d]_D,\qquad
 B=\dim V=\binom{pd-1}{d-1}.
\]
For the collected monomial expansion $g=\sum_\beta c_\beta x^\beta$,
choose integer representatives $\widehat c_\beta\in\{0,\ldots,p-1\}$,
set $G=\sum_\beta\widehat c_\beta x^\beta$, and define
\begin{equation}\label{fg:eq:delta}
 \Delta=\left(\frac{G^p-\sum_\beta
                    (\widehat c_\beta x^\beta)^p}{p}\right)\bmod p.
\end{equation}
The numerator is coefficientwise divisible by $p$. Changing the integer
representatives changes both $p$-th powers by multiples of $p^2$, so
the result is well defined. Define the $\Fp$-linear extractor $u$ by
\[
 u(x^\gamma)=
 \begin{cases}
 x^{(\gamma-(p-1)\mathbf1)/p},
       &\gamma\equiv(p-1)\mathbf1\pmod p,\\
 0,&\text{otherwise}.
 \end{cases}
\]
Nonnegative exponents in the first case are automatic. Let
$T(a)=u(\Delta a)$ and let $\ell$ select the coefficient with exponent
$\tau=(p-1)\mathbf1$. Since $\deg\Delta=pD$, the operator $T$ preserves
$V$: extracting from degree $(p+1)D$ subtracts $d(p-1)$ and divides by
$p$, returning degree $D$.
The cited criterion is
\begin{equation}\label{fg:eq:finite-operator}
 h(X)=\min\{j+1:\ell(T^jg)\ne0\},
 \qquad \min\varnothing=\infty.
\end{equation}
Indeed, the unique degree-$D$ monomial outside
$(x_1^p,\ldots,x_d^p)$ is $x^\tau$.

The representative convention in \eqref{fg:eq:delta} is compatible with
the Frobenius-image convention for $\Delta_1$ in the criterion.
Grouping monomials with the same exponent residue modulo $p$ changes
the divided difference only by a $p$-th-power polynomial. Replacing
$\Delta$ by $\Delta+H^p$ changes $T(a)$ by $H u(a)$.
For $a\in V$, $u(a)=\ell(a)$ is constant. Thus the two recurrences
agree until their first nonzero target, and agree forever when every
target is zero. This proves the required compatibility for both finite
and infinite height.

\begin{theorem}[A space bound for the full height]\label{fg:thm:space}
Fix $p$. Let $L$ be the complete bit length of a division-free arithmetic
circuit over $\Fp$, with addition, subtraction and multiplication gates,
and with its declared variable count $d$ and all indices encoded in
binary. Promise that it computes a homogeneous degree-$d$ equation in
$d\ge3$ variables defining a smooth geometrically connected projective
hypersurface. There is a deterministic algorithm which outputs its
finite Artin--Mazur height in binary, or outputs infinity, using
\[
 O_p\bigl(L^2\log(L+2)\bigr)
\]
working bits. Every finite height is at most $B=\binom{pd-1}{d-1}$.
The algorithm halts on all syntactically valid inputs and has arbitrary
output outside the promise. Infinite-height decision, and finite-height
or threshold queries with a binary-encoded queried height, therefore
belong to polynomial space.
\end{theorem}

\begin{proof}
Cayley--Hamilton gives a recurrence of order $B$ for
$a_j=\ell(T^jg)$. If $a_0,\ldots,a_{B-1}$ all vanish, every subsequent
term vanishes. Otherwise the first nonzero term has index below $B$.
Together with \eqref{fg:eq:finite-operator}, this proves the height
bound and the finite stopping rule. No characteristic polynomial needs
to be computed.

We give the storage analysis in detail. Every declared variable occurs
on a promised input, by the singular-coordinate-vertex argument in
Proposition~\ref{fg:prop:height-one}; hence $d=O(L)$. A total algorithm
first rejects $d<3$ or a declared dimension larger than its number of
variable occurrences. The bound $B\le2^{pd-1}$ gives $\log B=O_p(d)$. The integer
$B$ can be computed by successive exact divisions in the binomial
formula, using $O_p(d)$ bits. A degree-$D$ basis monomial is represented
by $d$ exponents of size at most $D$, using $O_p(d\log(d+2))$ bits.
Enumerate these weak compositions in a fixed order, with one such
counter, whenever a sum over the basis is required.

First compute one circuit coefficient. For a gate $v$ and exponent
vector $\alpha$ of total degree $O_p(d)$, write $C(v,\alpha)$ for its
coefficient modulo $p$. Constants, variables, sums and differences are
immediate. At a product gate with predecessors $v_1,v_2$, use
\[
 C(v,\alpha)=\sum_{0\le\beta\le\alpha}
           C(v_1,\beta)C(v_2,\alpha-\beta).
\]
Evaluate the summands and the two factors sequentially. Each recursive
frame stores its gate index, a constant number of exponent vectors,
the current summation vector and field elements. Its space is
$O_p(d\log(d+2)+\log(L+2))$. Gate indices decrease along every recursive
path, so at most $O(L)$ frames are active. The resulting space bound is
\begin{equation}\label{fg:eq:coefficient-space}
 O_p\bigl(Ld\log(d+2)+L\log(L+2)\bigr).
\end{equation}
This procedure is exact even if intermediate gates have very large
formal degree: a term contributing to $x^\alpha$ at a product uses
only componentwise smaller nonnegative exponents. Arithmetic remains
modulo the fixed prime. Recomputed subcircuits cost time but no table
of coefficients. A circuit for $g=F^{p-1}$ has $O_p(L)$ gates, so this
also computes each $c_\beta$.

Next compute one coefficient of $\Delta$. For $|\gamma|=pD$, enumerate
the ordered $p$-tuples of degree-$D$ monomials whose exponents sum to
$\gamma$. Accumulate their products of canonical coefficients
$\widehat c_\beta$ modulo $p^2$. Subtract
$\widehat c_{\gamma/p}^{p}$ if every coordinate of $\gamma$ is divisible
by $p$, and subtract zero otherwise. The resulting canonical residue
in $\{0,\ldots,p^2-1\}$ is divisible by $p$; divide it by $p$ and reduce
modulo $p$. This is exactly \eqref{fg:eq:delta}. Since $p$ is fixed,
the tuple counters add only $O_p(d\log(d+2))$ bits to
\eqref{fg:eq:coefficient-space}. The computation uses the coefficients
of the collected polynomial $g$, so no uncollected parse-tree
multiplicity is mistaken for a monomial coefficient.

Let $M$ be the matrix of $T$ in the degree-$D$ monomial basis. Its row
$\alpha$, column $\beta$ entry is
\begin{equation}\label{fg:eq:matrix-entry}
 M_{\alpha\beta}
 =[x^{p\alpha+(p-1)\mathbf1-\beta}]\Delta,
\end{equation}
with value zero if a coordinate is negative. The target has total
degree $pD$, so the preceding procedure supplies this entry in the
same space. The matrix itself is never stored. Its initial vector has
component $c_\beta$, and its target functional selects row $\tau$.

For $e\ge0$, compute an entry $P(e,\alpha,\beta)=(M^e)_{\alpha\beta}$
by balanced recursion. At $e=0$ return the Kronecker delta, and at
$e=1$ use \eqref{fg:eq:matrix-entry}. For $e\ge2$, put
$a=\lfloor e/2\rfloor$ and $b=e-a$, and use
\[
 P(e,\alpha,\beta)=\sum_{|\eta|=D}
                   P(a,\alpha,\eta)P(b,\eta,\beta).
\]
Again evaluate both factors and all summands sequentially. For $e<B$
the recursion has depth $O_p(d)$, and each frame uses
$O_p(d\log(d+2))$ bits. Its total stack is
$O_p(d^2\log(d+2))$. At a leaf, the coefficient procedure adds its
own stack from \eqref{fg:eq:coefficient-space}; these space costs add.

For each $j=0,\ldots,B-1$, evaluate
\[
 a_j=\sum_{|\beta|=D}P(j,\tau,\beta)c_\beta.
\]
Output $j+1$ at the first nonzero value, and infinity if all $B$ tests
vanish. The outer counters use only $O_p(d\log(d+2))$ further bits.
Combining the bounds and $d=O(L)$ proves the asserted working space.
The finite answer has $O_p(d)$ bits. A binary-specified query can first
be compared with $B$ before this same algorithm is used.

Every loop has an explicit finite endpoint, and recursive gate indices
or power exponents decrease. On a nonhomogeneous off-promise input,
the degree-$D$ coefficients of $F^{p-1}$ define the vector used in the
same finite computation. Thus the algorithm halts without recognizing
the geometric promise. Linearity here is over $\Fp$; no semilinear
operator over an extension field has been silently substituted.
\end{proof}

The proof also gives a deterministic time bound
$2^{O_p(L^2\log(L+2))}$. The coefficient recursion has depth $O(L)$
and at most $2^{O_p(d\log(d+2))}$ children per frame, and balanced
matrix powering has $2^{O_p(d^2)}$ leaves. Repeating for $B$ indices
does not change the stated bound. The space conclusion is stronger than
this exponential-time conclusion. It is an upper bound: no matching
hardness classification for infinite height follows from it.

\begin{corollary}[Polynomial time in sparse or fixed-dimensional input]
\label{fg:cor:elementary-p}
Fix $p$. Every fixed finite exact-height predicate and every fixed
height-threshold predicate is decidable in deterministic polynomial time
for an explicitly listed sparse equation in the geometric promise above.
For a fixed number of variables, the full height is computable in
deterministic polynomial time from a division-free circuit.
\end{corollary}
\begin{proof}
For a fixed queried height, only a fixed finite list of coefficients
$b_{p^j}$ and congruences modulo $p^j$ is needed, by
\eqref{fg:eq:all-height} and \eqref{fg:eq:exact-height}. If the sparse
equation has $t$ terms, expand its $(p^j-1)$-st power by enumerating
ordered tuples of terms. There are $t^{p^j-1}$ tuples, a polynomial
number because $p,j$ are fixed. Add exponent vectors in binary, retain
the required diagonal exponent, and accumulate coefficients modulo
$p^j$. First collect equal exponent vectors modulo $p$, remove zero
coefficients, and lift the resulting homogeneous coefficients to
$\{0,\ldots,p-1\}$; this fixes an integral homogeneous lift even if
the input list originally had repeated or canceling terms.
A promised smooth input has every variable present, so the
number of variables is bounded by the length of the explicit input.
All these operations therefore have polynomial bit cost.

If $d$ is fixed, truncate the circuit at degree $d$ in its $d$ variables.
The truncated ring has the constant number $\binom{2d}{d}$ of basis
monomials. Evaluate its gates modulo $p$, then form the collected
$g=F^{p-1}$ and the constant-size operator $T$ in
\eqref{fg:eq:finite-operator}. All relevant degree bounds and the
matrix size $B$ are constant. Test its first $B$ target coefficients
by ordinary matrix-vector iteration. The total cost is polynomial in
the circuit length. Both algorithms can reject the elementary size
failures above and halt with arbitrary output outside the promise.
\end{proof}

\subsection{The cohomological and zeta-function comparison}
There is an independent finite method through the zeta function.
For $n=d-2>0$, rational Witt cohomology vanishes outside degrees zero
and $n$. To see this, use the truncation sequences
\[
 0\longrightarrow\OO_X\longrightarrow W_s\OO_X
   \longrightarrow W_{s-1}\OO_X\longrightarrow0
\]
and the coherent vanishings. Induction gives the intermediate
vanishings, while $H^0(X,W_s\OO_X)=W_s(\Fp)$.
The degree-zero and degree-$n$ transition maps are surjective, the
latter by $H^{n+1}(X,\OO_X)=0$. Their inverse systems satisfy the
Mittag--Leffler condition, so passage to Witt vectors introduces no
additional cohomology in adjacent degrees.

The rational Dieudonn\'e module of $\Phi_X$ is
$H^n(X,W\OO_X)[1/p]$
\citep[II, Corollary~4.3, p.~110]{fg-artin-mazur}.
It is zero at infinite height and has dimension $h$ at finite height;
this dichotomy is stated explicitly in
\citet[\S3, p.~6 of the preprint]{fg-yobuko}.
The comparison with the slope interval $[0,1)$ in crystalline cohomology
is \citet[II, Corollary~3.5, (3.5.3), p.~616]{fg-illusie}; see also
\citet[Theorem~1.1 and (6.2)]{fg-bbe}. Consequently
\begin{equation}\label{fg:eq:zeta-infinity}
 h(X)=\infty\iff\bigl((1-T)Z(X,T)\bigr)^{<1}=1.
\end{equation}
The factor $1-T$ removes degree-zero cohomology. All remaining
slope-less-than-one factors come from the top Witt group, with one
parity, so they cannot cancel each other.

For an exact algorithm, expand the input circuit modulo the ideal of
all monomials of total degree greater than $d$. At most
$\binom{2d}{d}\le4^d$ monomials occur; truncation is a ring
homomorphism and is valid at every gate regardless of its formal
degree. Partition projective space by the first nonzero coordinate,
set it equal to one, and specialize the earlier coordinates to zero.
This gives $d$ disjoint affine hypersurfaces with at most $d-1$
variables and degree at most $d$. Zeta functions multiply across the
partition. Zero and nonzero constant equations give affine spaces and
empty sets and are handled directly. For each other chart,
\citet[Theorem~37, p.~29 of the preprint]{ff-lauder-wan} give the
deterministic bit bound
\[
 \widetilde O\left(2^{13r^2}a^{3r+7}
                   e^{3r^2+9r}p^{2r+4}\right)
\]
for the full zeta function of a degree-$e$ polynomial in $r$ variables
over $\F_{p^a}$, without a smoothness hypothesis. With $a=1$, $r<d$,
$e\le d$, this is $2^{O_p(d^2\log(d+2))}$.

Reduce $(1-T)Z(X,T)=A(T)/B(T)$ to coprime integer polynomials with
constant coefficient one. Products, exact gcds and divisions have
polynomial bit cost in the explicit zeta output. Coprimeness over
$\Q$ persists over $\Q_p$. For
$C(T)=1+\sum_i c_iT^i$, all reciprocal roots have valuation at least
one exactly when $p^i\mid c_i$ for every $i$. One direction follows
from elementary symmetric functions; for the other, $C(U/p)$ has
integral coefficients, so its reciprocal polynomial is monic integral
and its roots are integral. Apply this test to both $A$ and $B$ in
\eqref{fg:eq:zeta-infinity}. This gives the same deterministic
exponential-time bound, with an independent cohomological explanation
of what the infinite-height predicate detects.

The finite-height estimate
$h\le\dim H^{n-1}(X,\Omega_X^1)+1$ of
\citet[Corollary~2.3, p.~100]{fg-vdgk} gives another geometric stopping
bound. The operator argument above supplies a bound directly from the
degree and variable count, while controlling the storage needed for a
succinct input.

\section{Integral invariant derivatives at the counted height}\label{fg:sec:derivatives}
The height tests have a differential form on the integral Stienstra
lift. We prove the coefficientwise estimates for arbitrary lifts over
an unramified Witt ring; a specified Honda model is not a hypothesis.
These estimates occur in the earlier manuscript
\citet[\S\S2--3]{fg-invariant-derivatives}. The proofs here also give the
precise next-height test needed for the counting family.

Let $R=W(k)$ for a perfect field $k$ of characteristic $p$, and let
$\mathcal F$ be a one-dimensional commutative formal-group law over $R$.
Its normalized invariant derivation and logarithm are
\[
 D=v(T)\frac d{dT},\qquad v(T)=\partial_2\mathcal F(T,0)\in1+TR[[T]],
 \qquad L(T)=\int_0^T v(s)^{-1}\,ds.
\]
Associativity shows that $L(\mathcal F(T,U))=L(T)+L(U)$.
For the Stienstra coordinate, equation~\eqref{fg:eq:stienstra} gives
\[
 D=\left(\sum_{m\ge1}b_mT^{m-1}\right)^{-1}\frac d{dT}.
\]
For an integer $h\ge1$ put
\[
 q=p^h,\qquad c_h=\frac1{p-1}-\frac1{p^h-1},\qquad
 E_h(a)=\sum_{\substack{j\ge1\\h\nmid j}}
                    \left\lfloor\frac a{p^j}\right\rfloor.
\]

\begin{theorem}[Divisibility of invariant derivatives]\label{fg:thm:derivatives}
If the reduction of $\mathcal F$ has height at least $h$, then for every
$a\ge0$ and $f\in R[[T]]$,
\begin{equation}\label{fg:eq:derivative-bound}
 D^af\in
 p^{\max\{0,E_h(a),\lceil c_ha-(h-1)\rceil\}}R[[T]].
\end{equation}
If its height is exactly $h$, then for $j=hr+s$, $0\le s<h$,
\[
 v_p\bigl(D^{p^j}(T^{p^s})(0)\bigr)=E_h(p^j),\qquad
 v_p\bigl(D^{q^r}T(0)\bigr)=c_h(q^r-1).
\]
Thus $c_h$ is the optimal uniform asymptotic slope at exact height $h$.
The error $h-1$ in the affine bound is optimal for every fixed normalized
group satisfying the height lower bound.

At infinite height, $L$ and $L^{-1}$ are integral and
\[
 D^aR[[T]]\subseteq p^{v_p(a!)}R[[T]].
\]
This exponent is attained. The slope $1/(p-1)$ has no bounded uniform
affine error in this infinite-height case.
\end{theorem}

\begin{proof}
We first prove the affine estimate. By Lemma~\ref{fg:lem:height-log},
the coefficients of $L$ below degree $q$ are integral. The new coordinate
\[
 z=T+\sum_{2\le m<q}[T^m]L(T)\,T^m
\]
has an integral compositional inverse. In this coordinate the logarithm
is $z+bz^q+O(z^{q+1})$, with $pb\in R$. A strict integral coordinate
change conjugates $D$ and preserves every ideal $p^eR[[T]]$, so it
suffices to work in this coordinate.

Define integral translation operators $\Delta_i$ by
\[
 f(\mathcal F(z,u))=\sum_{i\ge0}(\Delta_i f)(z)u^i.
\]
Logarithmic addition gives
$\sum_i\Delta_i u^i=\exp(L(u)D)$. Coefficient comparison yields
\[
 \Delta_i=D^i/i!\quad(i<q),\qquad
 \Delta_q=D^q/q!+bD.
\]
Set $B=v_p(q!)-1=c_h(q-1)$ and write $q!=p^{B+1}w$, with $w$ a unit.
Then
\begin{equation}\label{fg:eq:operator-recursion}
 D^q=p^BuD+p^{B+1}A,\qquad u=-wpb\in R,
 \qquad A=w\Delta_q,
\end{equation}
where $A$ preserves the integral function ring.

For $0\le a<q$, the equality $D^a=a!\Delta_a$ gives factorial
divisibility. Writing $a=\sum_{i=0}^{h-1}a_i p^i$, Legendre's formula
gives
\[
 v_p(a!)-c_ha
 =\sum_{i=0}^{h-1}a_i
       \left(\frac{p^i}{q-1}-\frac1{p-1}\right)\ge1-h.
\]
Every parenthesized coefficient is nonpositive, so the minimum occurs
when all digits are $p-1$, at $a=q-1$.
For $a\ge q$, apply \eqref{fg:eq:operator-recursion} to $D^{a-q}f$.
Induction bounds the first term by $c_ha-(h-1)$ and the second by
that quantity plus $1-c_h\ge0$. This uses only that $A$ preserves
integral ideals; it needs no commutation of $A$ with $D$.
Integer-valued valuations and integrality prove the affine estimate.
For every normalized coordinate,
\[
 D^{q-1}(T^{q-1})(0)=(q-1)!,\qquad
 v_p((q-1)!)=c_h(q-1)-(h-1).
\]
This proves optimality of the error independently of the slope argument.

For the floor-sum estimate, adjoin $\pi$ with $\pi^{q-1}=p$ and work
in the discrete valuation ring $R'=R[\pi]$. The rescaled law
\[
 \mathcal F_\pi(X,Y)=\pi^{-1}\mathcal F(\pi X,\pi Y)
\]
is integral, and every coefficient of its $p$-series is divisible by
$p$. Below degree $q$ this follows from the height hypothesis; in
degree $m\ge q$ it follows from the factor $\pi^{m-1}$.
The recursion \eqref{fg:eq:log-recursion}, in every degree, therefore
makes its logarithm integral over $R'$, and its inverse is integral.
If $E=L^{-1}=\sum e_mT^m$, uniqueness gives
\[
 E_\pi(T)=\pi^{-1}E(\pi T),\qquad
 v_p(e_m)\ge-\frac{m-1}{q-1}.
\]
This assertion uses divisibility by $p$ of the entire rescaled
$p$-series, rather than only reduction modulo the ramified uniformizer.

Logarithmic differentiation gives
\[
 D^af(0)=a![T^a]f(E(T)).
\]
A product of $b\ge1$ exponential coefficients with total degree $a$
has valuation at least $-(a-b)/(q-1)$. Since the resulting coefficient
lies in the unramified field $R[1/p]$, its valuation is an integer.
Consequently
\[
 v_p(D^af(0))\ge
 v_p(a!)-\left\lfloor\frac{a-1}{q-1}\right\rfloor
 \quad(a\ge1).
\]
The same estimate holds coefficientwise: apply the evaluation argument
in $u$ to each coefficient in $T$ of $f(\mathcal F(T,u))$, using
$D_u^a f(\mathcal F(T,u))|_{u=0}=D_T^af(T)$.

Now write $a=\sum_j a_jp^j$. Compose these coefficientwise estimates
for the powers $(D^{p^j})^{a_j}$. The finite identities
\[
 \sum_j a_jv_p((p^j)!)=v_p(a!),\qquad
 \sum_j a_j\left\lfloor\frac{p^j-1}{q-1}\right\rfloor
       =\sum_{r\ge1}\left\lfloor\frac a{q^r}\right\rfloor
\]
give $v_p(D^af)\ge E_h(a)$. Together with the affine estimate this
proves \eqref{fg:eq:derivative-bound}.

Suppose the height is exactly $h$, and let $\gamma_q$ be the coefficient
of $T^q$ in $[p]_{\mathcal F}$. It is a unit. Modulo $\pi$ the rescaled
$p$-series divided by $p$ is $T+\overline\gamma_qT^q$.
Indeed, degrees below $q$ other than one have an additional positive
power of $\pi$, and degrees above $q$ do also. Divide the logarithm
recursion by $p$. All its terms involving the $i$-th power of the
$p$-series for $i\ge2$ vanish modulo $\pi$. Hence
\[
 \overline L_\pi(T)=T+\overline\gamma_qT^q,\qquad
 \overline E_\pi(T)=\sum_{r\ge0}\beta_rT^{q^r},
\]
where $\beta_0=1$ and $\beta_r=-\overline\gamma_q\beta_{r-1}^q$.
Every $\beta_r$ is nonzero. For $j=hr+s$, the coefficient of $T^{p^j}$
in $\overline E_\pi(T)^{p^s}$ is therefore nonzero. Undoing the scaling
gives
\[
 v_p\bigl([T^{p^j}]E(T)^{p^s}\bigr)
 =-\frac{p^j-p^s}{q-1}
 =-\left\lfloor\frac{p^j-1}{q-1}\right\rfloor.
\]
Multiplication by $(p^j)!$ proves both witness equalities. Letting
$r$ tend to infinity in the witness with $f=T$ excludes every larger
slope with bounded, or even sublinear, error for this fixed group.

At infinite height every coefficient of $[p]$ belongs to $pR$.
The logarithm recursion proves $L\in R[[T]]$ in every degree, and
its inverse is integral. In its coordinate, $D=d/dL$, whose $a$-th
iterate is coefficientwise divisible by $a!$. The function $L(T)^a$
attains $a!$. Finally
\[
 v_p(a!)=\frac a{p-1}-\frac{s_p(a)}{p-1}.
\]
At $a=p^r-1$ the digit error is $r$, which is unbounded. This proves
the infinite-height assertions.
\end{proof}

\begin{corollary}[A differential test for the next height]\label{fg:cor:derivative-height}
For a Stienstra lift whose reduction has height at least $r$, put $q=p^r$.
Then
\[
 \frac{D^qT(0)}{q!/p}\equiv
       -\frac{b_q}{p^{r-1}}\pmod p,
\]
and both quotients are integral. Moreover
\[
 h(X)\ge r+1
 \iff D^qT(0)\in q!R
 \iff D^qR[[T]]\subseteq q!R[[T]].
\]
For $r=2$, the valuation of $D^{p^2}T(0)$ is exactly $p$ at height
two and at least $p+1$ at height at least three.
\end{corollary}

\begin{proof}
Write $L=T+\sum\ell_mT^m$ and $E=T+\sum e_mT^m$.
Lemma~\ref{fg:lem:height-log} and formal inversion give
$e_m\in R$ for $m<q$, and $e_q+\ell_q\in R$. Since
$D^qT(0)=q!e_q$, this gives the signed congruence and the pointwise
equivalence. At height at least $r+1$, the floor sum in
Theorem~\ref{fg:thm:derivatives}, with parameter $r+1$ and order $q$,
is $v_p(q!)$. This proves the whole-operator inclusion. The converse
follows by testing $T$ at zero. Finally $v_p((p^2)!)=p+1$.
\end{proof}

For the counting family, $S=0\bmod p$ gives height at least two.
Every integral formal function on its Stienstra lift consequently
satisfies
\[
 D^af\in p^{\max\left\{0,
    \sum_{j\ge0}\left\lfloor a/p^{2j+1}\right\rfloor,
    \left\lceil pa/(p^2-1)-1\right\rceil\right\}}R[[T]].
\]
In particular $\overline D^{\,p}=0$. This reduced operator vanishes
at every height at least two, including infinity; the next distinction
is the additional factorial divisibility in
Corollary~\ref{fg:cor:derivative-height}.

\subsection{The base and the Honda-dual hypothesis}
The unramified base and the scope of the height hypothesis are essential.
Over $\Z_p[[u]]$ the law $X+Y+uXY$ has
\[
 D=(1+uT)\frac d{dT},\qquad D^aT=u^{a-1}(1+uT)\quad(a\ge1).
\]
At $u=0$ its reduction is additive, while $D^pT$ is not divisible by
$p$ over the entire parameter ring. If $u=\pi$ in a ramified discrete
valuation ring with $v_\pi(p)=e>p-1$, the residue group is again additive,
but $v_p(D^pT(0))=(p-1)/e<1$. These computations explain exactly why
height at one closed point, or reduction only modulo a ramified
uniformizer, does not give the theorem's coefficientwise $p$-ideals.

The directional result in \citet[\S6]{fg-invariant-derivatives} concerns
the Cartier dual of the specified Honda lift with logarithm
$\sum_{r\ge0}T^{p^{hr}}/p^r$. That model has $h-1$ specified primitive
directions, with affine slopes $p^i/(p^h-1)$ and error one for
$1\le i<h$. An integral identification of that model and its directions
is an additional hypothesis. Neither the counting construction nor
Stienstra's formula provides it, so that directional result is not used
in the preceding proofs.

Equal height need not even give an isomorphism over the prime field.
For the Fermat cubic over $\Z_5$,
\[
 b_5=0,\qquad b_{25}=24!/(8!)^3=9465511770,\qquad
 b_{25}/5=4\pmod5.
\]
Its height-two leading $p$-series coefficient is four. The standard
Honda law just specified has coefficient one. An isomorphism over
$\F_5$ would multiply that coefficient by $u^{1-25}=1$ for
$u\in\F_5^\times$, which cannot identify them. In particular their
integral lifts over $\Z_5$ are not isomorphic. This is an obstruction
to inferring the model identification, rather than an assertion that
every numerical directional estimate fails in every other model.

\InputIfFileExists{sections/legacy.tex}{}{}
\section{Uniform primes and a yes-or-no oracle}\label{sec:uniform}
This section makes the assignment-count shifts parsimonious and removes a circularity in choosing primes for the sparse-list construction.

\begin{lemma}[Unique gate extensions]\label{lem:tseitin}
Any fan-in-two Boolean circuit can be converted in linear time to a width-three CNF whose satisfying assignments correspond bijectively to primary inputs on which the circuit outputs one.
\end{lemma}
\begin{proof}
Give each gate a new output bit $t$. For an AND gate $t=u\wedge v$ use
\[
(\neg t\vee u),\quad(\neg t\vee v),\quad(t\vee\neg u\vee\neg v).
\]
For an OR gate $t=u\vee v$ use
\[
(t\vee\neg u),\quad(t\vee\neg v),\quad(\neg t\vee u\vee v).
\]
For $t=\neg u$ use $(t\vee u)$ and $(\neg t\vee\neg u)$. A constant is enforced by a unit clause. Inspecting the two or four input patterns proves that each collection forces exactly its indicated output. Induction in gate order gives one legal wire word for each primary input. A positive unit clause on the final output selects the accepted inputs. All clauses have width at most three.
\end{proof}

\begin{lemma}[An exact additive shift]\label{lem:shift}
Given a formula $\Phi$ on $n$ primary bits with $S$ solutions and an integer $0\le a<p$, a width-three formula with exactly $S+a$ solutions can be constructed with $O(|\Phi|+\log p)$ variables and clauses.
\end{lemma}
\begin{proof}
Let $k=\lceil\log_2p\rceil$, let $z$ be a $k$-bit word, and add a switch $b$. Build the Boolean circuit
\[
(\neg b\wedge\Phi(x)\wedge[z=0])\ \vee\
(b\wedge[x=0]\wedge[z<a]).
\]
The first branch contributes precisely $S$ primary words and the second precisely $a$. Every bit unused in a branch is fixed there, so there is no spurious power-of-two factor. To compute $z<a$, read the bits from most significant to least. Maintain bits $E$ and $L$ meaning that the prefixes are equal or that the prefix of $z$ is smaller; initialize $(E,L)=(1,0)$. At a constant bit $a_j=0$, replace $E$ by $E\wedge\neg z_j$ and leave $L$ unchanged. At $a_j=1$, replace
\[
L\ \text{by}\ L\vee(E\wedge\neg z_j),\qquad
E\ \text{by}\ E\wedge z_j,
\]
using the old $E$ in the first update. At most four elementary gates per bit suffice, including negation. This also handles $a=0$, when the final comparison is always false. Use binary trees for the conjunctions and apply Lemma~\ref{lem:tseitin}. All gate values remain uniquely determined even on an unselected branch, so the final CNF has exactly $S+a$ satisfying assignments.
\end{proof}

\begin{proposition}[Uniform reconstruction]\label{prop:uniform}
A scalar-evaluation oracle for the smooth circuit hypersurfaces, uniform over polynomially bounded primes, computes $\SAT$ with $n+1$ queries and polynomial additional work. A uniform yes-or-no oracle for nonvanishing, or equivalently for Frobenius splitting on these inputs, decides satisfiability with $n+1$ queries and computes $\SAT$ with
\[
O(n^2\log(n+2))
\]
queries. The reductions use only primes of polynomial magnitude, and every query is a proved smooth instance.
\end{proposition}
\begin{proof}
Take the first $n+1$ primes. Their product is strictly greater than $2^n$, and their largest is $O((n+1)\log(n+2))$, by \citet[Corollary to Theorem~3, (3.13), p.~69]{root-rosser-schoenfeld}. Deterministic sieving finds them in polynomial time. Theorem~\ref{thm:main} produces circuits in time polynomial in $|\Phi|+p$. Their residues determine the unique integer $S\in[0,2^n]$ by the Chinese remainder theorem.

For decision alone, if $S=0$ every prime gives zero; if $0<S\le2^n$, not all of these primes can divide $S$. Thus the disjunction of the nonvanishing answers decides satisfiability. At one fixed prime a satisfiable formula may have zero residue, so one prime alone does not suffice for this argument.

For exact reconstruction from a yes-or-no oracle, for each selected prime query the shifted formulas of Lemma~\ref{lem:shift} for all $0\le a<p$. Exactly one shift has zero scalar. If it is $a_0$, then $S\equiv-a_0\pmod p$. The number of calls is at most the sum of the first $n+1$ primes, bounded by $(n+1)$ times their largest, giving the displayed estimate. All shifted formulas have polynomial length. Chinese remainder arithmetic has polynomial bit cost since its total modulus has polynomial binary length.
\end{proof}

The uniformity here is in the numerical prime $p$, not merely in $\log p$. A uniform polynomial-time evaluator or decision algorithm in this model would compute $\#\mathrm P$ functions in deterministic polynomial time. The fixed-prime theorem makes the more limited modular-counting assertion.

\begin{proposition}[Separate tags give sparse lists]\label{prop:sparse-tags}
Let the local factors $P_f,Q_i$ in \eqref{eq:incidence} be $A_1,\ldots,A_s$, of degrees $k_i$ summing to $b$. If $p>s$, there is a polynomial-length explicit sparse polynomial whose smooth stabilization has coherent scalar $S\bmod p$. Both scalar and yes-or-no uniform reconstruction can be performed using only such sparse-list inputs at polynomially bounded primes.
\end{proposition}
\begin{proof}
Pad the factor list to $p-1$ entries by constants one, of degree zero. Partition the $b$ incidence coordinate indices into sets $I_i$ of cardinalities $k_i$. For $Z=\prod_{j=1}^b w_j$, put $M_i=\prod_{j\notin I_i}w_j$. Give each factor a different tag and set
\[
G_{\mathrm{list}}=-t_1^{p-1}M_1A_1+
             \sum_{i=2}^{p-1}t_i^{p-1}M_iA_i.
\]
Every summand has degree $b+p-1$, equal to the number of variables. Distinct tags force one choice of each summand in the $(p-1)$st power. Since each coordinate is omitted by exactly one $M_i$,
\[
\prod_iM_i=Z^{p-2},\qquad
h_p(G_{\mathrm{list}})=-(p-1)![Z]\prod_iA_i=S.
\]
There are at most $8s+(p-1-s)$ displayed terms before combining equal exponents. Padding and stabilization add only polynomially many terms and binary exponent bits. The tag attached to a factor is essential: a common tag would not force the local factors to be chosen separately.

For scalar queries, take the first $n+1$ primes greater than $s$. Their largest is at most the $(s+n+1)$st prime, so it is polynomially bounded and their product is greater than $2^n$.

For yes-or-no queries the primes must exceed the factor counts of \emph{all} shifted formulas. Here is a bound chosen in advance. Let $L\ge2$ dominate the original variable and clause counts and circuit description length. The construction in Lemma~\ref{lem:shift} can be made with at most
\[
100(L+\lceil\log_2p\rceil+3)
\]
clause and variable factors combined; this follows by using at most six gates per original clause, two gates per bit for each zero test, four per comparison bit, and at most three clauses per gate. Set $B=10^6(L+2)^2$ and take the first $n+1$ primes greater than $B$. They occur among the first $B+n+1$ primes. The same cited bound implies $p_j<j^2$ for $j\ge6$, giving
\[
p<(B+n+1)^2=:P.
\]
Thus $\lceil\log_2p\rceil\le\lceil\log_2P\rceil$, and the displayed factor bound is strictly smaller than $B$ for every $L\ge2$. For example, $\log_2(B+L+1)\le21+2\log_2(L+2)$, and substitution makes the inequality immediate. Consequently \emph{every} shift at \emph{every} selected prime satisfies $p>s$. Sieving through $P$, constructing the queries, and making at most $(n+1)P$ calls all take polynomial time in $L$. The bound is deliberately coarse; its purpose is to make prime selection independent of the shifts. The same Chinese remainder argument completes reconstruction.
\end{proof}

\InputIfFileExists{sections/presentations.tex}{}{}
\InputIfFileExists{sections/saturation_ci.tex}{}{}
\subsection{Recognizing relations over the prime field}
Geometric integrality permits degree reductions beyond those implied
by radicality. An algebraic element of the function field must already
belong to the prime field. A linearized polynomial relation can
therefore force a hidden factor to vanish, provided its coefficient
sum is nonzero. The following theorem recognizes such a relation with
an unknown number of lower Frobenius blocks and unknown coefficients.
Additive-polynomial composition and its skew-polynomial interpretation
are classical. In particular, \citet[\S6, Theorem~6.1, pp.~22--24 of
the author version]{sp-giesbrecht} gives deterministic decomposition
in time polynomial in the Frobenius degree at fixed characteristic,
including characteristic two. That input already specifies a
univariate additive polynomial. Here the sparse homogeneous equation
hides both its additive coefficients and its multivariate argument
$U/V$; the exponent gap below permits their recovery.

\begin{theorem}[Sparse recognition of a linearized relation]
\label{sp:thm:linearized-preprocessing}
Fix a prime $p$. Let a homogeneous sparse list with binary exponents
generate $K\subset S=\Fp[x_0,\ldots,x_N]$, and suppose that
$I=K^{\mathrm{sat}}$ has geometrically integral quotient $S/I$.
Consider a nonzero row admitting the exact identity
\begin{equation}
 F=U^Q+\sum_{i=0}^{s}a_iV^{Q-p^i}U^{p^i},
 \qquad Q=p^e,\quad 0\le s<e,
 \label{sp:eq:linearized-relation}
\end{equation}
where $U\ne0$ is homogeneous of degree $d\ge1$, $V$ is any
coefficient-one monomial of degree $d$, and $a_i\in\Fp$. Assume
\begin{equation}
 (d+1)p^s\le Q,\qquad
 \lambda:=1+\sum_{i=0}^{s}a_i\ne0.
 \label{sp:eq:linearized-gap}
\end{equation}
None of $U,V,e,s,d$ or the coefficient vector is given. The number
$s$ is not fixed, and zero lower coefficients are allowed.

There is a deterministic polynomial sparse-bit algorithm that finds
and verifies such an identity whenever one exists and replaces $F$
by $U$. Each accepted replacement satisfies
\[
 K\subseteq K'\subseteq I,\qquad (K')^{\mathrm{sat}}=I,
\]
introduces no coordinate index, does not increase the number of
terms, and divides the numerical degree by $Q\ge p$.
The algorithm may choose a minimum-degree replacement among its
passing candidates. Repeated use, interleaved with rank-one digit
reduction and batched valid coordinate-content division, terminates
in polynomial sparse-bit time. An ordinary whole Frobenius power
may also be handled by its existing root-extraction rule.
\end{theorem}
\begin{proof}
Put $R=S/I$ and $E=\operatorname{Frac}(R)$. If $z\in E$ is
algebraic over $\Fp$, then $\Fp(z)=\mathbb F_{p^b}$ for some $b$.
Tensoring its inclusion in $E$ with $\overline{\Fp}$ preserves
injectivity. For $b>1$ the source is a product of fields and has
nonzero zero divisors, whereas the target is a localization of the
domain $R\otimes_{\Fp}\overline{\Fp}$. Hence $b=1$.
This is the usual geometric-integrality condition
\citep[Definition~10.49.1 and Lemma~10.49.3]{sp-stacks-integral}.
The saturated homogeneous coordinate ring of a geometrically
integral projective scheme has this property: after base change,
its punctured affine cone is a $\mathbb G_m$-torsor over the
integral scheme, and its affine closure is integral.

Reduce \eqref{sp:eq:linearized-relation} modulo $I$.
If $V=0$, then $U^Q=0$ and $U=0$ in the domain $R$.
Otherwise $z=U/V\in E$ satisfies
\[
 z^Q+\sum_{i=0}^{s}a_i z^{p^i}=0.
\]
It is algebraic over $\Fp$ and therefore lies in $\Fp$.
The equation becomes $\lambda z=0$, forcing $z=0$.
Thus $U\in I$. Every term of $F$ is divisible by $U$,
which proves the ideal sandwich and equality of saturations.
This safety argument does not require the gap or monomial hypotheses;
those hypotheses will make the unknown relation recognizable.

Write $D=\deg F$ and $T=|\operatorname{supp}F|$.
Enumerate $Q=p^e\le D$ with $Q\mid D$, set $d=D/Q$, and
enumerate $0\le s<e$ satisfying \eqref{sp:eq:linearized-gap}'s
degree inequality. There are $O(\log^2(D+1))$ pairs $(e,s)$.
For each pair, take the \emph{$Q$-section} $W$ of $F$: select
the terms whose exponent vectors are coordinatewise divisible by
$Q$, and divide those vectors by $Q$, retaining their coefficients.
Reject a zero or wrong-degree section.

To determine what a true relation gives, write
$U=\sum_u c_u x^u$ and $V=x^v$, with $|u|=|v|=d$.
A term of its $i$th lower block has exponent
\[
 (Q-p^i)v+p^iu=Qv+p^i(u-v).
\]
Every coordinate of $p^i(u-v)$ has absolute value at most
$dp^s\le Q-p^s<Q$. It is coordinatewise divisible by $Q$
only when $u=v$. If $c_v=0$ when $V$ is absent from $U$, then
\begin{equation}
 W=U+(\lambda-1)c_vV.
 \label{sp:eq:leading-section}
\end{equation}
Because $\lambda\ne0$, $W$ has exactly the support of $U$.
No membership in $I$ is inferred from taking a section alone.

First suppose $V$ occurs in this support. Enumerate its at most
$T$ possibilities and the $p-1$ possible nonzero values of
$\lambda$. Leave all other coefficients of $W$ unchanged and
divide its $V$ coefficient by $\lambda$ to propose $U$.
For these fixed candidates form the sparse polynomials
\[
 B_i=V^{Q-p^i}U^{p^i}\quad(0\le i\le s),\qquad H=F-U^Q.
\]
Solve over $\Fp$ the coefficient equations
\begin{equation}
 H=\sum_{i=0}^{s}a_i B_i,
 \qquad \sum_{i=0}^{s}a_i=\lambda-1.
 \label{sp:eq:linearized-system}
\end{equation}
Use the union of the supports of $H$ and the $B_i$ as the row
indices of this linear system. Its unknowns are only the $s+1$
scalars $a_i$. A consistent system gives one coefficient vector;
there is no enumeration of $p^{s+1}$ vectors. Verify the complete
identity and every degree condition before accepting. The true
support-present representation is among the candidates.

For the other case set $U=W$ and again put $H=F-U^Q$.
If $H=0$, ordinary whole-power extraction applies. It also admits
a relation with all lower coefficients zero, $s=0$, and any
support monomial of $U$ as $V$, when $d+1\le Q$.
Otherwise select a surviving exponent vector $\alpha$ of $H$
and put $q=p^s$, $M=Q-q$. In a true representation it has the form
\[
 \alpha=Mv+\beta,\qquad
 \beta=(q-p^i)v+p^iu\ge0,\qquad |\beta|=qd\le M.
\]
Cancellation may delete terms but cannot produce a new exponent,
so this description holds for any chosen surviving term.
Consequently $\lfloor\alpha/M\rfloor$ equals either $v$ or
$v+e_j$ for one coordinate unit vector $e_j$. The latter requires
$qd=M$ and $\beta=Me_j$. Enumerate that quotient and the vectors
obtained by subtracting one coordinate unit, retaining only
nonnegative vectors of total degree $d$. There are at most $v_0+1$
candidates, where $v_0$ is the number of active coordinates in
the input row. For each candidate $V$ and each
$\lambda\in\Fp^*$, solve \eqref{sp:eq:linearized-system} and
verify the full identity. This includes the boundary case of
equality in the degree gap and the true absent-support representation.

The searches have polynomial sparse-bit cost. Each $B_i$ is obtained
by multiplying exponent integers and adding a monomial shift;
its support has at most $T$ terms. The coefficient matrix has at
most $(s+2)T+T+1$ rows and $s+1\le\log_p D$ columns.
Gaussian elimination, sparse collection, and exact verification
are polynomial in the written input length. Intermediate exponent
integers have polynomial bit length, since a candidate identity is
homogeneous of degree $D$. Each positive coordinate exponent of
$V$ in the nonzero-remainder case occurs in that remainder, because
$M>0$. Thus absent coordinate identifiers need never be enumerated.

Every accepted $U$ is a section with at most one coefficient changed,
so it has no more than $T$ terms and introduces no coordinate.
All passing candidate outputs can be compared by degree, with a
fixed ordering to break ties. Each new step decreases degree by
at least a factor two. The known digit operation decreases it by
at least the fraction $1/(v_0+2)$ and also does not increase
support or active coordinates. Batched eligible monomial content
is removed at most once between such steps. Hence there are
$O(v_0\log(D_0+1))$ successful degree-reducing steps from initial
degree $D_0$, proving polynomial termination and the combined cost.
\end{proof}

\subsection{An affine relation with no prime-field root}
A nonzero constant term can force both the inner polynomial and its
monomial denominator to vanish. The resulting replacement has two
rows, so its complexity must be measured by their total support.

\begin{theorem}[Sparse recognition of an affine relation]
\label{sp:thm:affine-preprocessing}
Fix $p$ and let a homogeneous sparse binary-exponent list generate
$K\subset S=\Fp[x_0,\ldots,x_N]$, with $I=K^{\mathrm{sat}}$ and
$S/I$ geometrically integral. Suppose a nonzero row has a representation
\begin{equation}
 F=U^Q+\sum_{i=0}^s a_iV^{Q-p^i}U^{p^i}+\beta V^Q,
 \qquad Q=p^e,\quad 0\le s<e,
 \label{sp:eq:affine-relation}
\end{equation}
where $U$ is homogeneous, $V$ is a coefficient-one monomial of the
same degree $d\ge1$, and
\[
 (d+1)p^s\le Q,\qquad a_i\in\Fp,\qquad
 1+\sum_i a_i=0,\qquad \beta\in\Fp^\times.
\]
None of these data is given. A deterministic polynomial sparse-bit
algorithm finds and verifies a representation whenever one exists,
normalizing $U$ so that $V$ is absent from its support, and replaces
$F$ by the nonzero rows among $U,V$. Zero normalized $U$ is allowed.
The generated ideals satisfy
\[
 K\subseteq K'\subseteq I,\qquad (K')^{\mathrm{sat}}=I.
\]
The total output support is at most the support of $F$, no coordinate
index is introduced, and every output row has degree $\deg(F)/Q$.
The algorithm can minimize this degree among its passing candidates.
Its repeated use with the preceding support-nonincreasing operations
has polynomial bit cost, including when the number of rows grows.
\end{theorem}
\begin{proof}
The change $U\mapsto U-cV$, $c\in\Fp$, changes the right side of
\eqref{sp:eq:affine-relation} by
$-c(1+\sum_i a_i)V^Q=0$. Subtracting the coefficient of $V$ therefore
gives the stated normalization without changing any other parameter.

As in Theorem~\ref{sp:thm:linearized-preprocessing}, a lower-block
exponent $Qv+p^i(u-v)$ is coordinatewise divisible by $Q$ only when
$u=v$: the coordinates of $u-v$ have absolute value at most $d$, while
$Q/p^i\ge d+1$. Thus the divided pure $Q$-section is exactly
\[
 W=E_Q(F)=U+\beta V.
\]
Its support is the disjoint union of that of $U$ and the monomial $V$.
For each power $Q$ dividing $D=\deg F$, and each $s$ satisfying the
gap, enumerate $V$ from the support of $W$. Its coefficient is the
only possible $\beta$, and deleting that term gives $U$. Solve the
sparse coefficient system
\[
 \sum_{i=0}^s a_i V^{Q-p^i}U^{p^i}
       =F-U^Q-\beta V^Q,
 \qquad \sum_{i=0}^s a_i=-1.
\]
There are $O(\log D)$ unknowns and $O(|\operatorname{supp}F|\log D)$
coefficient equations. Deterministic elimination and an exact sparse
identity check accept precisely the passing representations. The
true normalized representation is among them. If $U=0$, all block
columns are zero; the system accepts exactly when $F=\beta V^Q$,
and the sum constraint still has a solution. Ordering $Q$
decreasingly and fixing the monomial and pivot orders gives the
minimum output degree with deterministic ties.

For safety, suppose that the image of $V$ in $S/I$ were nonzero.
Then $t=U/V$ in its fraction field would be algebraic over $\Fp$.
Geometric integrality and the constant-field argument in
Theorem~\ref{sp:thm:linearized-preprocessing} give $t\in\Fp$.
The relation becomes
\[
 0=t^Q+\sum_i a_it^{p^i}+\beta
   =\left(1+\sum_i a_i\right)t+\beta=\beta,
\]
a contradiction. Hence $V\in I$, and the original relation gives
$U^Q\in I$, so $U\in I$. Every term of $F$ belongs to $(U,V)$.
The ideal sandwich and equality of saturations follow.

The section identity gives
\[
 |\operatorname{supp}U|+|\operatorname{supp}V|
   =|\operatorname{supp}W|\le|\operatorname{supp}F|.
\]
All output exponents are original section exponents divided by $Q$,
so no coordinate index is added. The degree decreases by a factor
$Q\ge p$. Every candidate block is obtained by a Frobenius power and
a monomial shift; neither large numerical powers nor coefficient
vectors are enumerated. This proves the cost of one recognition pass.

For iteration, let $T_0$ be the original total support, $D_0$ the
maximum original degree, and $v$ the number of occurring coordinates.
Every intermediate nonzero list has at most $T_0$ rows. Use
\[
 \Phi=\sum_{H\text{ a nonzero row}}
      |\operatorname{supp}H|\bigl(1+\log\deg H\bigr).
\]
An affine step decreases $\Phi$ by at least $\log p$, because its
total output support does not increase and both outputs have the
same smaller degree. A nontrivial rank-one step decreases it by at
least $1/(v+2)$, using the degree estimate of
Lemma~\ref{sci:lem:deflation}; the other preceding degree reductions
have the same or a stronger bound. Thus the number of these steps
is polynomial in $T_0,v,\log(D_0+1)$. Full eligible coordinate content
is removed in maximal batches between them. Each batch has at most
$v$ removals per row, since changing one coordinate does not alter
the other coordinates' minimum exponents. This also has polynomial
cost. Current supports, coefficient sizes, coordinate labels and
exponent lengths stay bounded by the original input parameters, so
all intervening recognition scans are polynomial.
\end{proof}

The conclusion forces a coordinate variable into $I$, since the
positive-degree monomial $V$ belongs to the prime ideal. It therefore
does not give an additional reduction on a projective scheme lying
in no coordinate hyperplane. A homogeneous relation with
$1+\sum_i a_i=0$ and $\beta=0$ needs a different output: the relation
allows every prime-field value of $U/V$, and does not force $U$ or
$V$ separately into $I$.

\subsection{Retaining all prime-field components}
A homogeneous linearized relation with zero coefficient sum does not
specify one prime-field value of its argument. It nevertheless permits
a smaller equation that keeps every such value. Additive composition
is classical \citep[\S6]{sp-giesbrecht}; the issue here is recognizing
the hidden argument and controlling support through repeated reductions.

\begin{theorem}[A zero-sum relation with controlled support growth]
\label{sp:thm:zero-sum-preprocessing}
Fix $p$ and let a homogeneous sparse binary-exponent list generate
$K\subset\Fp[x_0,\ldots,x_N]$, with $I=K^{\mathrm{sat}}$ and $S/I$
geometrically integral. Suppose a nonzero row admits
\[
 F=U^Q+\sum_{i=0}^{s}a_i V^{Q-p^i}U^{p^i},
 \qquad Q=p^e>p,\quad 0\le s<e,
\]
where $U$ is homogeneous, $V$ is a coefficient-one monomial of the
same degree $d\ge1$, and
\[
 (d+1)p^s\le Q,\qquad a_i\in\Fp,
 \qquad 1+\sum_i a_i=0.
\]
None of these data is given. A deterministic polynomial sparse-bit
algorithm recognizes and exactly verifies such a representation whenever
one exists, normalizes $U$ to exclude the monomial $V$, and replaces
$F$ by the nonzero polynomial
\begin{equation}
 G=U^p-V^{p-1}U.
 \label{sp:eq:zero-sum-output}
\end{equation}
It satisfies
\[
 K\subseteq K'\subseteq I,
 \qquad (K')^{\mathrm{sat}}=I,
 \qquad \deg G=\frac{p\deg F}{Q}<\deg F.
\]
No coordinate index is added, and the output has at most twice the
input support. A deterministic convention can minimize the output
degree among passing representations. Despite this possible increase,
repeated use has polynomial total sparse-bit cost, also when
interleaved with the preceding coordinate, digit, linearized and
affine reductions. If $T_0$ is the original total support and $D_0$
the largest original degree, the total support throughout any such
interleaved computation is at most
\[
 16T_0\max\{1,\log_p D_0\}.
\]
\end{theorem}
\begin{proof}
Put $D=\deg F$. The translation $U\mapsto U-cV$, $c\in\Fp$,
leaves the relation unchanged because its coefficient sum is zero.
Remove the $V$-term of $U$. The normalized $U$ is nonzero, since
otherwise $F=0$. The same gap calculation as before gives
\[
 U=E_Q(F),
\]
where $E_Q$ is the divided coordinatewise $Q$-section. Indeed,
$Qv+p^i(u-v)$ is coordinatewise divisible by $Q$ only when $u=v$.
Thus the inner polynomial is determined at every candidate scale.

Enumerate powers $Q>p$ dividing $D$, and the admissible values of
$s$, all in polynomial number in $\log(D+1)$. Set $d=D/Q$, reject a
zero section, and form $R=F-U^Q$. A genuine normalized representation
has $R\ne0$. Otherwise $U/V$ in the rational function field would
satisfy the nonzero polynomial $\sum_i a_iT^{p^i}$, because
$\sum_i a_i=-1$. The polynomial ring is geometrically integral, so
its algebraic constants are $\Fp$. This would make $U$ a scalar
multiple of $V$, contrary to its normalization.

For completeness of the monomial search, put $q=p^s$ and $M=Q-q$,
and choose any exponent $\gamma$ in $R$. At least one contributing
lower-block term has
\[
 \gamma=(Q-p^i)v+p^iu
        =Mv+\beta,
 \qquad \beta=(q-p^i)v+p^iu\ge0,
 \qquad |\beta|=qd\le M.
\]
Therefore the coordinatewise floor $\lfloor\gamma/M\rfloor$ equals
$v$, or equals $v+e_j$ for a single coordinate $j$. Enumerate these
possibilities and keep the nonnegative vectors of degree $d$ whose
monomial is absent from $U$. This includes the equality case $qd=M$.
For each candidate solve
\[
 \sum_{i=0}^{s}a_iV^{Q-p^i}U^{p^i}=R,
 \qquad \sum_{i=0}^{s}a_i=-1
\]
by coefficientwise linear algebra, and verify the exact sparse identity.
There are $O(\log D)$ unknowns and
$O(|\operatorname{supp}F|\log D)$ equations. Frobenius powers and
monomial shifts construct every column without dense expansion.
The true representation passes. All accepted representations meet
the stated conditions. Decreasing order of $Q$ minimizes
$pD/Q$, with fixed further orders resolving ties. Every coordinate
of $V$ occurs in $R$, since its lower-block exponents are at least
$Mv$ coordinatewise; hence no new index is introduced.

Let $\overline U,\overline V$ be the images in $S/I$. If
$\overline V=0$, the original relation gives $\overline U=0$.
Otherwise $\overline U/\overline V$ is algebraic over $\Fp$,
hence belongs to $\Fp$. In either case $G\in I$.
To prove the other ideal inclusion, write
\[
 L(T)=\sum_{k=0}^e c_kT^{p^k},\qquad
 c_e=1,\quad \sum_{k=0}^e c_k=0,
 \qquad b_j=-\sum_{k=0}^j c_k.
\]
Here $c_k=a_k$ for $k\le s$ and is zero at the other lower indices.
Adjacent coefficients give
\[
 L(T)=\sum_{j=0}^{e-1}b_j(T^p-T)^{p^j}.
\]
Homogenizing to degree $Q$ in $U,V$ yields the exact identity
\begin{equation}
 F=\sum_{j=0}^{e-1}b_j
       V^{Q-p^{j+1}}G^{p^j}.
 \label{sp:eq:zero-sum-divisibility}
\end{equation}
All exponents of $V$ are nonnegative. Thus $F\in(G)$, proving
$K\subseteq K'\subseteq I$ and equality after saturation.
The polynomial $G$ is nonzero: if it vanished in the ambient
polynomial ring, the nonconstant rational function $U/V$ would be a
root of $T^p-T$, again impossible. Formula~\eqref{sp:eq:zero-sum-output}
gives its degree and its support bound.

It remains to control iteration. Let $T_0$ be the original total
support, $D_0$ the largest original degree, and $v$ the number of
initially occurring coordinates. Every current nonzero row belongs
to the proper homogeneous ideal $I$, so its degree is positive.
The gap and $Q>p$ give
\[
 D=dQ>d^2,\qquad D_{\mathrm{new}}=pd<p\sqrt D,
 \qquad D_{\mathrm{new}}\le D/p.
\]
Along a lineage put $z=\log_p D$. At a zero-sum step,
\[
 z_{\mathrm{new}}<1+\tfrac12z_{\mathrm{old}},
 \qquad z_{\mathrm{new}}\le z_{\mathrm{old}}-1.
\]
Intervening earlier operations only decrease degree. After $j\ge1$
zero-sum steps the first inequality therefore gives
$z_j<2+(z_0-2)/2^j$. After
$\lceil\log_2\max\{1,z_0\}\rceil$ such steps, if that many occur,
the degree is below $p^3$; at most three further steps are possible
by the second inequality. Thus every lineage has at most
\[
 k_0=\left\lceil\log_2\max\{1,\log_p D_0\}\right\rceil+3
 \quad\text{zero-sum steps},\qquad
 2^{k_0}\le16\max\{1,\log_p D_0\}.
\]

Give each original term a formal token. Across a
support-nonincreasing operation, assign all new terms injectively
to old tokens, including terms in different affine output rows.
Across a zero-sum step assign at most two new tokens to each old
token. Unassigned tokens terminate. These assignments only count
terms and assert no matching of polynomial coefficients. A token
follows the degree of its current row. Branching occurs only at a
zero-sum step, so each initial token has at most $2^{k_0}$ descendant
leaves, including terminated tokens. The current tokens form an
antichain in this forest. Consequently the total support, across
all rows and all splits, satisfies
\begin{equation}
 T_{\max}\le T_0 2^{k_0}
 \le16T_0\max\{1,\log_p D_0\}.
 \label{sp:eq:zero-sum-support-budget}
\end{equation}

To count all non-coordinate operations, suppress only
coordinate-content transitions in this forest and retain every
other transition, including unary transitions of earlier rules.
Every nontrivial remaining step makes each output degree at most
$\rho=1-1/(v+2)$ times its input degree: this is the digit guarantee,
while the other steps divide degree by at least $p$. No operation
introduces a coordinate, so the active count never exceeds $v$.
Since $-\log\rho\ge1/(v+2)$, every root-to-leaf path has at most
\[
 h_0=\left\lfloor(v+2)\log D_0\right\rfloor
\]
edges. A tree with at most $2^{k_0}$ leaves and height $h_0$ has at
most $2^{k_0}(h_0+1)$ vertices, because its root-to-leaf paths cover
every vertex. Each successful non-coordinate row operation has a
nonzero output and can be charged to one of its new token edges;
distinct operations create distinct edges. The total number of
these operations is therefore
\[
 O\!\left(T_0(v+2)(1+\log D_0)^2\right).
\]
This height-times-leaves count includes all unary earlier-rule
transitions, so no vertex bound is inferred from branching depth
alone.

Perform maximal eligible coordinate-content cleanup initially and
between successive non-coordinate operations. Removing one
coordinate's complete content does not create content in another
coordinate, so a batch has at most $v$ removals per current row.
The number of rows is at most $T_{\max}$, and the number of batches
is bounded by one plus the operation count just proved. Thus
coordinate cleanup also has polynomial cost. Coefficients remain
in $\Fp$, coordinate indices never grow, and exponent lengths are
bounded by the original degree bit length. Every intermediate list
has polynomial sparse-bit size. Complete recognition scans have
polynomial cost in that size; a successful scan is charged to a
successful operation or cleanup batch, and a complete unsuccessful
scan terminates. This proves the full interleaved bit bound.
\end{proof}

Support can increase in one allowed step. Over $\F_3$, let $m\ge3$,
$d=3^{2(m-1)}$, $V=y^d$, and
\[
 U=\sum_{j=0}^{m-1}(-1)^j
       x^{3^{2j}}y^{d-3^{2j}},\qquad Q=3^{2m+1}.
\]
Take $a_0=a_2=1$, with the other lower coefficients zero. The two
lower blocks telescope to two endpoint terms, so $F$ has $m+2$
terms. The two exponent sets in $U^3-V^2U$ consist of odd and even
powers of three and are disjoint, giving $2m$ terms. At $m=3$ this
is a five-term input and a six-term output. A different passing
representation may give a smaller output; the example shows why a
uniform iteration proof cannot simply assume support nonincrease.

\begin{example}[A strict family for the zero-sum reduction]
\label{sp:ex:zero-sum-family}
Over $\F_2$, let $n\ge5$ be odd and put
\[
 g=\sum_{j=0}^{n-1}x_j^n+x_0^{n-1}x_1,
 \qquad \mathcal M=\operatorname{supp}(g),
 \qquad Q=2^r>n.
\]
Here a member of $\mathcal M$ denotes its coefficient-one monomial.
Let $K$ be generated by the displayed rows
\begin{equation}
 g^Q-V^{Q-1}g,\qquad
 g^{2Q}-V^{2Q-1}g,
 \qquad V\in\mathcal M.
 \label{sp:eq:zero-sum-family}
\end{equation}
Then $K^{\mathrm{sat}}=(g)$, whose projective hypersurface in
$\PP^{n-1}$ is smooth, geometrically integral and anticanonical.
The preceding coordinate, digit, nonzero-sum and affine reductions
leave this displayed list unchanged, and its least principal
complete-intersection excess is $n(Q-1)$.
Theorem~\ref{sp:thm:zero-sum-preprocessing}, choosing minimum output
degree, replaces its rows by $g(g+V)$, of degree $2n$ and with an
$n$-term degree-$n$ cofactor. The designated sparse-cofactor algorithm
therefore computes the scalar in polynomial bit time.
For $n=2r+1$, $r\ge3$, the input length is polynomial in $r$ while
the earlier residual numerical cofactor degrees are exponential in $r$.
\end{example}
\begin{proof}
The partial derivatives of $g$ are
\[
 g_{x_0}=x_0^{n-1},\qquad
 g_{x_1}=x_1^{n-1}+x_0^{n-1},\qquad
 g_{x_j}=x_j^{n-1}\quad(j\ge2).
\]
Their only common affine zero is the origin, proving projective
smoothness. Positive-degree factors would intersect projectively
and yield a singular point; repeated factors are likewise excluded.
Thus $g$ is geometrically irreducible. Its degree is the number $n$
of homogeneous coordinates, so it is anticanonical.

The cofactors $g^{P-1}-V^{P-1}$, $P=Q,2Q$, have no common
projective zero. Where $g=0$, vanishing of all cofactors forces all
$V$ to vanish, including every $x_j^n$. Where $g\ne0$, their
vanishing gives
\[
 (V/g)^{Q-1}=(V/g)^{2Q-1}=1.
\]
The two exponents are coprime, so $V/g=1$ for all $V\in\mathcal M$.
But $|\mathcal M|=n+1$ is even, contradicting
$1=g/g=\sum_{V\in\mathcal M}V/g$. The cofactor ideal sheaf is
therefore the unit sheaf, proving the scheme-theoretic saturation.
The least nonzero degree in $K$ is $nQ$, attained by a row of
scale $Q$, so its principal excess is $nQ-n$.

Fix a row of scale $P=Q$ or $2Q$. Its pure top terms include
$x_j^{nP}$ for every pure monomial different from $V$.
Only $V^P$ cancels between the two blocks: any further cancellation
would require $P(\alpha-v)=u-v$, with every coordinate on the
right of absolute value less than $P$. At least $n-1$ pure top
terms remain, so the monomial content is one.

Consider a digit scale $A=2^a\le P$. Every pure top term has low
digit zero. Choose a pure monomial $x_j^n$ whose coordinate is
outside the support of $V$. Its mixed term has $j$-exponent $n$,
which is odd, and therefore gives a nonzero low row. Since the
common total degree $nP$ is divisible by $A$, a column of high
total degree $nP/A$ can occur only in low row zero. A surviving
pure top column and any nonzero entry in the nonzero low row form
a triangular two-by-two minor with nonzero diagonal. The digit
rank is at least two.

If $P<A\le nP$, write $A=Pb$, with $b$ a power of two satisfying
$2\le b\le n$. Choose two indices $j\ge2$ whose pure monomials
are not $V$; $n\ge5$ guarantees two choices. Their top terms have
distinct low rows $P(n\bmod b)e_j$ and distinct high columns
$\lfloor n/b\rfloor e_j$. Both factors are nonzero because $n$
is odd and $b\le n$. They form an identity minor. The cross
monomial involves only indices zero and one. Every remaining mixed
term has a coordinate not divisible by $P$, since $P>n$ and
$u\ne v$; it cannot enter those low rows. Thus the minor persists.
For $A>nP$ the high factor is constant. No nontrivial digit
reduction applies at any scale.

In characteristic two a nonzero additive coefficient sum equals
one. Its top-section formula forces the inner polynomial of any
passing earlier relation at scale $A$ to be $E_A(F)$. Because $n$
is odd and $A$ divides $nP$, necessarily $A\le P$. Here
\[
 E_A(F)=g^{P/A}-H,
 \qquad H=E_A(V^{P-1}g).
\]
The term $V$ contributes the nonzero monomial $V^{P/A}$ to $H$.
The other pure monomials have different residues modulo $A$.
The cross monomial shares a residue with a pure monomial only for
$x_1^n$, and then only if $A\mid n-1$. Thus $H$ is nonzero and
involves either one coordinate or only $x_0,x_1$. It cannot be a
multiple of $g$, which has positive degree in every other coordinate.
Consequently $E_A(F)\notin(g)$, contradicting the safety conclusion
of any passing nonzero-sum relation. An affine no-root relation is
also impossible: it would put a positive-degree monomial in $(g)$,
although this prime ideal contains no coordinate. None of the
earlier rules changes the row, so their whole displayed-row closure
leaves the list unchanged.

For the new rule take $U=g+V$, $s=0$, and $a_0=1$. The degree gap
is $n+1\le P$, and its output is
\[
 (g+V)^2-V(g+V)=g(g+V).
\]
No smaller output degree is possible, since $P$ is the largest
power of two dividing $nP$. At this scale $U=E_P(F)=g+V$ is fixed.
If a different monomial $W$ gave a passing output, safety modulo
$(g)$ would imply $V(V+W)=0$, hence $V+W\in(g)$. This is a
polynomial of degree $n$ with at most two terms, whereas $g$ has
$n+1\ge6$ terms. It must be zero, so $W=V$. Thus the
minimum-degree output is precisely the stated row. Its cofactor
$g+V$ has degree $n$ and $n$ terms, meeting
Theorem~\ref{sp:thm:sparse-cofactor} with explicit polynomial bounds.

Finally, there are $2(n+1)$ original rows, each with $2n$ terms and
exponent bit length $O(r+\log n)$. Setting $n=2r+1$ gives polynomial
input length in $r$, while $n(Q-1)$ is exponential. This comparison
concerns the specified preprocessing and cofactor conditions; it
asserts no hardness of the scalar on this family.
\end{proof}

\subsection{Explicit relations and strict examples}
The general recognizer includes two useful identities with short sparse
forms. The following corollary states their hypotheses explicitly.

\begin{corollary}[Two explicit prime-field relations]
\label{sp:thm:constant-field-preprocessing}
Fix a prime $p$. Let a homogeneous sparse list with binary exponents
generate $K\subset S=\Fp[x_0,\ldots,x_N]$, and suppose that
$I=K^{\mathrm{sat}}$ has geometrically integral quotient $S/I$.
For a nonzero listed row $F$, consider the identities
\begin{align}
 F&=U^P+aV^{P-1}U,
       &&a\in\Fp\setminus\{0,-1\},                 \label{sp:eq:constant-two}\\
 F&=U^{P^2}+V^{P^2-P}U^P-V^{P^2-1}U,              \label{sp:eq:constant-three}
\end{align}
where $P=p^r$, $r\ge1$, $U$ is nonzero and homogeneous of degree
$1\le d<P$, and $V$ is any monomial of degree $d$ with coefficient
one. The polynomials $U,V$, the integer $r$, and the scalar $a$ are
not part of the input.

There is a deterministic algorithm of polynomial sparse-bit cost that
finds a valid identity of either displayed type whenever one exists,
and verifies the identity before replacing $F$ by $U$. It may choose
a replacement of minimum degree among all the identities it finds.
Every accepted replacement gives
\[
 K\subseteq K'\subseteq I,\qquad (K')^{\mathrm{sat}}=I.
\]
It introduces no new coordinate index and does not increase the number
of terms. The degree decreases by $P$ in
\eqref{sp:eq:constant-two} and by $P^2$ in
\eqref{sp:eq:constant-three}. Repeated replacements, interleaved with
rank-one digit reduction and any valid coordinate-content divisions,
terminate in polynomial sparse-bit time. The first identity has no
allowed value of $a$ when $p=2$; the second applies in that
characteristic as well.
\end{corollary}
\begin{proof}
For \eqref{sp:eq:constant-two}, use
Theorem~\ref{sp:thm:linearized-preprocessing} with $Q=P$, $s=0$
and $a_0=a$. Its gap condition is $d+1\le P$, and its coefficient
sum is $1+a\ne0$. For \eqref{sp:eq:constant-three}, write $P=p^r$,
take $Q=P^2$, $s=r$, $a_r=1$, $a_0=-1$, and set all other lower
coefficients to zero. The gap again says $d+1\le P$, and the total
coefficient sum is one. This includes characteristic two.
The recognition, minimum-degree choice and iteration bounds are
those of the theorem; neither $V$ nor its membership in the support
of $U$ is required as input.
\end{proof}

\begin{example}[A strict improvement in growing dimension]
\label{sp:ex:constant-field-family}
Fix $p=7$, let $n\ge3$ with $7\nmid n$, and choose $P=7^r>n$.
In $\Fp[x_0,\ldots,x_{n-1}]$ put
\[
 g=\sum_{i=0}^{n-1}x_i^n,\qquad
 F_j=gx_j^{n(P-1)}+g^P\quad(0\le j<n),\qquad K=(F_j).
\]
Then $K^{\mathrm{sat}}=(g)$, and $g$ defines a smooth geometrically
integral anticanonical hypersurface. Every row has $2n-1$ terms and
degree $nP$. No coordinate-content division or nontrivial rank-one
digit reduction changes any row. The least degree of a nonzero
homogeneous equation in $K$ is $nP$, giving least hypersurface excess
$n(P-1)$. Theorem~\ref{sp:thm:linearized-preprocessing} replaces
the rows by scalar multiples of $g$, giving excess zero.
\end{example}
\begin{proof}
The partial derivatives $n x_i^{n-1}$ have no common projective zero.
A reducible projective hypersurface of positive dimension has
intersecting geometric components, so this smooth hypersurface is
geometrically integral. Its degree $n$ is anticanonical in
$\PP^{n-1}$.

Write $F_j=gh_j$, with $h_j=x_j^{n(P-1)}+g^{P-1}$. The identity
\[
 \sum_jx_j^n h_j=2g^P
\]
shows that a common projective zero of the $h_j$ would have $g=0$
and then every $x_j=0$. Thus the cofactors contain a power of the
irrelevant ideal, proving $K^{\mathrm{sat}}=(g)$.
The expansion of $F_j$ has coefficient two at $x_j^{nP}$,
coefficient one at every other pure power, and the $n-1$ distinct
cross terms $x_i^n x_j^{nP-n}$. In particular its monomial content is
one.

Consider a digit radix $Q=7^s$. If $Q\le P$, write $n=aQ+b$,
where $1\le b<Q$. The pure terms have zero low digit and pure high
monomials of degree $nP/Q$. A cross term has low monomial
$x_i^b x_j^{Q-b}$ and high monomial
$x_i^a x_j^{nP/Q-a-1}$, of degree $nP/Q-1$.
The nonzero pure row and a cross row therefore give a nonzero
two-by-two minor with columns of different total degrees. If
$P<Q\le nP$, write $Q=Pq$ and $n=aq+b$, with $a\ge1$ and
$1\le b<q$. Pure terms have distinct low monomials $x_i^{bP}$ and
distinct high monomials $x_i^a$. Cross terms have low monomials
$x_i^n x_j^{bP-n}$ involving two coordinates, since $P>n$.
They do not alter the pure rows, which again have rank at least two.
For $Q>nP$ every high factor is constant. This covers every possible
digit scale.

The ideal is generated in degree $nP$, and a nonzero row attains that
degree. This proves the stated minimum in $K$, not in its saturation.
Identity~\eqref{sp:eq:constant-two} holds with $U=g$, $V=x_j^n$,
and $a=1$. Every accepted replacement belongs to $(g)$, so it has
degree at least $n$. Minimum-degree selection therefore gives a
degree-$n$ associate of $g$. Taking, for example, $n=7r+1$ for
$r\ge2$ gives growing dimension, polynomial written sparse size in
$r$, and exponentially growing former excess.

For any fixed $c\ge2$, adjoin $c-1$ variables $w_i$ and include the
linear rows $w_i$. The saturation becomes $(w_1,\ldots,w_{c-1},g)$
in $\PP^{n+c-2}$, with anticanonical degree sum $n+c-1$.
Before the new reduction, every equation of degree less than $nP$
belongs to the height-$c-1$ linear ideal. A regular $c$-sequence must
therefore have degree sum at least $nP+c-1$, and the linear rows
together with one $F_j$ attain it. After reduction the excess is zero
and the deterministic linear-prefix theorem applies.
\end{proof}

\subsection{Exponential support in a Cartier iterate}
Smooth saturation alone does not bound the number of monomials
reached by repeatedly applying the sparse Cartier operator. The next
family survives coordinate and rank-one digit reduction, while
Theorem~\ref{sp:thm:linearized-preprocessing} recognizes its
prime-field relation. It separates a limitation of explicit support
enumeration from the scalar computation itself.

Multiply-then-section operators also appear in algorithms for
diagonals of rational functions. \citet[Definition~1.7,
Proposition~1.9, Algorithm~1 and Remark~2.2, pp.~249--254]{sp-rowland-yassawi}
construct polynomial numerator states whose finiteness is controlled
by numerical coordinate degrees. The coefficient algorithms of
\citet[Theorems~3.4 and~4.1, author pp.~9--11]{sp-bccd} have
logarithmic dependence on the requested coefficient index, together
with polynomial dependence on numerical degree and height.
\citet[Theorem~1.2 and Corollary~1.4, pp.~4--5; \S2.3]{sp-effective-christol}
bound an invariant linear space by lattice points in a Newton-polytope
region. Those numerical bounds can be large for a sparse list with
binary exponents. Conversely, a monomial-support graph is different
from a minimal automaton or an invariant linear space. The example
below asserts a lower bound only for the specified explicit monomial
representation, and supplies no lower bound for those compressed
representations.

\begin{proposition}[A dense iterate from one monomial]
\label{sp:prop:cartier-states}
Over $\mathbb F_7$, let $P=7^r$, $r\ge1$, and set
\[
 g=x^3+y^3+z^3,\qquad b=3(P-1),\qquad
 f_j=gx_j^b+g^P\quad(x_0=x,x_1=y,x_2=z).
\]
The ideal $K=(f_0,f_1,f_2)$ has smooth geometrically integral
anticanonical saturation $(g)$ and sparse binary input length
$\Theta(r)$. Each cofactor $h_j=f_j/g$ has exactly $28^r$ terms.
Every nonzero constant combination of the rows, even over a field
extension, has monomial content one and has no nontrivial rank-one
base-$7^s$ digit reduction. The least value of $\deg F-3$ for a
nonzero homogeneous $F\in K$ is $b$.

Let $T_f(q)=\mathcal C_7(f_0^6q)$ on homogeneous forms of degree $b$.
Then
\[
 T_f^r((xyz)^{P-1})=(2x^3+y^3+z^3)^{P-1},\qquad
 T_f(h_0)=6h_0.
\]
The first polynomial has exactly $28^r$ nonzero terms. Consequently
an algorithm that explicitly enumerates the reachable monomial states
from $(xyz)^{P-1}$ must handle exponentially many states in the
written sparse input length. This is not a lower bound for computing
the scalar or for using compressed polynomial representations.
\end{proposition}
\begin{proof}
Example~\ref{sp:ex:constant-field-family}, with $n=3$, proves the
saturation, geometry, five-term rows and minimum degree. For nonzero
$c_0,c_1,c_2\in\mathbb F_7$,
\[
 (c_0x^3+c_1y^3+c_2z^3)^{P-1}
  =\prod_{i=0}^{r-1}
    (c_0x^{3\cdot7^i}+c_1y^{3\cdot7^i}
                         +c_2z^{3\cdot7^i})^6.
\]
Each factor has $\binom82=28$ terms, with nonzero multinomial
coefficients modulo seven. After dividing the exponents by three,
their digits at place $i$ are the three chosen nonnegative integers
of sum six. Uniqueness of base-seven expansion prevents collisions
between distinct digit choices. This proves the support count $28^r$.
In $h_j=x_j^b+g^{P-1}$ only one pure coefficient changes from one to
two, so no term disappears.

For a nonzero constant vector $(a_0,a_1,a_2)$, put
$\sigma=a_0+a_1+a_2$. Its row combination is
\[
 F=\sum_j(a_j+\sigma)x_j^{3P}
       +\sum_{j\ne i}a_jx_j^{3P-3}x_i^3.
\]
The pure part is nonzero: its vanishing would give $a_j=-\sigma$
for every $j$ and hence $4\sigma=0$, forcing all coefficients to
vanish. If only one $a_j$ is nonzero, all three pure powers occur.
If at least two are nonzero, then for any coordinate $x_k$ an active
$j\ne k$ and the third index $i$ supply a cross term omitting $x_k$.
Thus the monomial content is one in either case.

At radix $Q=7^s\le P$, write $D=3P/Q$. The zero-residue section is
$\sum_j(a_j+\sigma)X_j^D$. For every active $a_j$, the two
residues $(Q-3)e_j+3e_i$ have sections $a_jX_j^{D-1}$.
These nonzero residues are distinct. If $t$ coefficients $a_j$ are
nonzero, their sections span the $t$ independent monomials
$X_j^{D-1}$, and the nonzero degree-$D$ section is independent of
them. The digit matrix has exact rank $t+1\ge2$.
For $Q>P$, one has $Q\ge7P>3P$, so the high factor is constant.
All these arguments remain valid over extension fields.

The monomial definition of Cartier extraction gives
$\mathcal C_Q(A^QB)=A\mathcal C_Q(B)$ and
$\mathcal C_Q\mathcal C_{Q'}=\mathcal C_{QQ'}$ for powers of seven.
The operator preserves degree $b$, since
$\bigl(6\cdot3P+b-18\bigr)/7=b$.
Induction therefore gives
\[
 T_f^s(q)=\mathcal C_{7^s}(f_0^{7^s-1}q).
\]
Put $A=2x^3+y^3+z^3$ and $B=y^3+z^3$, so that
$f_0=A^P+x^bB$. In the binomial expansion of $f_0^{P-1}$,
a term with $k$ factors from $x^bB$, $1\le k<P$, has $x$-exponent
congruent to $-3k$ modulo $P$, after pulling out its $P$th-power
factor. After multiplication by $(xyz)^{P-1}$ it cannot survive
$\mathcal C_P$, since $P\nmid3k$. The $k=0$ term gives
$A^{P-1}$. This proves the exact iterate, including cancellations.
Its $28^r$ nonzero monomials each have a nonzero transition path in
the matrix product defining $T_f^r$. They are therefore distinct
reachable states. Their ordinary explicit binary exponent lists
require $\Omega(r28^r)$ bits, whereas the input has length
$\Theta(r)$.

Finally,
\[
 T_f(h_0)=\mathcal C_7(g^6h_0^7)
          =h_0\mathcal C_7(g^6).
\]
The only surviving monomial of $g^6$ is $x^6y^6z^6$, with coefficient
$6!/(2!2!2!)=90\equiv6\pmod7$. This proves the eigenvalue assertion.
The equation $g$ and a short expression for the dense iterate are
explicit in this family; their existence is fully compatible with
the lower bound on expanded monomial output.
At this fixed prime and ambient dimension, evaluation of the original
rows at the $57$ points of $\PP^2(\mathbb F_7)$ is itself a
polynomial-time scalar computation. Thus the example concerns the
specified representation of the operator, not the unrestricted scalar
problem.
\end{proof}

\subsection{Certified removal of coordinate factors}
The finite fibres can certify further degree reductions before the
bounded-excess algorithm is used. All powers of a coordinate may be cancelled uniformly from equations
in the prime ideal $I_X$ whenever that coordinate does not vanish
identically on $X$. The unknown coordinates that do
vanish span a space of dimension at most $c-1$; the following algorithm
finds certificates without first computing $I_X$.

Fix a deterministic preprocessing convention on the displayed generating
list of $K$. For a set
$H\subseteq\{0,\ldots,N\}$, process each nonzero row by dividing out its
full monomial content in coordinates outside $H$. Then apply the
rank-one digit operation of Lemma~\ref{sci:lem:deflation}, choosing the
largest admissible $r\ge1$. Then apply a minimum-degree recognized
replacement from Theorems~\ref{sp:thm:linearized-preprocessing},
\ref{sp:thm:affine-preprocessing}, and~\ref{sp:thm:zero-sum-preprocessing}, breaking ties by a fixed
ordering. Normalize each output row by its first nonzero coefficient
and repeat on the resulting list until no operation applies. Let
$\mathcal P_H(K)$ denote the ideal of the resulting list. Here the digit
operation writes $f=AB^{p^r}$ with all exponents of $A$ less than $p^r$
and $B$ nonconstant, and replaces it by $AB$. It is recognized by the
rank-one coefficient matrix of low and high exponent digits. Zero rows
are skipped. This convention defines the processed ideal even when the
speculated coordinate divisions are not valid on $X$.
Discard any speculative branch that produces a nonzero constant.

\begin{theorem}[Certified coordinate division before saturation recovery]
\label{sci:thm:certified-preprocessing}
Fix a prime $p$, a codimension $c\ge2$, and $E\ge0$. Let a homogeneous
sparse list with binary exponents generate
$K\subset\Fp[x_0,\ldots,x_N]$, with
$K^{\mathrm{sat}}=I_X$, where $X$ is a smooth geometrically integral
anticanonical global complete intersection of dimension $N-c\ge1$.
Set
\[
 H_X=\{j:x_j\in I_X\},\qquad K^\circ=\mathcal P_{H_X}(K).
\]
Suppose that, after a finite field extension, $K^\circ$ contains a
regular $c$-sequence of degree sum at most $N+1+E$. Neither $H_X$ nor
that sequence is part of the input. For every $a\ge1$, a randomized
algorithm of bit cost polynomial in the written input length and $a$
returns the coherent Frobenius scalar or a failure symbol, with failure
probability at most $2^{-a}$ and no incorrect outputs on the promise.
Repetition has expected polynomial cost. Its only random choices are
projection centres. Every coordinate division used for the final
scalar is certified from the original equations on a deterministic
grid of at most $N+1$ affine fibres for each tested tuple and centre.
\end{theorem}
\begin{proof}
Write $I=I_X$, $n=N-c$, and $R=N+1+E$. The space $I_1$ has dimension
at most $c-1$. Otherwise all $c$ minimal complete-intersection
generators would be linear, contradicting their degree sum
$N+1\ge c+2$. Consequently $|H_X|\le c-1$, and there are
$\sum_{j=0}^{c-1}\binom{N+1}{j}$ possible sets to examine.
Every coordinate occurs in the original list. Indeed, if $x_j$ were
absent, multiplication by it would be injective on $S/K$, so saturation
would give $K=I$. The coordinate point $e_j$ would then be a cone vertex
whose tangent space has dimension $N-\dim I_1\ge N-c+1$, contradicting
smoothness of the $(N-c)$-dimensional variety. Thus $N+1$ itself is
bounded by the written sparse input length, even with sparse coordinate
labels.
For the true set, every division is valid: if $x_j\notin I$ and
$x_j^b u\in I$, primality gives $u\in I$. Digit reduction is valid by
radicality, and the recognized prime-field relations are valid by
geometric integrality. Every old row belongs to the ideal generated by its replacement rows, so the
ideal sandwich $K\subseteq K^\circ\subseteq I$ proves
$(K^\circ)^{\mathrm{sat}}=I$.

We first establish the certificate used even on a wrong branch. Let
$k/\Fp$ be a finite extension, let $t_0$ be any linear form over $k$,
and let $B$ be a finite-dimensional quotient of
\[
 S_k/(K_k,t_0-1).
\]
If the image of $x_j$ in $B$ is nonzero, then $x_j\notin I$. Indeed,
$x_j\in K^{\mathrm{sat}}$ would give
$\mathfrak m^b x_j\subset K$ for some $b$, and hence
$t_0^b x_j\in K_k$; the image would vanish after setting $t_0=1$.
The certificate is a nonzero algebra element. It need not be a unit,
and nilpotents are not discarded.

For each possible $H$, form the speculative list $\mathcal P_H(K)$.
Discard a branch containing a nonzero constant. Run the explicit candidate construction of
Lemma~\ref{sci:lem:profile-candidates} on all profiles of sum between
$N+1$ and $R$. For each candidate tuple, sample a centre and apply the
exact degree-bounded Macaulay certificate described in
Proposition~\ref{sci:prop:random-designated}.
A successful certificate proves that the tuple is regular and finite
over its projection base, whether or not the speculative list has the
correct saturation. Complete its coordinates as $x=Uz+Vt$, where
$t=(t_0,\ldots,t_n)$. For
\[
 \Gamma_1=\{0,e_1,\ldots,e_n\}\subset k^n
\]
with $e_i$ the standard coordinate vectors,
form the full finite algebra
\[
 A_b=k[z_1,\ldots,z_c]/
       (f_i(Uz+V(1,b)):1\le i\le c).
\]
Its dimension is $\Delta=\prod_i\deg f_i\le R^c$.
Evaluate every original input row in $A_b$ by binary powering and
quotient by the ideal of those values, obtaining $B_b$. This is exact
finite-dimensional linear algebra: if $v_1,\ldots,v_\Delta$ is an
algebra basis, the vectors $v_j u_i$ span the ideal generated by the
values of the original rows. Compute the image of every coordinate in
this quotient. Let $C$ collect the indices having a nonzero image in
at least one of these algebras, over all branches examined.
Each $B_b$ is a quotient of $S_k/(K_k,t_0-1)$, so every index in $C$
is certified safe even if its branch was wrong.

For completeness, take the true branch and one of its regular
bounded-excess candidates with a successful centre. Its ideal $J$ is
contained in $I$. The graded module $M=S_k/I_k$ is finite free over
$k[t_0,\ldots,t_n]$, with a homogeneous basis of nonnegative degrees,
by the finite-free argument of Lemma~\ref{sci:lem:fibres}.
On $t_0=1$, the original ideals $K_k$ and $I_k$ agree: a power of the
irrelevant ideal annihilates their quotient, and $t_0$ is invertible.
Therefore the certificate algebra is exactly
\[
 B_b=M\otimes_{k[t_0,\ldots,t_n]} k_{(1,b)}.
\]
In a homogeneous free basis, a degree-one coordinate has coefficient
polynomials of degrees at most one, because every basis degree is
nonnegative. If $x_j\notin I$, at least one coefficient is nonzero.
Dehomogenization at $t_0=1$ is injective on a fixed homogeneous degree,
and evaluation at $0,e_1,\ldots,e_n$ detects every nonzero affine linear
polynomial over every field. Thus one of these fibres certifies $j$.
On this event $C$ is exactly the complement of $H_X$.

Process the original list again, dividing only by coordinates in $C$
and applying digit and prime-field relation reductions with the fixed
convention. All these
operations now preserve saturation, independently of whether $C$ is
complete. Run Theorem~\ref{sci:thm:random-excess} on this safe list.
Every successful output is correct. Once the true branch has a good
centre, its list is exactly $K^\circ$, so the assumed bounded-excess
promise applies to this final run.

Assign failure budgets $2^{-a-1}$ to certificate discovery and to the
final scalar calculation. Only one regular candidate on the true branch
is needed for discovery; wrong branches cannot produce a false
certificate. The centre resultant has degree at most $cR^c$, so the
same finite-field sampling and repetition bound used in
Theorem~\ref{sci:thm:random-excess} gives its required failure budget.
A union bound between the two stages proves the stated probability.
Both stages return a distinguished failure when their bounded trials
are exhausted; no unbounded search is hidden in this assertion.

It remains to check the binary-degree cost of preprocessing.
Theorem~\ref{sp:thm:zero-sum-preprocessing} bounds both total support
and the number of steps polynomially through the entire interleaving,
including affine row splitting and zero-sum support increases.
These combinatorial bounds apply on speculative branches up to rejection
as well: all surviving rows are homogeneous of positive degree,
regardless of membership in $I$.
Its token count uses the stronger square-root degree decrease at
every potentially support-increasing step. Full eligible monomial
content is removed in maximal batches between degree reductions.
Testing every possible $r$ uses only
binary exponent arithmetic and sparse rank-one tests. The number of
branches, profiles, algebra dimensions, coordinate tests and fibre
points is polynomial for fixed $c,E$. Original large exponents occur
only in binary powers of finite-algebra elements. Finite-field
construction and scalar descent have the already established polynomial
bit bounds. This completes the proof.
\end{proof}

\begin{corollary}[A deterministic processed linear prefix]
\label{sci:cor:certified-linear}
Under the hypotheses of Theorem~\ref{sci:thm:certified-preprocessing},
if $c\ge2$ and $\dim(K^\circ)_1=c-1$, the scalar is computable
deterministically in polynomial bit time.
\end{corollary}
\begin{proof}
Enumerate the same possible protected-coordinate sets. On each branch
whose degree-one part has dimension $c-1$, run the bounded-loop
centre construction of Theorem~\ref{sci:thm:linear-prefix}, testing
its exact centre certificate and returning failure if any required
step fails. The true branch has a regular bounded-excess tuple and
succeeds. Its degree-one fibre grid certifies every safe coordinate.
All certificates from other branches remain sound by the preceding
proof. Reprocessing with their union therefore gives $K^\circ$.
Apply the deterministic linear-prefix algorithm to that list. Each
stage has a fixed polynomial bound even on a branch failing its
promises; only verified centres and coordinate certificates are used.
\end{proof}

\begin{example}[Unbounded excess before certified division]
\label{sci:ex:certified-fermat}
Fix $p$ and choose arbitrarily large $N\ge3$ with $p\nmid N$ and
$N\equiv1\pmod{p-1}$, for example $N=1+p(p-1)k$, $k\ge1$.
In $\Fp[w,x_1,\ldots,x_N]$ put
\[
 F=\sum_{i=1}^N x_i^N,\quad I=(w,F),\quad
 z=(w,x_1,\ldots,x_N),\quad
 K=(z_j^{N-1}w,z_j^{N-1}F:0\le j\le N).
\]
Then $K^{\mathrm{sat}}=I$, its anticanonical variety is smooth and
geometrically integral, and the least complete-intersection excess
after digit reduction alone is exactly $N-1$. Certified coordinate
division gives excess zero and the deterministic case of
Corollary~\ref{sci:cor:certified-linear}.
\end{example}
\begin{proof}
On every projective coordinate chart a cover factor is a unit, so the
saturation is $I$. The nonzero gradient entries $Nx_i^{N-1}$ prove
smoothness of the degree-$N$ hypersurface in $w=0$. Two distinct
positive-degree factors would intersect in projective space and give
a singular point; repeated factors are also impossible. Thus it is
geometrically integral. The degrees $1,N$ sum to $N+1$.

The row $w^N$ digit-reduces to $w$: a positive monomial exponent keeps
its residue modulo $p-1$, and a terminal positive exponent lies between
one and $p-1$. Every $F$-row is unchanged by a nontrivial rank-one digit
test. To see this, fix $P=p^r$ and write $N=dP+b$ when $P\le N$.
Then $d\ge1$ and $1\le b<P$, since $p\nmid N$.
For $x_j^{N-1}F$, the terms with $i\ne j$ have distinct low digits
$x_j^{b-1}x_i^b$ and distinct high digits $x_j^d x_i^d$.
No other term shares these rows or columns of the coefficient matrix,
so two such terms give rank at least two. They exist because $N\ge3$.
The $w$-cover row has the same property with its $N$ distinct
$x_i$ terms. If $N<P\le2N-1$, the internal row has one term with high
digit $x_j$ and the remaining terms have high digit one; their low
supports are disjoint, so the rank is again two. The $w$-cover row
has constant high factor in this range. For $P>2N-1$ all high factors
are constant. These are all possible digit tests.

After digit reduction the ideal therefore contains $w$, and every
remaining row outside $(w)$ has degree $2N-1$. Every homogeneous
element of smaller degree lies in $(w)$. A regular pair consequently
has degree sum at least $1+(2N-1)=2N$, attained by
$(w,x_1^{N-1}F)$. Its excess is $N-1$. On the other hand,
$H_X=\{0\}$, where index $0$ denotes the coordinate $w$, and safe
division removes $x_i^{N-1}$ from an internal
$F$-row. The fully processed ideal is $(w,F)$, of excess zero and
with a one-dimensional degree-one part, as asserted.
\end{proof}

\InputIfFileExists{sections/structured.tex}{}{}
\InputIfFileExists{sections/arrangements.tex}{}{}
\InputIfFileExists{sections/fixed_fields.tex}{}{}
\InputIfFileExists{sections/motives.tex}{}{}
\InputIfFileExists{sections/small_zeta.tex}{}{}
\InputIfFileExists{sections/questions.tex}{}{}

\section*{Method and records}
The mathematics and text were produced by GPT-6 Astra (ChatGPT Pro mode). Separate instances of the same model re-proved specified statements before reading their proofs; other statements received targeted proof checks and finite computations. Earlier records identify checks of the binary incidence and generic smoothing mechanisms by Claude Opus 5.5 at effort high, and statement-only readings of the subsequent constructions. Claude Opus 5.5 at effort max proposed the height-one formal-group observation; the geometric deduction, higher-height assertions and derivative applications are proved and separately checked here. The event register distinguishes these checks from formal verification. No human has checked this paper. The text and the original accompanying programs are dedicated to the public domain under CC0 1.0 Universal, \url{https://creativecommons.org/publicdomain/zero/1.0/}. The finite computations use CoreChecks, SparsePresentationChecks, StructuredVerify, StructuredAlgorithms, FixedFieldReplay, FixedInvariantsReplay, DistributedPaddingChecks, SparseLineGCDChecks, SparseDigitDeflationChecks, SaturationCIChecks, SlicedColonChecks, WalshTransferChecks, PrimeFieldPreprocessingChecks, and CartierSupportChecks. These programs, the statement register, check events, and reproducible certificates are organized for the workbench at \url{\workbenchurl}; the permanent address of the records is \url{\recordurl}.

\begin{thebibliography}{99}
\bibitem[Katz(1972)]{core-katz}
Katz, N. M. (1972). Algebraic solutions of differential equations ($p$-curvature and the Hodge filtration). \emph{Inventiones Mathematicae, 18}, 1--118. Algorithm (2.3.7.14), Special Case (2.3.7.17), pp. 35--36. \url{https://web.math.princeton.edu/~nmk/old/algsoln.pdf}.

\bibitem[Wan(2008)]{core-wan}
Wan, D. (2008). Modular counting of rational points over finite fields. \emph{Foundations of Computational Mathematics, 8}, 597--605. Theorem 1.1, p. 3 of the author manuscript. \url{https://www.math.uci.edu/~dwan/sparse.pdf}.

\bibitem[Huang et~al.(2018)]{core-hlyy}
Huang, A., Lian, B., Yau, S.-T., \& Yu, C. (2018). \emph{Hasse--Witt matrices, unit roots and period integrals} (arXiv:1801.01189). Corollary 2.5 and Remark 2.6, p.~9 of the author manuscript. \url{https://arxiv.org/abs/1801.01189}.

\bibitem[Knutson(2009)]{core-knutson}
Knutson, A. (2009). Frobenius splitting, point-counting, and degeneration. arXiv:0911.4941. Theorems 1--2. \url{https://arxiv.org/abs/0911.4941}.

\bibitem[Poonen(2004)]{core-poonen}
Poonen, B. (2004). Bertini theorems over finite fields. \emph{Annals of Mathematics, 160}(3), 1099--1127. Section 2.3. \url{https://arxiv.org/abs/math/0204002}.

\bibitem[Wan and Zhang(2025a)]{root-wz-betti}
Wan, D., \& Zhang, D. (2025a). \emph{Betti number bounds for varieties and exponential sums} (arXiv:2501.12623v1). Theorem 1.1.2, p.~2; Corollary 5.5.8, p.~38. \url{https://arxiv.org/abs/2501.12623v1}.

\bibitem[Wan and Zhang(2025b)]{root-wz-visibility}
Wan, D., \& Zhang, D. (2025b). Revisiting Dwork cohomology: Visibility and divisibility of Frobenius eigenvalues in rigid cohomology. \emph{Compositio Mathematica, 161}, 1215--1249. Theorem 1.2; Theorem 1.3, p.~2, in arXiv:2207.12633v2. \url{https://doi.org/10.1112/S0010437X25007110}.

\bibitem[DiPippo and Howe(1998)]{root-dipippo-howe}
DiPippo, S. A., \& Howe, E. W. (1998). Real polynomials with all roots on the unit circle and abelian varieties over finite fields. \emph{Journal of Number Theory, 73}, 426--450. Corrigendum, \emph{83} (2000), 182. Theorems 1.1--1.2, pp.~1--2 of the corrected author manuscript. \url{https://arxiv.org/abs/math/9803097v3}.

\bibitem[OpenAI(2026)]{root-openai-permanent}
OpenAI. (2026). Circuit and formula lower bounds for the permanent. In \emph{Ten advances in mathematics and theoretical computer science} (Chapter~5, updated August~6, 2026). Theorems 1.1--1.3, pp.~115--116. \url{https://cdn.openai.com/pdf/ten-proofs-oai.pdf}.

\bibitem[Rosser and Schoenfeld(1962)]{root-rosser-schoenfeld}
Rosser, J. B., and Schoenfeld, L. (1962).
Approximate formulas for some functions of prime numbers.
\emph{Illinois Journal of Mathematics, 6}(1), 64--94.
\url{https://doi.org/10.1215/ijm/1255631807}.

\bibitem[Yobuko(2019)]{fg-yobuko}
Yobuko, F. (2019).
Quasi-Frobenius splitting and lifting of Calabi--Yau varieties in
characteristic $p$.
\emph{Mathematische Zeitschrift, 292}(1--2), 307--316.
\url{https://doi.org/10.1007/s00209-018-2198-7}.
The cited preprint is arXiv:1704.05604v2 (2018).

\bibitem[Artin and Mazur(1977)]{fg-artin-mazur}
Artin, M., \& Mazur, B. (1977). Formal groups arising from algebraic varieties.
\emph{Annales scientifiques de l'\'Ecole Normale Sup\'erieure}, 4e s\'erie,
\emph{10}(1), 87--131. \url{https://doi.org/10.24033/asens.1322}.

\bibitem[Stienstra(1987)]{fg-stienstra}
Stienstra, J. (1987). Formal group laws arising from algebraic varieties.
\emph{American Journal of Mathematics, 109}(5), 907--925.
\url{https://doi.org/10.2307/2374494}.

\bibitem[Hazewinkel(1977)]{fg-hazewinkel}
Hazewinkel, M. (1977). Constructing formal groups I: The local one dimensional
case. \emph{Journal of Pure and Applied Algebra, 9}, 131--149.
\url{https://ir.cwi.nl/pub/9959/9959D.pdf}.

\bibitem[Silverman(2009)]{fg-silverman}
Silverman, J. H. (2009). \emph{The arithmetic of elliptic curves}
(2nd ed., Graduate Texts in Mathematics, Vol.~106). Springer.
\url{https://doi.org/10.1007/978-0-387-09494-6}.

\bibitem[Integral invariant derivatives(2026)]{fg-invariant-derivatives}
\emph{Integral invariant derivatives at every formal height: Optimal affine
error, divided-power distributions, Cartier duality, and binomial homotopy
types}. (2026, November 21). Unpublished research manuscript.
The one-dimensional estimates are in \S\S2--3; the specified Honda-dual
directions are in \S6. Complete proofs of the estimates used here appear
in Section~\ref{fg:sec:derivatives}.

\bibitem[Kawakami et~al.(2025a)]{fg-kty-one}
Kawakami, T., Takamatsu, T., \& Yoshikawa, S. (2025).
\emph{Fedder type criteria for quasi-$F$-splitting I}
(arXiv:2204.10076v3; revision of November~21, 2025).
\url{https://arxiv.org/abs/2204.10076v3}.

\bibitem[Batubara et~al.(2025)]{fg-bgp}
Batubara, R., Garzella, J. J., \& Pan, A. (2025).
\emph{K3 surfaces of any Artin--Mazur height over $\mathbb F_5$ and
$\mathbb F_7$ via quasi-$F$-split singularities and GPU acceleration}
(arXiv:2502.12428v3).
\url{https://arxiv.org/abs/2502.12428v3}.

\bibitem[Illusie(1979)]{fg-illusie}
Illusie, L. (1979). Complexe de de Rham--Witt et cohomologie cristalline.
\emph{Annales scientifiques de l'\'Ecole Normale Sup\'erieure},
4e s\'erie, \emph{12}(4), 501--661.
\url{https://doi.org/10.24033/asens.1374}.

\bibitem[Berthelot et~al.(2007)]{fg-bbe}
Berthelot, P., Bloch, S., \& Esnault, H. (2007).
On Witt vector cohomology for singular varieties.
\emph{Compositio Mathematica, 143}(2), 363--392.
\url{https://arxiv.org/abs/math/0510349}.

\bibitem[van der Geer and Katsura(2003)]{fg-vdgk}
van der Geer, G., \& Katsura, T. (2003).
On the height of Calabi--Yau varieties in positive characteristic.
\emph{Documenta Mathematica, 8}, 97--113.
\url{https://doi.org/10.4171/DM/140}.

\bibitem[Kawakami et~al.(2025b)]{fg-kty-two}
Kawakami, T., Takamatsu, T., and Yoshikawa, S. (2025b).
\emph{Fedder type criteria for quasi-$F$-splitting II}.
arXiv:2511.17270v1.
\url{https://arxiv.org/abs/2511.17270v1}.

\InputIfFileExists{refs_presentations.tex}{}{}
\InputIfFileExists{refs_saturation.tex}{}{}
\InputIfFileExists{refs_structured.tex}{}{}
\InputIfFileExists{refs_fixed.tex}{}{}
\end{thebibliography}
\end{document}
